% concatenated from MixedHybrid.tex
\documentclass[12pt]{amsart}
%\usepackage{amsart}
%
%\usepackage{stackengine}
%\usepackage{scalerel}

\usepackage{amsmath}
\usepackage{amssymb}
\usepackage{amstext}
\usepackage{amscd}
\usepackage[frenchb,english]{babel}
\selectlanguage{english}
\usepackage{enumitem}
\usepackage{greekctr}
\usepackage[matrix,arrow,ps]{xy}
\usepackage[bookmarks=true,colorlinks=true,citecolor=blue,urlcolor=blue,linkcolor=magenta]{hyperref}

\def\labelitemi{--}

%\newtheorem{teo}{Theorem}[section]
\newtheorem{theorem}{Theorem}[section]
%\newtheorem{prop}[teo]{Proposition}
\newtheorem{proposition}[theorem]{Proposition}
%\newtheorem{lem}[teo]{Lemma}
\newtheorem{lemma}[theorem]{Lemma}
%\newtheorem{cor}[teo]{Corollary}
\newtheorem{corollary}[theorem]{Corollary}
%\newtheorem{conj}[teo]{Conjecture}
\newtheorem{conjecture}[theorem]{Conjecture}
%\newtheorem{exe}[teo]{Example}
%\newtheorem{defi}[teo]{Definition}
\newtheorem{definition}[theorem]{Definition}
%\newtheorem{ques}[teo]{Question}
%\newtheorem{rem}[teo]{Remark}


\newcommand{\mupp}{\mu^\ddagger}

\newcommand{\OK}{{\cal O}_K}
\newcommand{\IM}{\mbox{Im}}
\newcommand{\Pic}{\mbox{Pic}}
\newcommand{\Hom}{\mbox{Hom}}
\newcommand{\Spec}{\mbox{Spec}}
\newcommand{\supp}{\mbox{supp}}
\newcommand{\ord}{\mbox{ord}}
\newcommand{\pgcd}{\mbox{pgcd}}
\newcommand{\De}{\Delta_{E/K}}
\newcommand{\oomega}{\omega_{{\cal X}/\OK}}
\newcommand{\Oomega}{\Omega^1_X}
\newcommand{\iii}{\sqrt{-1}}

\newcommand{\Alb}{{\rm Alb}}
\newcommand{\PGL}{{\rm PGL}}
\newcommand{\End}{{\rm End}}
\newcommand{\GSp}{{\rm GSp}}
\newcommand{\GL}{{\rm GL}}
\newcommand{\SL}{{\rm SL}}
\newcommand{\Sh}{{\rm Sh}}
\newcommand{\Lie}{{\rm Lie}}
\newcommand{\intt}{{\rm int}}
\newcommand{\Aut}{{\rm Aut}}
\newcommand{\Supp}{{\rm Supp}}
\newcommand{\Disc}{{\rm Disc}}
%\DeclareMathOperator{\Disc}{Disc}

\newcommand{\Gal}{{\rm Gal}}
\newcommand{\tors}{{\text{tors}}}
\newcommand{\Res}{{\rm Res}}
\newcommand{\MT}{{\rm MT}}
\newcommand{\der}{{\rm der}}
\newcommand{\ab}{{\rm ab}}
\newcommand{\im}{{\rm im}}
\newcommand{\ad}{{\rm ad}}
\newcommand{\TT}{{\mathbb T}}
\newcommand{\FF}{{\mathbb F}}
\newcommand{\F}{{\mathbb F}}
\newcommand{\CC}{{\mathbb C}}
\newcommand{\RR}{{\mathbb R}}
\newcommand{\R}{{\mathbb R}}
\newcommand{\ZZ}{{\mathbb Z}}
\newcommand{\Z}{{\mathbb Z}}
\newcommand{\QQ}{{\mathbb Q}}
\newcommand{\Q}{{\mathbb Q}}
\newcommand{\C}{{\mathbb C}}
\newcommand{\NN}{{\mathbb N}}
\newcommand{\DD}{{\mathbb D}}
\newcommand{\HH}{{\mathbb H}}
\newcommand{\PP}{{\mathbb P}}
\newcommand{\Hcal}{{\mathcal H}}
\newcommand{\GG}{{\mathbb G}}
\newcommand{\SSS}{{\mathbb S}}
\newcommand{\AAA}{{\mathbb A}}
\newcommand{\A}{{\mathbb A}}
\newcommand{\UUU}{{\mathbb U}}
\newcommand{\lto}{\longrightarrow}
\newcommand{\SU}{{\rm SU}}
\newcommand{\Sp}{{\rm Sp}}
\newcommand{\codim}{{\operatorname{codim}}}

\renewcommand{\Gal}{{\operatorname{Gal}}}



\def\Fp{\mathfrak{p}}
\def\FS{\mathfrak{S}}
\def\FG{\mathfrak{G}}
\def\FH{\mathfrak{H}}
\def\FK{\mathfrak{K}}
\def\FT{\mathfrak{t}}
\def\Fl{\mathfrak{l}}
\def\Fs{\mathfrak{s}}
\def\Fg{\mathfrak{g}}
\def\Fh{\mathfrak{h}}
\def\Fu{\mathfrak{u}}
\def\Ff{\mathfrak{t}}
\def\Fm{\mathfrak{m}}
\def\Fz{\mathfrak{z}}

\def\FJ{{\cal G}}

\newcommand{\cHH}{{\mathcal H}}
\newcommand{\eps}{{\varepsilon}}

\newcommand{\cB}{{\mathcal B}}
\newcommand{\cK}{{\cal K}}
\newcommand{\cT}{{\cal T}}
\newcommand{\cH}{{\mathcal{H}}}
\newcommand{\cE}{{\cal E}}
\newcommand{\cF}{\mathcal{F}}
\newcommand{\cA}{{\mathcal A}}
\newcommand{\cQ}{{\cal Q}}
\newcommand{\cP}{{\cal P}}
\newcommand{\cD}{{\cal D}}
\newcommand{\cG}{{\cal G}}
\newcommand{\cS}{\mathcal{S}}
\newcommand{\cO}{\mathcal{O}}
\newcommand{\cL}{{\mathcal L}}
\newcommand{\cZ}{{\mathcal Z}}
\newcommand{\cX}{{\mathcal X}}
\newcommand{\cY}{{\mathcal Y}}
\newcommand{\clambda}{{\tilde{ \lambda}}}

\newcommand{\lan}{\langle}
\newcommand{\ol}{\overline}
\newcommand{\wh}{\widehat}
\newcommand{\ran}{\rangle}
\newcommand{\vol}{{\mbox{vol}}}
\newcommand{\doo}{\frac{\deg (\omega_{X/O_K})}{[K:\QQ]}}
\newcommand{\nd}{\frac{\log N_{K/\QQ}\Delta_E}{[K:\QQ]}}
\newcommand{\hfi}{\sum_{\sigma}
\frac{1}{12}\frac{\log
\vert\Delta(\tau_{\sigma})\IM(\tau_{\sigma})^6\vert}{[K:\QQ]}}
\newcommand{\OOO}{{\cal O}}
\newcommand{\ox}{\overline{x}}
\newcommand{\ocL}{\overline{\cal L}}
\newcommand{\oK}{\overline{K}}
\newcommand{\oL}{\overline{L}}
\newcommand{\oc}{\mbox{\^c}}
\newcommand{\oY}{\overline{Y}}
\newcommand{\oQ}{\overline{\QQ}}
\newcommand{\oCH}{\widehat{CH}}
\newcommand{\oE}{\overline{E}}
\newcommand{\os}{\overline{\sigma}}
%\renewcommand{\refname}{R\'ef\'erences}
\renewcommand{\contentsname}{Table des mati\`eres}
\newcommand{\vo}{\mbox{vol}}
\newcommand{\hh}{{\cal H}}
\newcommand{\tg}{\widetilde{g}_{Ar}}
\newcommand{\spec}{\mbox{Spec}}
%\newtheorem{prop}{Proposition}[section]
%\newtheorem{lem}{Lemme}[section]
%\newtheorem{cor}{Corollaire}[section]
%\newtheorem{teo}{Th\'eor\`eme}[section]
%\newtheorem{defi}{D\'efinition}[section]
\newtheorem{Aprop}{Proposition}
\renewcommand{\theAprop}{\Alph{Aprop}}
\newtheorem{Alem}[Aprop]{Lemme}
\newtheorem{Acor}[Aprop]{Corollaire}
\newtheorem{Ateo}[Aprop]{Th\'eor\`eme}
\newtheorem{Adefi}[Aprop]{D\'efinition}
\newcommand{\GaN}{\Gamma_0(N)}
\newcommand{\vp}{\varphi}
\newcommand{\tz}{{\Bbb T}}
\newcommand{\tnz}{{\Bbb H}}
\newcommand{\tq}{{\Bbb T}_{\Bbb Q}}
\newcommand{\tc}{{\Bbb T}_{\Bbb C}}
\newcommand{\tr}{{\Bbb T}_{\Bbb R}}
\newcommand{\dt}{{\mbox{det}}}
%\newcommand{\Fp}{{\bf{\mbox{F}}}_p}
\newcommand{\szz}{{\cal S}_{\Bbb Z}(2,\GaN)}
\newcommand{\sz}{{\cal S}_{\Bbb Z}}
\newcommand{\snz}{{\cal S}_{\Bbb Z}^n}
\newcommand{\snq}{{\cal S}_{\Bbb Q}^n(2,\GaN)}
\newcommand{\snr}{{\cal S}_{\Bbb R}^n(2,\GaN)}
\newcommand{\scc}{{\cal S}_{\Bbb C}(2,\GaN)}
\newcommand{\sr}{{\cal S}_{\Bbb R}(2,\GaN)}
\newcommand{\Jnz}{J_0(N)_{\Bbb Z}^{n}}
\newcommand{\Jnq}{J_0(N)_{\Bbb Q}^{n}}
\newcommand{\lie}{{\mbox{Lie}}}
\newcommand{\g}{{\mbox{g}}}
\renewcommand{\refname}{Bibliography.}
\renewcommand{\contentsname}{Table of contents.}
\newcommand{\diag}{{\rm diag}}
\newcommand{\Frob}{{\rm Frob}}
\newcommand{\sm}{{\rm sm}}
%\newcommand{\der}{{\rm der}}
\newcommand{\wt}[1]{{\widetilde{#1}}}
\newcommand{\ul}{\underline}

\DeclareMathOperator{\coker}{\operatorname{coker}}
%\DeclareMathOperator{\Hom}{\operatorname{Hom}}

\newcommand{\tens}{\otimes}
\usepackage{mathtools}
\DeclareMathOperator{\Radu}{{\mathrm{Rad}^{\mathrm{u}}}}
\DeclarePairedDelimiterX{\Nm}[1]{\left\lVert}{\right\rVert}{#1}
\DeclarePairedDelimiterX{\norm}[1]{\lVert}{\rVert}{#1}
\newcommand{\sous}{\backslash}

\newcommand{\abs}[1]{\left|#1\right|}


\newenvironment{prf}[1]{\trivlist
\item[\hskip \labelsep{\it
#1.\hspace*{.3em}}]}{~\hspace{\fill}~$\square$\endtrivlist}
%\newenvironment{proof}{\begin{prf}{\bf Proof}}{\end{prf}}
\title[Hybrid conjecture in a mixed Shimura variety]{Hybrid conjecture in a mixed Shimura variety}
\author{Rodolphe Richard}
\email{Rodolphe.Richard@normalesup.org}
\address{The University of Manchester,
Alan Turing Building,
Oxford Rd,
Manchester
M13 9PL,
UK}
\author{Andrei Yafaev}
\email{a.yafaev@ucl.ac.uk}
\address{UCL, Department of Mathematics, Gower Street, WC1E 6BT, London, UK}
\date{\today}
\begin{document}
\setcounter{tocdepth}{1}
\setcounter{secnumdepth}{4}
\maketitle

\begin{abstract}The authors previously formulated the hybrid conjecture, unifying  André-Pink-Zannier and André-Oort conjectures, and proved it in Shimura varieties of abelian type. We study its analogue for mixed Shimura varieties, and consider the prime example, the universal abelian scheme~$\mathcal{A}_g\to \mathbb{A}_g$. 


In a radical departure from the Pila-Zannier strategy, typically applied to such questions, we employ instead a combination of equidistribution and o-minimality

Our main result strictly includes the following:  the Hybrid Conjecture, in particular the André-Pink-Zannier and André-Oort conjectures, for~$\mathbb{A}_g$; the mixed André-Oort conjecture for~$\mathcal{A}_g$; and Manin-Mumford conjecture for arbitrary abelian varieties. It also yields an analogue of the ``Manin-Mumford in arithmetic pencil", a result of \cite{BRU}, for abelian schemes over a variety.

The mixed hybrid conjecture in~$\mathcal{A}_g$ also encompasses the Mordell-Lang conjecture. We actually reduce the mixed hybrid conjecture for~$\mathbb{A}_g$ to its "mordellic" part. 

We also prove, Galois-theoretic results: uniform variants on the Ribet's Kummer theory of Abelian varieties, and Serre's theorem on Lang's conjecture.
\end{abstract}

\tableofcontents

\section{Introduction}

Let~$g\in\Z_{\geq0}$, let~$\mathbf{A}_g$ be the moduli space of principally polarised abelian varieties of dimension~$g$ over~$\C$.
We have~$\mathbf{A}_g(\C)\simeq Sp(2g,\Z)\backslash \mathfrak{H}_g$ where~$\mathfrak{H}_g$ is
 the Siegel upper half-space. Let~$\Gamma:=\ker (Sp(2g,\Z)\to Sp(2g,\Z/(3)))$ or any other torsion free congruence subgroup. Then $\mathbf{A}_\Gamma:=\Gamma\backslash \mathfrak{H}_g$ is a \emph{fine}
moduli space of abelian varieties equipped with some appropriate level structure of type~$\Gamma$. There exists a corresponding universal abelian 
scheme~$\mathcal{A}_\Gamma\to  \mathbf{A}_\Gamma$. 

Then~$\mathbf{A}_g$ is the archetype of a Shimura variety and~$\mathcal{A}_\Gamma$ is the archetype of a mixed Shimura variety.

In~\cite{RY:LMS1}, the authors introduced the hybrid conjecture, which encompass both the André-Oort conjecture and the André-Pink-Zannier conjecture, using the notion of \emph{hybrid orbit} in a Shimura variety, and, in~\cite{RY:LMS2} they proved it for Shimura varieties of abelian type.
Here, we define the analogue in the mixed Shimura variety~$\mathcal{A}_\Gamma$, which we call the \emph{mixed hybrid orbit}. 

This notion unifies the hybrid conjecture for~$\mathbf{A}_g$ -- which includes the André-Oort conjecture (AO) and the generalised André-Pink-Zannier conjecture (APZ) --  with the Mordell-Lang conjecture (ML) -- which includes the Mordell conjecture (MC) and the Manin-Mumford conjecture (MM) -- for an arbitrary abelian variety. This also generalises the mixed André-Oort conjecture (MxAO) for~$\mathcal{A}_\Gamma$.

Our main results are Theorems~\ref{thm:Main} and~\ref{thm:MMonHyb}. Theorem~\ref{thm:MMonHyb} implies all the former results:~(HC), (AO), (APZ), (MM), (MxAO), except the mordellic part of the Mordell-Lang conjecture:  (MC) and thus (ML). Theorem~\ref{thm:MMonHyb} is a particular case of Theorem~\ref{thm:Main}, which reduces the ``mixed hybrid conjecture" for~$\mathcal{A}_\Gamma$ to its "mordellic part".

This is a combination of two results, interesting on their own. 

A first result, Theorem~\ref{thm:weak} uses a combination of equidistribution and o-minimality, introduced by the first named author and E. Ullmo in their ergodic-theoretic proof of the geometric part of the André-Oort conjecture. Theorem~\ref{thm:weak} implies the analogue, for abelian scheme over a variety of characteristic zero, of a result of~\cite{BRU}, ``Manin-Mumford conjecture in artihmetic pencils'', which concerns abelian schemes over rings of algebraic integers.

Secondly are results on Galois representations attached to abelian varieties. Theorem~\ref{thm:Serre hybrid} is a reinforcement of Serre's Theorem~\ref{thm:Serre eAK} on a conjecture of Lang on homotheties of the Tate modules. Our version is uniform on a hybrid orbit of an abelian variety, which generalises the isogeny class by allowing CM factors and passing to quotient and to powers (cf. Lemma~\ref{lem:truc}). Our next result is Theorem~\ref{thm:uniform Kummer1}, which is uniform version, over a hybrid orbit of an abelian variety, of Ribet's Kummer theory~Th.~\ref{thm:Kummer} for Abelian varieties. We actually prove stronger results. Theorem~\ref{thm:Serre unif principal} goes beyond Lang's conjecture, considering any element in the centre of the Mumford-Tate group, not only homotheties. Theorem~\ref{thm:uniform Kummer1} proves uniformity for arbitrary abelian extension, not only those coming by CM theory. We also give a new formulation, Theorem~\ref{thm:uniform Faltings}, of Tate uniform integral conjecture, derived from Faltings' theorem on Tate conjectures, of independent interest. 
These Galois theoretic results are not only over number fields, but also for finitely generated extension of~$\Q$. 

Because our approach is based on equidistribution, it may lead to refinements, though we did not explore that question. For instance, for cases in which one can prove topologic refinements of the Hybrid conjecture in the Simura variety~$\mathbf{A}_g$ using equidistribution (e.g. \cite{Duke,Ilya,RY:SHecke}), does our approach deduce to corresponding refinements in the mixed Shimura variety~$\mathcal{A}_g$?

\subsection{Main statements}
When~$A$ is a principally polarised abelian variety over~$\C$, and~$Q\in A(\C)$, and~$\lambda$ is a $\Gamma$-level structure\footnote{There is~$N\in\Z_{\geq1}$ such that~$\Gamma\geq \Gamma(N):=\ker (Sp(2g,\Z)\to Sp(2g,\Z/(N))$, and a~$\Gamma$-level structure on~$A$ amounts to an right coset~$\lambda\in \left(Isom(\Z/(N))^{2g},A[N])\right)/\Gamma$.} on~$A$, we denote by~$A\in \mathbf{A}_g(\C)$, by~$(A,\lambda)\in \mathbf{A}_\Gamma(\C)$ and~$(A,Q,\lambda)\in \mathcal{A}_\Gamma(\C)$ the corresponding moduli point.


%A level structure on~$A$ is an isomorphism~$\lambda:(\Z/(3))^{2g}\to A(\C)[3]$ such that the Weyl pairing~$w_3:A(\C)[3]^2\to \exp(\frac{1}{3}2\pi i\Z)$ and the standard symplectic form on~$(\Z/(3))^{2g}$ satisfy~$w_3(\lambda(P),\lambda(P')=\exp$... 

In the following, we consider that abelian varieties of dimension zero have CM.
\begin{definition}[Hybrid, mixed hybrid and mordellic orbits]
\label{def:mixed hybrid orbits}
Let~$A$ be an abelian variety over~$\C$ and let~$P\in A(\C)$. %Let~$\Sigma_g(A)\subseteq \mathbf{A}_g$.
We define
\begin{subequations}
%\begin{multline}
%\Sigma_g(A):=\{B\in \mathbf{A}_g|\exists n\in \Z_{\geq1},
%\exists \phi:A^n\to B,\\ B/\phi(A)\text{ has CM}\}.
%\end{multline}
\begin{multline}
\Sigma_\Gamma(A):=\{(B,\lambda)\in \mathbf{A}_\Gamma|\exists n\in \Z_{\geq1},
\exists \phi:A^n\to B,\\ B/\phi(A^n)\text{ has CM}\}.
\end{multline}
\begin{multline}\label{eq:def Hgamma}
H_\Gamma(A,P):=\{(B,Q,\lambda)\in \mathcal{A}_\Gamma|\exists d,n\in \Z_{\geq1}, \exists \phi:A^n\to B, 
\\B/\phi(A^n)\text{ has CM and }d\cdot Q=\phi(P,\ldots,P)\}.
\end{multline}
\begin{multline}
M_\Gamma(A,P):=\{(B,Q,\lambda)\in \mathcal{A}_\Gamma|\exists n\in \Z_{\geq1},\exists \phi:A^n\to B,\\ B/\phi(A^n)\text{ has CM and }Q=\phi(P,\ldots,P)\}.
\end{multline}
\end{subequations}
We say that~$H_\Gamma(A,P)\subseteq \mathcal{A}_g$ is the \emph{mixed hybrid orbit of~$(A,P)$
in~$\mathcal{A}_\Gamma$}
and that~$M_\Gamma(A,P)\subseteq H_\Gamma(A,P)\subseteq \mathcal{A}_\Gamma$ is the \emph{mordellic orbit of~$(A,P)$ in~$\mathcal{A}_\Gamma$}.

%More generally, 
%let~$\alpha:\mathcal{A}\to S$ be a principally polarised abelian scheme of relative dimension~$g$ over a complex algebraic variety~$S$. We define
%\begin{subequations}
%\begin{equation}\label{subeq1}
%\Sigma_\mathcal{A}(A):=\{s\in S(\C)|\mathcal{A}_s\in \Sigma_g(A)\}.
%\end{equation}
%\begin{multline}\label{subeq2}
%H_\mathcal{A}(A,P):=\{Q\in \mathcal{A}|\exists d,n\in \Z_{\geq1}, \exists \phi:A^n\to \mathcal{A}_{\alpha(Q)}, 
%\\\mathcal{A}_{\alpha(Q)}/\phi(A)\text{ has CM and }d\cdot Q=\phi(P,\ldots,P)\}.
%\end{multline}
%\begin{multline}\label{subeq3}
%M_\mathcal{A}(A,P):=\{Q\in \mathcal{A}|\exists n\in \Z_{\geq1},\exists \phi:A^n\to \mathcal{A}_{\alpha(Q)},\\ \mathcal{A}_{\alpha(Q)}/\phi(A)\text{ has CM and }Q=\phi(P,\ldots,P)\}.
%\end{multline}
%\end{subequations}
\end{definition}

%In the above definition, we don't require the morphisms of abelian variety~$\phi$ to be compatible, when it makes sense, with polarisations and level structure.
%For instance, $\Sigma_g(A)\subseteq \mathbf{A}_g$ is a (countable) union of unpolarised isogeny classes.



The relation with the Hybrid orbit defined in~\cite{RY:LMS1,RY:LMS2} is the following.
\begin{lemma}\label{lem:truc}
Let~$(GSp(2g),\pm\mathfrak{H}_g)$ denote Siegel's Shimura datum. Let~$M_A$ be the Mumford-Tate group of a principally polarised~$A\in \mathbf{A}_g(\C)$ and let~$x_A:\C^\times\to M_A(\R)$
be its Hodge cocharater, and let~$X_A=M_A(\R)\cdot x_A$.
Let~$\Sigma_{\pm\mathfrak{H}_g}(M_A,X_A,x_A)$ be as in~\cite[Def. 6]{RY:LMS1}. Then
\[
\Sigma_g(A)\subseteq GSp(2g,\Z)\backslash\Sigma_{\pm\mathfrak{H}_g}(M_A,X_A,x_A).
\]
\end{lemma}
\begin{proof}Let~$A$ be an abelian variety over~$\C$. Let~$A\to \bigoplus_{i=1}^d A_i^{n_i}$ be an isogeny where the~$A_i$ are simple and pairwise non-isogeneous and~$n_i\geq 1$. Let~$x_i:=x_{A_i}^{ad}:\C^\times \to MT(A_i)\to M_i:=MT(A_i)^{ad}$ be the morphism induced by the Hodge cocharacter for~$A_i$. We recall that the image of
\[
(x_1,\ldots,x_d):\C^\times\to M_1\times\ldots\times M_d
\]
is~$\Q$-Zariski dense. Let~$\xi=\bigoplus x_i^{\oplus n_i}:\C^\times \to M_1^{n_1}\times\ldots\times M_d^{n_d}$.
Let~$M'_i\leq M_i^{n_i}$ be the image of~$g\mapsto (g,\ldots,g):M_i\to M_i^{n_i}$. Then the image of~$\xi$ is $\Q$-Zariski dense in
\[
M':=M_1'\times\ldots\times M_d'\leq M_1^{n_1}\times\ldots\times M_d^{n_d}.
\]
We claim that there is a natural isomorphism~$\iota:M^{ad}_A\simeq M'$ such that~$\xi=\iota\circ x^{ad}_A$ where~$x^{ad}_A:\C^\times\to M_A\to M_A^{ad}$ is induced by the Hodge cocharacter of~$A$. Using the isomorphisms~$M_i\to M'_i$ one can identify~$x_A^{ad}$ with~$(x_1,\ldots,x_d):\C^\times\to M_1\times\ldots \times M_d$.

 Consider an abelian variety~$B$ over~$\C$ and let similarly~$B\to \bigoplus_{j=1}^e B_j^{m_j}$ be an isogeny, such that  the~$B_j$ simple and pairwise non-isogeneous.

From the above discussion one sees that there exists
\[
\phi:M^{ad}_A\to M^{ad}_B
\]
 such that~$x_B^{ad}=\phi\circ x_A^{ad}$ if and only: for every~$1\leq j\leq e$ such that $M^{ad}_{B_j}$ is not trivial, there is~$1\leq i\leq d$ such that~$A_i$ is isogeneous to~$B_j$. This is readily equivalent to the property
\[
\exists n\in \Z_{\geq1},
\exists \phi:A^n\to B,\\ B/\phi(A^n)\text{ has CM}.
\]
The Lemma follows.
\end{proof}
The proof of Lemma~\ref{lem:truc} gives the following interpretation of hybrid orbits for abelian varieties.
The abelian variety~$B$ is in te hybrid orbit of~$A$ if and only if every simple quotient of~$B$ is CM or is also a quotient of~$A$.

The main result of~\cite{RY:LMS2} implies the following. We refer to~\cite{RY:LMS2} for the definition of a weakly special subvariety in a pure Shimura variety.
\begin{theorem}[Hybrid conjecture for~$\A_\Gamma$ \cite{RY:LMS2}]\label{thm:IHES}
Let~$\Sigma\subseteq \Sigma_g(A)$ be a subset. Then the irreducible components of~$\ol{\Sigma}^{Zar}$ are weakly special subvarieties.
\end{theorem}


Theorem~\ref{thm:IHES} establishes part~\eqref{conj:0} of the following Conjecture.
\begin{conjecture}[Mixed hybrid conjecture for universal abelian schemes]
\label{conj}
Let~$A$ be an abelian variety over~$\C$, and~$P\in A(\C)$. Let~$H_{\Gamma}(A,P)$ be as in~\eqref{eq:def Hgamma} and let~$\Sigma\subseteq H_{\Gamma}(A,P)$ be a subset. Let~$V\subseteq \ol{\Sigma}^{Zar}$ be an irreducible component and let~$S\subseteq \A_\Gamma$ be the image of~$V$.
\begin{enumerate}\setcounter{enumi}{0}
\item \label{conj:0} Then~$S\subseteq \A_\Gamma$ is a weakly special subvariety.
\item \label{conj:1}
Assume that for some~$s\in S$ we have~$V\cap (A_s)_{tors}\neq \emptyset$. Then there exists~$d\in \Z_{\geq1}$ and an abelian subscheme~$\mathcal{B}\leq \mathcal{A}_{\Gamma}|_S$ such that~$d\cdot V=\mathcal{B}$. Equivalently,~$V$ is an irreducible component of~$\mathcal{B}+\mathcal{A}_{\Gamma}|_S[d]$.
\item \label{conj:2}
In general, let~$\eta\in S$ be the generic point of~$S$ and~$V_\eta\subseteq \mathcal{A}_\eta$ be the fibres of~$V\subseteq \mathcal{A}_\Gamma$ over~$\eta$.  
Then there exists~$d\in \Z_{\geq1}$ and~$\phi:A\to \mathcal{A}_\eta$ and an abelian subvariety~$B\leq A_\eta$ such that~$d\cdot V_\eta=\phi(P)+B\subseteq \mathcal{A}_\eta$.
\end{enumerate}
\end{conjecture}
\subsubsection{}\label{sec:mordell lang case}

If~$S$ is a singleton~$\{s\}$, and~$A=\mathcal{A}^n_s$ and~$P=(P_1,\ldots,P_n)$, then Conjecture~\ref{conj} implies the \emph{Mordell-Lang conjecture} for the finitely generated subgroup
\[\{Q\in \mathcal{A}_s|\exists m\in\Z_{\geq1},\phi_1,\ldots,\phi_n\in\End(\mathcal{A}_s), mQ=\phi_1(P_1)+\ldots+\phi_n(P_n)\}\leq \mathcal{A}_s.\]
\subsubsection{}

Assume~$A=0$. Therefore~$P=0$ and~$\Sigma$ is the set of~$(B,Q,\lambda)$ in $\mathcal{A}_\Gamma$ such that~$B$ is CM and~$Q$ is a torsion point. That case of Conjecture~\ref{conj} is the \emph{mixed André-Oort conjecture} for~$\mathcal{A}_\Gamma$. It implies Manin-Mumford conjecture only for CM abelian varieties.
\subsubsection{}\label{sec:conj for P0}	
Assume~$P=0$ for an arbitrary~$A$. Then Conjecture~\ref{conj} implies (\cite{Pink})
\begin{itemize}
\item the \emph{Hybrid conjecture} in~$\mathbf{A}_g$ (\cite{RY:LMS2}), and in particular
\begin{itemize}
\item  the \emph{André-Oort Conjecture} for~$\mathbf{A}_g$ (\cite{PT,Tsimerman,AGHM,YZ}), and
\item  the \emph{generalised André-Pink-Zannier Conjecture} for~$\mathbf{A}_g$ (\cite{RY:IHES});
\end{itemize}
\item the \emph{Manin-Mumford conjecture} for arbitrary abelian varieties;
\item the \emph{mixed André-Oort Conjecture} in~$\mathcal{A}_g$ (cf.\,\cite{Gao}).
\end{itemize}

\subsubsection{Results}
Our main result Theorem~\ref{thm:Main} is the following.

\begin{theorem}\label{thm:Main}
Let~$A$,~$P$,~$\Sigma$, $V$ and $S$ be as in Conjecture~\ref{conj} and $M_\Gamma(A,P)$ be as in Definition~\ref{def:mixed hybrid orbits}.
% be an abelian variety over~$\C$, and~$P\in A(\C)$. Let~$H_{\Gamma}(A,P)$ be as in~\eqref{subeq2}. Let~$V\subseteq \mathcal{A}_\Gamma$ be an irreducible component of the Zariski closure of some subset~$\Sigma\subseteq H_{\Gamma}(A,P)$. 
%
%Let~$S$ be the image of~$V$ in~$\mathbf{A}_\Gamma$ and let~$\mathcal{A}_{\Gamma,S}\to S$ be restriction to~$S$ of the universal abelian scheme.

Then there exists~$d\in \Z_{\geq1}$ and an abelian subscheme~$\mathcal{B}\leq \mathcal{A}_{\Gamma}|_{S}$ such that~$d\cdot V=W+\mathcal{B}$ with
\[
W=\ol{W\cap M_\Gamma(A,P)}^{Zar}.
\]
\end{theorem}
When~$P=0$, the set~$M_g(A,P)$ is contained in the neutral section~$\mathbf{A}_g\to \mathcal{A}_g$ and for every~$n$, the point~$Q_n$ has finite order in~$B_n$. 
Theorem~\ref{thm:Main} gives the following, which encompasses all the Conjectures mentionned in \S~\ref{sec:conj for P0}.
\begin{corollary}\label{thm:MMonHyb} Conjecture~\ref{conj} holds if~$P=0$. We are moreover in the case~\eqref{conj:1} of Conjecture~\ref{conj}.
\end{corollary}
	
Note that Theorem~\ref{thm:Main} reduces Conjecture~\ref{conj} to its ``Mordellic part''.
\begin{corollary}It is enough to prove Conjecture~\ref{conj} for~$\Sigma$ such that~$\Sigma\subseteq M_\Gamma(A,P)$.
\end{corollary}
\subsection{The Mordellic case}
Two of the main strategies used to tackle problems of Zilber-Pink type are: the Pila-Zannier strategy; and Equidistribution.
Both approaches rely on "large Galois orbits".
Moreover, the Pila-Zannier strategy needs quantitative lower bounds, in terms of the height of the points one considers, for the cardinality of the Galois orbits.

Let us consider the case~\S\ref{sec:mordell lang case}, when~$P$ is not a torsion point. Let~$K/\Q$ be a finitely generated extension such that~$A$ as a model over~$K$ such that~$P\in A(K)$, and let~$g=\dim(A)$. Then~$M_\Gamma(A,P)$ contains the infinite subset~$\{(A,n\cdot P,\lambda)|n\in\Z\}\subseteq \mathcal{A}_\Gamma(K)$, and~$H_\Gamma(A,P)$ contains 
\[
\{(A,Q,\lambda)|\exists p\in \Z,q \in\Z\smallsetminus\{0\}, q\cdot Q=p\cdot P\}
\]

When~$gcd(p,q)=1$, the height of~$(A,Q,\lambda)$ is roughly~$\max\{p;q\}$. Then
\begin{itemize}
\item Pila-Zannier strategy applies to sequences~$(A,Q_n,\lambda)$ such that~$q_n\to \infty$ and~$p_n\leq f(q_n)$, with~$f(x)=a+x^b$ for some~$a,b\in\R$, or some other explicit~$f:\R\to \R$ in finer variant of the strategy.
\item Equidistribution can only apply to sequences~$(A,Q_n,\lambda)$ such that~$q_n\to \infty$.
\end{itemize}
One can of course find, for every~$f:\R\to\R$, a sequence~$(p_n,q_n)$ such that~$p_n>f(q_n)$ for~$n\gg0$.

In this setting, the methods of this article apply to all sequences for which~$q_n\to +\infty$. This is what allow us to reduce the general case to the case where~$q_n$ is bounded, or ultimately~$q_n=1$. Theorem~\ref{thm:Main} is something the Pila-Zannier strategy may fail to achieve.

In the case~$\#S=1$, Theorem~\ref{thm:Main} reduces the Mordell-Lang conjecture for abelian varieties to its Mordellic part~\cite{Fal91} by Faltings. This reduction was originally achieved in~\cite{McQ} by McQuillan (which also covers semi-abelian varieties), using a method introduced by Hindry for the Manin-Mumford conjecture.

In that context Theorem~\ref{thm:Main} recovers results of~\cite{McQ} using an entirely new method. Our method  builds on the Equidistribution proof~\cite{RWeyl} of Manin-Mumford by Weyl criterion, and combines equidistribution with o-minimality following a trick introduced by~\cite{RU}. 

\subsection{Strategy}

Our proof of Theorem~\ref{thm:Main} combines two ingredients.
\subsubsection{Equidistibution} The first is based on Theorem~\ref{thm:equi0}, which is adapted from a trick introduced by the first author in Ullmo in~\cite{RU}. 

The equidistribution method, cf. \cite{U}, traditionally proceeds as follows. We consider a sequence of subsets~$E_n\subseteq S$ in an ambiant variety~$S$. The goal is to determine~$V:=\ol{\bigcup E_n}^{Zar}\subseteq S$. One constructs, for each~$n$, a measure~$\mu_n\in Prob(S)$ with support~$E_n$. One studies the limits~$\mu_\infty$. A difficulty is that one needs to establish:
\begin{enumerate}
\item \label{enum:equi1}
that the limit measure~$\mu_\infty$ is non zero.
\item \label{enum:equi2}
that~$\Supp(\mu_\infty)$ contains~$\bigcup E_n$.
\end{enumerate}
One can then infer that~$V=\ol{Supp(\mu_\infty)}^{Zar}$.

Theorem~\ref{thm:equi0} allows us to obtain information on~$V$ when~\eqref{enum:equi2} fails, or even when~$\mu_\infty=0$. This argument can be applied when the situation, for some o-minimal structure, is related to definable families of definable subsets.

A typical example is studied in this article, and this aticle can serve as an introduction to this new paradigm of the equidistribution method.

In the universal abelian scheme~$S=\mathcal{A}_\Gamma$, consider a sequence~$s_n\in \A_\Gamma$, and subsets~$E_n\subseteq A_{s_n}$ in the corresponding fibers. Note that~$E_n\not\subseteq \Supp(\mu_\infty)$, except when~$s_n$ is stationnary. Note that~$\A_\Gamma$ is not compact, and when~$s_n$ diverges to infinity, we necessarily have~$\mu_\infty=0$.

Using these ideas, we obtain Theorem~\ref{thm:weak}.
Theorem~\ref{thm:weak} is an analogue of~\cite[1.9]{BRU}. The latter studies "Manin-Mumford conjecture" in arithmetic pencils, whereas we study Mordell-Lang conjecture for an abelian scheme over a complex algebraic variety.

\subsubsection{Galois representations}

In \S\ref{sec:Galois} we consider two Galois representations: the Tate module attached to an abelian variety, and the affine Tate module (cf. \S\ref{sec:Kummer}) attached to~$(A,P)$ for some~$P\in A$.

The first main result of~\S\ref{sec:Galois}  Theorem~\ref{thm:Serre hybrid} is a stronger variant of Serre's Theorem~\ref{thm:Serre eAK} that is uniform for abelian varieties ranging through a hybrid orbit. Theorem~\ref{thm:Serre eAK} concerns homoteties~$\lambda\in GL(1,\wh{\Z})$ of a Tate module which are in the image of the Galois representation. 
Actually, we obtain a Theorem~\ref{thm:Serre unif principal}
which concerns all the elements~$\lambda\in Z(M)(\wh{\Z})$ in the centre~$Z(M)$ of the Mumford-Tate group of~$A$.

The second main result Theorem~\ref{thm:uniform Kummer1} is a stronger variant of Ribet's Theorem~\ref{thm:Kummer} on Kummer theory for abelian varieties, that is uniform for abelian varieties ranging through a hybrid orbit (cf. Theorem~\ref{thm:Gal En}). We adapt a presentation~\cite[App. 2]{Hindry} of a proof of Ribet's theorem. The critical difference is our use of Sah's theorem Sah's~\cite[App. 2, Lem.~D]{Hindry} in the proof of Proposition~\ref{prop:412}. We also needed to fully develop Hindry's arguments and make them quantitative. This allows us to obtain a precise quantitative form~\eqref{eq:Kummer explicite} of Ribet's theorem. For instance, we obtain (cf. \S\ref{sec:CD}) a stronger form of a Theorem of Rémond that is a "key tool" (\cite[p. 16]{CD}) in a recent article~\cite{CD} of Checcoli and Dill. 

We also give a reformulation Theorem~\ref{thm:uniform Faltings} of Faltings' theorem on Tate's conjecture.

\subsubsection{Combining the ingredients} In \S\ref{sec:proof}, we prove theorem~\ref{thm:Main}. We consider~$(A_{s_n},Q_n)\in H_\Gamma(A,P)$. Using \S\ref{sec:Galois}, we construct subsets~$E_n\subseteq A_{s_n}\cap Gal(\ol{K}/K)\cdot (A_{s_n},Q_n)$. This allows us to apply Theorem~\ref{thm:weak}. We finish the proof by comparing the conclusion of Theorem~\ref{thm:weak} with the property~$(A_{s_n},d\cdot Q_n)\in M_\Gamma(A,P)$.

\subsection{Acknowledgements}The first named author was supported by the University of Manchester and Grant EP/Y020758/1.


\section{Equidistribution in o-minimal families}
In this section, we use the term \emph{definable} with respect to a chosen~$o$-minimal structure (see~\cite{VdD}).
\begin{theorem}\label{thm:equi0}
Let~$X\times Y$ be a definable set, let~$y_n\in Y$ be a sequence, let~$\nu_n$ be a sequence of probability measures in~$X$, and assume that~$\nu_n$ has a limit~$\nu_\infty$ and that
\begin{equation}\label{eq:equi1}
\forall n\gg0, Supp(\nu_n)\subseteq Supp(\nu_\infty),
\end{equation}
and that~$Supp(\nu_\infty)$ is definable and that for any definable subset~$A\subseteq Supp(\nu_\infty)$ such that~$\dim(A)<\dim(Supp(\nu_\infty))$ one has~$\nu_\infty(A)=0$. 

\begin{itemize}
\item
Let~$V\subseteq X\times Y$ be a definable subset such that
\begin{equation}\label{eq:equi2}
\forall n\gg0, Supp(\nu_n)\times\{y_n\}\subseteq V.
\end{equation}
Then, for~$n\gg0$, 
\begin{equation}\label{eq:equi3}
\dim\Bigl(Supp(\nu_\infty)\times\{y_n\}\,\cap\,V\Bigr)=\dim(Supp(\nu_\infty)).
\end{equation}In particular,~$V$ contains a non-empty open subset of~$(Supp(\nu_\infty)+x_n)\times \{y_n\}$.
\item
Assume that~$X$ has a definable Lie group structure~$(x_1,x_2)\mapsto x_1\cdot x_2$, and let~$x_n\to x_\infty$ be a convergent sequence in~$X$.

Let~$V\subseteq X\times Y$ be a definable subset such that
\begin{equation}\label{eq:equi4}
\forall n\gg0, (Supp(\nu_n)\cdot x_n)\times\{y_n\}\subseteq V.
\end{equation}
Then, for~$n\gg0$, 
\begin{equation}\label{eq:equi6}
\dim\Bigl(\left(Supp(\nu_n)\cdot x_n\right)\times\{y_n\}\,\cap\,V\Bigr)=\dim(Supp(\nu_\infty)).
\end{equation}
In particular,~$V$ contains a non-empty open subset of~$(Supp(\nu_\infty)+x_n)\times \{y_n\}$.
\end{itemize}
\end{theorem}
\begin{proof}Without loss of generality we may assume that~$X$ is a bounded subset of~$\R^n$.
We apply~\cite[Th.~4.3, Cor.~4.4]{RU} with~$K=\ol{X}$; with~$b_i=0\in B:=\{0\}$ and~$A(b)=Supp(\nu_\infty)$; with~$b'_i=y_n$ in~$Y=B'$; with~$A'(b'_i)\times \{b'_i\}=V_i:=V\cap \left(X\times \{b'_i\}\right)$. The assumption~$\nu_\infty(A)=0$ of Theorem~\ref{thm:equi} implies the assumption~\cite[(4.3)]{RU}, and~\eqref{eq:equi1} implies the assumption~\cite[(4.3)]{RU}. Our assumption~\eqref{eq:equi2} gives~$\nu_n(V_n)=1$, and this implies the assumption~\cite[(4.4)]{RU}. Let~$d:=\dim(Supp(\nu_\infty))$. Then the assumption~\cite[(4.4)]{RU} is satisfied.

We deduce
\[
\forall n\gg0, d=\dim(\Supp(\nu_\infty)\cap V_n).
\]
This proves~\eqref{eq:equi3}.

Let us now prove~\eqref{eq:equi6}.

Let~$Y'=X\times Y$ and~$y'_n=(x_n,y_n)$. Let~$V'=\{(v',x,y)|(v'\cdot x,y)\in V\}$, which is definable. Let~$V'_{(x,y)}$ be such that~$V'\cap X\times \{(x,y)\}=V'_{(x,y)}\times \{(x,y)\}$ and let~$V_y$ be such that~$V_y\times\{x\}=V\cap \{y\}$. We have~$V'_{(x,y)}\cdot x=V_y$. By~\eqref{eq:equi4}, we have
\[
Supp(\nu_n)\cdot x_n\subseteq V_{y_n}= V'_{(x_n,y_n)}\cdot x_n.
\]
Therefore~$Supp(\nu_n)\subseteq V'_{(x_n,y_n)}$. 
We apply~\eqref{eq:equi3} to~$V'\subseteq X\times Y'$ and~$y'_n=(x_n,y_n)$. We have
\begin{equation}
\forall n\gg0, \dim\Bigl(Supp(\nu_\infty)\times\{(x_n,y_n)\}\,\cap\,V'\Bigr)=\dim(Supp(\nu_\infty)).
\end{equation}
In other terms,~$\dim(Supp(\nu_\infty)\cap V'_{(x_n,y_n)})=\dim(Supp(\nu_\infty))$, or
\[
\dim(Supp(\nu_\infty)\cap V_{y_n}\cdot x_n^{-1})=\dim(Supp(\nu_\infty)).
\]
We observe~$\dim(Supp(\nu_\infty)\cap V_{y_n}\cdot x_n^{-1})=\dim(Supp(\nu_\infty)\cdot x_n\cap V_{y_n})$. This implies~\eqref{eq:equi6} and concludes the proof of Theorem~\ref{thm:equi}.
\end{proof}

\subsection{Some equidistribution statements}\label{sec:31}

Consider~$K=\R^N/\Z^N$ and~$e\in\Z_{\geq1}$.
For a real Lie subgroup~$G\leq K$, let~$\nu_{G}$ be the Haar probability measure associated with~$G$, viewed as a probability measure on~$K$ concentrated on~$G$. For~$k\in K$, we denote by~$\nu_{G+k}$ the image of~$\nu_G$ by the translation~$k'\mapsto k+k':K\to K$. Let~$L_e:=\{\lambda^e|\lambda\in\wh{\Z}^{\times}\}$ act on~$K_{tors}=\Q^N/\Z^N$. For~$G\leq K$ a real Lie subgroup, and~$t\in K_{tors}$ and~$k\in K$ we define
\[
E(G,t,k)=\bigcup_{t'\in L_e\cdot t} G+t'+k
\]
and
\[
\mu_E:=\frac{1}{\#L_e\cdot t}\sum_{t'\in L_e\cdot t} \nu_{G+t'+k}
\]

The following implies that the set of measures of the form~$\mu_{E(G,t,k)}$ is a compact subset in the space of probabilities on~$K$.
\begin{theorem}\label{thm:equi}
Let~$G_i$ be a sequence of real Lie subgroups of~$K$, let~$t_i\in K_{tors}=\Q^N/\Z^N$ be
a sequence of torsion points, and let~$k_i$ be a sequence in~$K$.

Assume that the sequence~$\mu_{E(G_i,t_i,k_i)}$ has a limit~$\mu_\infty$ in the space of probabilities on~$K$.

Then there exists~$G_\infty\leq K$, and~$t_\infty\in K_{tors}$ and~$k_\infty\in K$ such that
\[
\mu_\infty=\mu_{E(G_\infty,t_\infty,k_\infty)}.
\]
Moreover, if the sequence~$k_i$ has a limit~$k_\infty$ in~$K$, then
\[
\mu_{E(G_i,t_i,0)}\to \mu_{E(G_\infty,t_\infty,0)}
\]
and
\[
\forall i\gg0, E(G_i,t_i,0)\subseteq E(G_\infty,t_\infty,0).
\]
\end{theorem}

Let~$\ast$ denote the convolution operation on probability measures on~$K$.
We observe that
\[
\mu_{E(G,t,k)}=
\mu_{E(G,0,0)}
\ast
\mu_{E(0,t,0)}
\ast
\mu_{E(0,0,k)}.
\]
Assume that~$\mu_{E(G_i,t_i,k_i)}$ has a limit~$\mu_\infty$. Since~$K$ is compact,
after possibly extracting a subsequence, there exist limits
\[
\alpha=\lim \mu_{E(G_i,0,0)}\qquad
\beta=\lim \mu_{E(0,t_i,0)}\qquad
\gamma=\lim \mu_{E(0,0,k_i)}.
\]
By continuity, we have~$\mu_\infty=\alpha\ast\beta\ast\gamma$. We note that~$
\alpha=\lim \mu_{E(0,0,k_i)}$ is the Dirac mass~$\delta_{k_i}$ and converges if and only if~$k_i$
has a limit~$k_\infty$ in~$K$. We have then
\[
\gamma=\delta_{k_{\infty}}=\mu_{E(0,0,k_\infty)}.
\]

By~\cite{RWeyl}, we have
\[
\beta=\mu_{E(B,t_\infty,0)}
\]
for some~$B\leq K$ and~$t_\infty\in K_{tors}$ such that
\[
\forall i\gg0, L_e\cdot t_i\subseteq B+L_e\cdot t_\infty.
\]

Moreover, we claim that, for some Lie subgroup~$G_{\infty}\leq K$, we have
\[
\alpha=\mu_{E(G_\infty,0,0)}\qquad\forall i\gg0, G_i\leq G_\infty.
\]
\begin{proof}Let~$X=Hom(K,U(1))\simeq \Z^n$. For~$\chi\in X$, let~$m_\chi(\nu_{G_i})=\int \chi(x) d\nu_{G_i}(x)$.
Let~$\Lambda_i:=\{\chi\in X|\chi(G_i)=\{1\}\}$. Then~$\chi\mapsto m_\chi(\nu_{G_i}):X\to \C$ is the indicator function~$1_{\Lambda_i}$ of~$\Lambda_i$. By Weyl criterion the function~$f:\chi\mapsto m_\chi(\alpha):X\to \C$ is the pointwise limit of~$1_{\Lambda_i}$. For~$\chi,\chi'\in X$ and~$i\in\Z_{\geq1}$ 
\[
m_\chi(\nu_{G_i})\in\{0;1\}.
\]
It follows that~$f(\chi)\in\{0;1\}$ and~$f(\chi)=1\Leftrightarrow\exists j_{\chi}\in\Z_{\geq1} \forall i\geq j_{\chi}, \chi\in\Lambda_i$.
Assume~$f(\chi)=f(\chi')=1$. For~$i\geq \max\{j_\chi;j_\chi'\}$, we have~$\chi',\chi\in\Lambda_i$.
Since each~$\Lambda_i$ is a group, we have
\[
\forall i\geq \max\{j_\chi;j_\chi'\}, \chi-\chi'\in \Lambda_i.
\]
We deduce~$f(\chi-\chi')=1$. Moreover, for~$\chi=1:X\to U(1)$, we have~$f(1)=1$. We deduce that~$\Lambda=\{\chi|f(\chi)=1\}$ is a subgroup of~$X$. Let~$G_\infty:=\{k\in K|\forall \chi\in\Lambda,\chi(k)=1\}$. Then
\[
m_\chi(\alpha)=f(\chi)=1_{\Lambda}(\chi)=m_\chi(\nu_{G_\infty}).
\]
By Weyl criterion we have
\[
\lim \mu_{E(G_i,0,0)}
=
\alpha
=\nu_{G_{\infty}}.
\]
Let~$\chi_1,\ldots,\chi_r$ be generators of~$\Lambda$. Then, for~$i\geq \max\{j_{\chi_1};\ldots;j_{\chi_r}\}$, we have
\[
\chi_1,\ldots,\chi_r \in \Lambda_i.
\]
Since~$\Lambda_i$ is a subgroup, we have~$\Lambda\leq \Lambda_i$. This implies~$G_i\leq G_{\infty}$.
\end{proof}


\section{On weak Mordell-Lang in geometric families}

\begin{theorem}\label{thm:weak}
Let~$A\to S$ be an abelian scheme over an algebraic variety over~$\C$ and~$e\in\Z_{\geq1}$. We assume that, 
\begin{equation}\label{eq:hypmonodromy}
\text{for some~$s\in S$, the monodromy action of~$\pi_1(S(\C))$ on~$\End(A_s)$ is trivial.}
\end{equation}
Let~$e\in\Z_{\geq1}$ and let~$L_e:=\{\lambda^e|\lambda\in\wh{\Z}^\times\}$ act on every commutative torsion group.

Let~$s_n$ be a generic sequence in~$S$ and  let~$A_n$ denote the fibre above~$s_n$. For every~$n$, let~$a_n\in (A_n)_{tors}$  be torsion point, let $G_n\leq A_n$  an algebraic subgroup and let~$k_n\in A_n$ be point of~$A_n$. Let
\[E_n:=G_n+L_e\cdot a_n+k_n
=\{g+\lambda^e\cdot a_n+k_n|g\in G_n, \lambda\in \wh{\Z}\}\subseteq A_n(\C)
\]
and let~$V$ be the Zariski closure of~$\bigcup E_n$ in~$A$.% Then for every irreducible component~$C$ of~$V$, there exists~$d\in\Z_{\geq1}$ such that~$[d](C)$ is an abelian subscheme of~$A\to S$.

Then for every irreducible component~$C$ of~$V$, there exists an abelian subscheme~$B\leq A$ such that~$C=C+B$
and~$\#E_n/B_{s_n}=O(1)$.
\end{theorem}


\subsection{Definable trivilialisation}
We recall the following.
\begin{theorem}[{\cite[Theorem~1.2]{PS}}]\label{thm:PS} Let~$\alpha:A\to S$ be an abelian scheme over an irreducible complex algebraic variety.

There exists a~$\R_{an,\exp}$-definable map~$\tau:A(\C)\to U(1)^{2g}$ such that for every~$s\in S$, the map~$\tau_s:A_s(\C)\to A(\C)\to U(1)^{2g}$ is a Lie group isomorphism. We say that~$\tau$ is a definable trivialisation of~$A\to S$
\end{theorem}
\begin{proof}Let~$S'\to S$ be a morphism of algebraic varieties, and let~$\alpha':A':=A\times_S S'\to S'$ be the abelian $S'$-scheme obtained by base change. 

We observe that~$\tau:A(\C)\to U(1)^{2g}$ is a definable trivialisation of~$A\to S$, then~$A'\to A\to U(1)^{2g}$ is a trivialisation of~$A'\to S'$. 

If~$S'\to S$ is surjective, then, by definable choice, there exists a definable section~$\sigma:S(\C)\to S'(\C)$ of~$S'(\C)\to S(\C)$. We note that~$A_s=A'_{\sigma(s)}$ by definition. If~$\tau':A'(\C)\to U(1)^{2g}$ is a trivialisation of~$A'\to S'$, then the map~$\tau:A(\C)\to U(1)^{2g}$ defined by~$\forall s\in S, \forall a\in A_s(\C), \tau(a)=\tau'(a)$ is a definable trivialisation of~$A\to S$.

Let~$S'\to S$ be an étale cover such that~$A'\to S'$ admits a polarisation and a full level structure of level~$N=3$.
Then there is a moduli map~$m:S'\to \mathbf{A}_{g,\psi}(N)$ such that~$A'\to S'$, the pullback of~$A$ by~$S'\to S$, is also the pullback of the universal abelian scheme~$\mathcal{A}\to \mathbf{A}_{g,\psi}(N)$. By the observation above, it is enough to treat the case where~$A\to S$ is~$\mathcal{A}\to \mathbf{A}_{g,\psi}(N)$.

The Theorem is then~\cite[Theorem~1.2]{PS}.
\end{proof}

We will prove the following.
\begin{theorem}\label{thm:locus}
Let~$A\to S$ and~$\tau:A\to U(1)^{2g}$ be as in Theorem~\ref{thm:PS}.
Let~$s_n$ be a generic sequence in~$S$ and~$B_n\leq A_n:=A_{s_n}$ be abelian subvarieties such that~$H:=\tau_{s_n}(B_n)\leq  U(1)^{2g}$ does not depend on~$n$. Assume~\eqref{eq:hypmonodromy}.

Then, possibly after extracting an infinite subsequence, the Zariski closure~$\ol{\bigcup B_n}^{Zar}$ is an abelian subscheme~$B\hookrightarrow A\to S$.
\end{theorem}
\subsubsection{}
We recall some facts while introducing some notations. Let~$\psi:\Z^{2g}\times\Z^{2g}\to \Z$ be a bilinear form such that its~$\Q$-linear extension~$\psi_\Q:\Q^{2g}\times\Q^{2g}\to \Q$ is a nondegenerate symplectic form. We denote~$\Gamma(N):=GSp(\psi_\Q)\cap \ker(GL(2g,\Z)\to GL(2g,\Z/(N))$ the principal congruence subgroups. Let~${H_g}$ be the upper Siegel half-space and~$X:=\pm {H}_g$. An element~$\tau\in{H}_g$ is a symmetric~$g\times g$ matrix~$\tau=\Re(\tau)+i\cdot \Im(\tau):\C^g\to \C^g$ over~$\C$, such that~$\Im(\tau)$ is the matrix of a positive definite quadratic form on~$\R^g$. To~$\tau\in \mathcal{H}_g$ we attach the morphism of abelian groups:
\begin{equation}\label{eq:iota}
\iota_\tau:(I_g,\tau):\Z^{2g}\to \C^g,
\end{equation}
where~$I_g$ is the injection~$\Z^{g}\leq \C^g$.
Then~$\Lambda_\tau:=\iota_\tau(\Z^{2g})$ is a lattice of~$\C^g$ and~$A_\tau:=\C^g/\Lambda_\tau$ is an abelian variety, and the symplectic form~$\Lambda_\tau\times \Lambda_\tau\to 2\pi i\Z$ induced by~$2\pi i\psi$ is a polarisation on~$A_\tau$. The spaces
\[
\mathbf{A}_{g,\psi}(N):=\Gamma(N)\sous{H}_g
\]
are connected components of a Shimura varieties, for the Shimura datum~$(GSp(\psi_\Q),\pm H_g)$, which are moduli spaces of abelian varieties with polarisation of type~$\psi$ 
and level structure of level~$N$, and are smooth for~$N\gg 0$.

Consider the action of~$\C^\times$ on~$\C^g$ by homotheties and consider the structure of~$\C^g$ as a~$\R$-linear vector space. It defines a~$\R$-linear representation~$\C^\times \to GL_\R(\C^g)\simeq GL(2g,\R)$. Consider the isomorphism~$\iota_\tau\tens\R:\Z^{2g}\tens\R\to \C^g$. We deduce conjugated representation
\[
\C^\times \to Aut(\Z^{2g}\tens\R)\simeq GL(2g,\R) 
\]
and this representations factors via some~$h_\tau:\C^\times\to GSp(\psi_\Q)(\R)$.

Choose an~$\R$-Lie groups isomorphism~$K:=U(1)^{2g}\to \R^{2g}/\Z^{2g}$. We can then identify~$\Z^{2g}\simeq H^{1}(K;\Z)$, and isomorphisms~$K\simeq A_\tau$ of~$\R$-Lie groups corresponds to isomorphisms~$\Z^{2g}\simeq H_1(A_\tau;\Z)\simeq \Lambda_\tau$.  Let~$A_{H_g}:=\sqcup_\tau A_\tau\to H_g$ be the structural morphism of the holomorphic family of complex tori~$(A_\tau)_{\tau\in H_g}$. There is a natural quotient map~$\C^g\times H_g\to A_{H_g}$ and the topology on~$A_{H_g}$ is the quotient topology. Observe that~$H^1(A_{H_g}/H_g;\Z)\simeq \sqcup_\tau \Lambda_\tau\hookrightarrow \C^g\times H_g\to H_g$ is an étale cover and that~$\bigsqcup_{\tau\in H_g}:\iota_\tau:\Z^{2g}\times H_g\to H^1(A_{H_g}/H_g;\Z)$ is a tivialisation of this étale cover. Note also that~$\lambda:L:=\bigsqcup Isom(\Z^{2g},H_1(A_\tau;\Z))\to H_g$ is also an étale cover, and even a (right)~$GL(2g,\Z)$-torsor over~$H_g$. Let~$D\subseteq {H_g}$ be a connected subset. Every section~$\sigma:D\to L$ of~$\lambda$ over~$D$ defines a family of isomorphisms~$(j_{\sigma}(\tau):K\simeq A_\tau)_{\tau\in D}$. The isomorphisms~$\iota_\tau$ induce a distinguished section~$\sigma_0$.
\paragraph{}\label{para:facts}

We use the following facts:
\begin{itemize}
\item the map 
\[
\bigsqcup_{\tau \in D} j_\sigma(\tau):K\times D\to \bigsqcup_{\tau \in D} A_\tau
\]
is an homeomorphism if and only~$\sigma$ is continuous.
\item if~$D$ is connected and~$\sigma$ is continuous, the map~$\sigma$ is unique up to the action of~$GL(2g,\Z)$. More precisely there exists~$\gamma\in GL(2g,\Z)$ such that~$\forall \tau\in S, j_\sigma(\tau)=j_{\sigma_0}(\tau)\circ\gamma$.
\end{itemize}

We record a last few facts. A connected~$\R$-Lie subgroup has Lie subalgebra~$\mathfrak{b}=H_1(B;\Z)\tens\R$, and that~$B$ is an abelian subvariety if and only if~$\mathfrak{b}$ is stable under the action of~$\C^\times$. When this is the case, the symplectic form form induced by~$\psi_\Q$ on~$H_1(B;\Z)\tens\Q$ is  nondegenerate. Finally, the stabiliser in~$GSp(\psi)_\Q$ of a non degenerate subspace are all Levi subgroups, and are in particular reductive groups. 

\subsubsection{}

Before proving Theorem~\ref{thm:locus}, let us first establish the following.
\begin{lemma}\label{lem:locus}
Let~$H_g$ be Siegel upper half-space and~$\iota_\tau$ be as in~\eqref{eq:iota}. Let~$\Lambda\leq \Z^{2g}$ be a subgroup. Let~$X'$ be the set of~$\tau\in H_g$ such that~$\iota_\tau(\Lambda)\tens\R\leq \C^g$ is a~$\C$-linear subspace.

Assume that~$X'\neq\emptyset$.
Then there exist a Shimura subdatum~$(L,X)\leq (GSp(2g),\pm H_g)$ such that~$X'=X\cap H_g$.
\end{lemma}
\begin{proof}[Proof of Lemma~\ref{lem:locus}]
Consider~$\tau\in X'$. Let~$\psi:\Z^{2g}\times \Z^{2g}\to \Z$ denote the symplectic form associated with~$GSp(2g)$.
Then~$\C^g/\iota_\tau(\Lambda)\leq \C^g/(\Z^g+\tau\times\Z^g)$ is an abelian subvariety. We deduce that~$\Lambda$ is non-degenerate for~$\psi$. It follows that the stabiliser~$L\leq GSp(2g)_\Q$ of~$\Lambda\tens\Q$ is a Levi subgroup. In particular it is a reductive group. 

Let~$h_\tau:\C^\times\to GSp(2g)$ be the Hodge cocharacter corresponding to~$\tau$. Then~$h(\C^\times)\leq L(\R)^+$.
By~\cite[Prop. 3.2]{CU}, the pair~$(L,L(\R)\cdot \tau)$ is a Shimura subdatum of~$(GSp(2g),\pm H_g)$.

Consider~$\tau'\in X'$. Then~$L(\R)^+\cdot \tau$ and~$L(\R)^+\cdot \tau'$ are two~$L(\R)^+$ orbits which are symmetric subspaces of~$H_g$. This implies that there exists~$z\in Z_{GSp(2g,\R)^+}(L)$ such that
\[
L(\R)^+\cdot \tau'=z\cdot L(\R)^+\cdot \tau=L(\R)^+\cdot z\cdot \tau.
\]
But~$\tau$ is fixed by~$Z_{GSp(2g,\R)^+}(h_\tau(\C^\times))\geq Z_{GSp(2g,\R)^+}(L)$.
Therefore
\[
L(\R)^+\cdot \tau'=L(\R)^+\cdot \tau.
\]
Since~$\tau'\in X'$ is arbitrary, we have~$X'\subseteq L(\R)^+\cdot \tau$. On the other hand, one sees easily that~$X'\subseteq  L(\R)^+\cdot \tau$. The Lemma follows.
\end{proof}
We can now prove~Theorem~\ref{thm:locus}.
\begin{proof}[Proof of Theorem~\ref{thm:locus}]
We claim that there exists a definable partition~$S=S_1\sqcup\ldots\sqcup S_k$ such that
\begin{enumerate}
\item \label{enu1} the maps~$A_{S_i}:=\bigsqcup_{s\in S_i} A_s\to A\to U(1)^{2g}$ are continuous;
\item \label{enu2} there exists a definable~$\wt{S_i}\subseteq H_g$ such that~$\wt{S_i}\to H_g\to S_i$ is an homeomorphism;
\item \label{enu3} each~$S_i$ is connected and simply connected.
\end{enumerate}
\begin{proof}[Proof of the claim] By~\cite[Ch. 9]{VdD} we can find~$S_i$ such that~\eqref{enu1} holds true. After refining the partition, we may assume~\eqref{enu2}. By~\cite[Ch. 8]{VdD}, we may refine the partition in such a way that the~$S_i$ are simplexes, and in particular are connected and simply connected. This proves~\eqref{enu3}. 
\end{proof}
For~$s\in S_i$, corresponding to~$\wt{s}\in \wt{S_i}$, we have an isomorphism~$A_{\wt{s}}=\C^g/\Lambda_{\wt{s}}\to A_s$.
We choose an isomorphism~$K:=U(1)^{2g}\simeq \R^{2g}/\Z^{2g}$. We can identify the Lie algebra of~$U(1)^{2g}$ with~$\R^{2g}$ and~$H^1(K;\Z)$ with~$\Z^{2g}\leq\R^{2g}$. The map~$A_{\wt{s}}\to A_s\to A\to U(1)^{2g}$ induces a Lie algebra isomorphism
\[
\C^g\to \R^{2g}
\]
which induces an isomorphism~$j_{\wt{s}}:\Lambda_\wt{s}$ to~$\Z^{2g}$. We claim that there exists~$\phi\in Aut(U(1)^{2g})$ such that
\[
\forall \wt{s}\in \wt{S}, \iota_\wt{s}\circ\phi={j_{\wt{s}}}^{-1}
\]
\begin{proof}[Proof of the claim] See~\ref{para:facts}.
\end{proof}
Extracting a subsequence, we may assume that, for some~$i\in \{1;\ldots;k\}$, we have
\[
\forall n\in\Z_{\geq0}, s_n\in S_i.
\]
Let~$H$ as in the statement of Theorem~\ref{thm:locus}.
Let~$\Lambda=H^1(H;\Z)\leq \Z^{2g}=H^1(K;\Z)$ be the subgroup corresponding to~$H$.

Let~$X'\subseteq H_g$ be as in Lemma~\ref{lem:locus}. By assumption~$H=\tau_{s_n}(B_n)$ for an abelian 
subvariety~$B_n\leq A_{s_n}$. This implies~$\wt{s_n}\in X'$. Therefore,~$s_n$ is in the image~$Z$ of~$X'$ by~$H_g\to \mathbb{A}_g$. Lemma~\ref{lem:locus} implies that~$Z$ is a special subvariety of~$\mathbb{A}_g$. In particular~$Z\cap S$ Zariski closed in~$S$.

But~$s_n$ is a generic sequence in~$S$. Our extracted subsequence is thus generic in~$S$. This implies that~$\{s_n\}\subseteq  Z\cap S$ is Zariski dense in~$S$. As~$Z$ is closed, we have~$S\subseteq Z$.

We can deduce that~$\wt{S_i}$ is contained in~$X'$. By definition of~$X'$, the image of
\[
H\cap X'\to (\R^{2g}/\Z^{2g})\times X'\to \bigsqcup_{\tau \in X'} \C^{g}/\Lambda_\tau
\]
is an holomorphic family of abelian varieties. We deduce that its image in the universal abelian variety~$\mathcal{A}_g|_Z$ restricted to~$Z$ is an abelian sub-$Z$-scheme. Let~$B\to S$ be the restriction to~$S$ of this abelian subscheme. By construction~$B_n$ is equal to the fibre~$B_{s_n}$.

We deduce that~${\bigsqcup B_{s_n}}$ is the inverse image of the dense subset~$\{s_n\}$ by the map~$B\to S$.
Therefore
\[
\ol{\bigsqcup B_{s_n}}^{Zar}=\ol{B}^{Zar}=B.\qedhere
\]
\end{proof}
\subsection{Proof of Theorem~\ref{thm:weak}} Let~$a_n$,~$G_n$,~$k_n$,~$e$,~$E_n$ be as in the statement  of Theorem~\ref{thm:weak}, and let~$C$ be an irreductible component of the Zariski closure~$\ol{\bigcup E_n}^{Zar}$ in~$A$. Then, extracting a subsequence, we may assume that there exists~$c_n$ is a generic sequence in~$C$ such that~$c_n\in E_n\cap C$ for every~$n$. This implies that for every extraction~$\phi:\Z_{\geq0}\to \Z_{\geq0}$, the variety~$C$ is a component of~$\ol{\bigcup E_{\phi(n)}}^{Zar}$.

For every~$n$, let~$B_n\leq B'_n\leq A_n$ be the biggest abelian subvariety such that~$B'_n+E_n\subseteq \ol{\bigcup E_n}^{Zar}$. Substituting~$G_n$ with~$G_n+B'_n$, we may assume that~$B_n:=G^0_n$ is identical to~$B'_n$.

Let~$a'_n\in (A_n)_{tors}\cap (G_n+L_e\cdot a_n)$ be a torsion point of minimal order. We have
\[
G_n+L_e\cdot a_n=G_n+L_e\cdot a'_n
\]
and without loss of generality we may assume~$a_n=a'_n$.

Let~$\tau:A(\C) \to U(1)^{2g}$ be a definable trivialisation as in Theorem~\ref{thm:PS}. Let us define~$H_n:=\tau(G_n)$ and~$t_n:=\tau(a_n)$ and~$\kappa_n:=\tau(k_n)$. Then, with the notations of~\S\ref{sec:31}, we have
\[
\tau(E_n)=E(H_n,t_n,\kappa_n)
\]
and we can consider
\[
\mu_n:=\mu_{E(H_n,t_n,\kappa_n)}\text{ and }
\nu_n:=\mu_{E(H_n,t_n,0)}.
\]
Extracting a subsequence, we may assume that~$\kappa_n$ has a limit~$\kappa_{\infty}$ in~$U(1)^{2g}$ 
and that the sequence of probability measures~$\nu_n$ has a limit~$\nu_{\infty}$. By Theorem~\ref{thm:equi}
we have
\[
\nu_\infty=\mu_{E(H_\infty,t_\infty,0)}
\]
with~$t_\infty\in U(1)_{tors}^{2g}$ and~$H_\infty\leq U(1)^{2g}$ a real Lie subgroup such that
\begin{equation}\label{eq:truc}
\forall n\gg0, 
E(H_n,t_n,0)\subseteq E(H_\infty,t_\infty,0):= H_\infty +L_e\cdot t_{\infty}.
\end{equation}
Let~$V:=\ol{\bigcup E_n}^{Zar}$. Recall that~$C$ is an irreducible component of~$V$. Let~$a\mapsto s(a)$ denote the structural map~$A\to S$. Let~$V'\subseteq S(\C)\times U(1)^{2g}\times U(1)^{2g}$
be the image of~$V\times U(1)^{2g}$ by
\[
(v,g)\mapsto (s(v),g, \tau(v)-g).
\]
Note that~$V'\subseteq S(\C)\times U(1)^{2g}\times U(1)^{2g}$ is a definable subset. For every~$n\in\Z_{\geq0}$, we deduce from~$E_n\subseteq V(\C)$ that
we have
\begin{equation}\label{eq:ttruc}
\{s_n\}\times \{\kappa_n\}\times E(H_n,t_n,0)\subseteq V'.
\end{equation}
We apply Theorem~\ref{thm:equi0} with~$Y=S(\C)\times U(1)^{2g}$ and~$X=S(\C)\times U(1)^{2g}$ and~$y_n=(s_n,\kappa_n)$ and~$V\subseteq X\times Y$ corresponding to~$V'$ above. Then the assumptions~\eqref{eq:equi1} and~\eqref{eq:equi2} are satisfied by~\eqref{eq:truc} and~\eqref{eq:ttruc}. Note that~$\nu_\infty$ is a finite linear combination of translates of the Haar measure on~$H_\infty$, and thus the assumption~$\dim(A)<\dim(H_\infty)\Rightarrow \nu_\infty(A)=0$ is satisfied.
 
We deduce from~\eqref{eq:equi6} that, for~$n\gg0$, the set~$V'$ contains a non empty subset of~$\{s_n\}\times \{\kappa_n\}\times E(H_\infty,t_\infty,0)$ of dimension~$\dim(H_\infty)$. By construction
\[
V'\cap \left( \{s_n\}\times \{\kappa_n\}\times U(1)^{2g}\right)=\{s_n\}\times \{\kappa_n\}\times \tau((V\cap A_n)-k_n).
\]
Let~$G'_n\leq A_n(\C)$ be the real Lie subgroup such that~$\tau(G_n)=H_\infty$. 
Using the definable bijection
\[
(V\cap A_n)-k_n\simeq V'\cap \left(\{s_n\}\times \{\kappa_n\}\times U(1)^{2g}\right)
\] 
we deduce that~$V\cap A_n$ contains a subset of~$E'_n:=G'_n+L_e\cdot a_n+k_n$ of dimension~$\dim(H_\infty)=\dim(G_n)=\dim(E'_n)$. As in Theorem~\ref{thm:equi0}, we deduce that
\[
V\cap A_n
\]
contains, for the archimedean topology, an non empty open subset of~$E'_n$. As~$V\cap A_n$ is a Zariski closed subset,
we deduce that~$V\cap A_n$ contains a component of~$\ol{E'_n}^{Zar}=\ol{G'_n}^{Zar}+L_e\cdot a_n+k_n$.

Using similar arguments, one proves that~$V\cap A_n$ contains every component of~$\ol{G'_n}^{Zar}+L_e\cdot a_n+k_n$.

By our assumption~$B_n=B_n'$, we have 
\[
B_n=(\ol{G'_n}^{Zar})^0.
\]

We apply Theorem~\ref{thm:locus} to~$H=H^+_\infty$. Then, after possibly extracting a subsequence, we deduce that
\[
\ol{\bigcup B_n}
\]
is an abelian subscheme~$B$ of~$A\to S$.

Possibly extracting a subsequence, we may assume that~$c_n\in C$ does not belong to the intersection of~$C$ with the union of the other irreducible components of~$V$. Thus the inclusion~$c_n+B_n\subseteq V$ and the irreducibility of~$c_n+B_n$ imply that~$c_n+B_n\subseteq C$.

By Proposition~\ref{prop:truc} below, this proves~$C+B=C$.

\begin{proposition}\label{prop:truc} Let~$C$ be a subvariety of an abelian scheme~$A\to S$ and let~$B\leq A$ be an abelian subscheme. Let~$s_n$ and~$c_n$ be a generic sequence in~$S$ and in~$C$ with~$c_n$ in the fibre~$C_n:=C_{s_n}$. Assume 
\[
\forall n\geq 0, c_n+B_{s_n}\subseteq C.
\]
Then
\[
\forall s\in S, \forall c\in C_s, c+B_s\subseteq C_s.
\]
In other terms
\[
C=C+B.
\]
\end{proposition}
\begin{proof}
Let~$a:A\to S$ denote the structural map. It is enough to prove that the set
\[
\Sigma:=\{c\in C| c+B_{a(c)}\subseteq C\}
\]
describes a closed subvariety of~$C$: from~$\{c_n\}\subseteq \Sigma$, we conclude~$C=\ol{\{c_n\}}^{Zar}\subseteq \ol{\Sigma}^{Zar}\subseteq \ol{C}^{Zar}=C$.

For~$N\in \Z_{\geq0}$, let
\[
\Sigma_N:=\{c\in C| c+B_{a(c)}[N]\subseteq C\}.
\]
Recall that~$\bigcup_{N\in\Z_{\geq1}}B_{a(c)}[N]$ is dense in~$B_{a(c)}$. We deduce that for every~$c\in C$,
\[
\forall N\in\Z_{\geq0}, c+B_{a(c)}[N]\subseteq C\Rightarrow c+B_{a(c)}\subseteq C.
\]
It follows that~$\Sigma=\bigcap_{N\in\Z_{\geq0}} \Sigma_N$. It is thus enough to prove that~$\Sigma_N$ describes a closed subvariety of~$C$.


Recall that~$S':=B[N]$ is an étale cover of~$S$. Consider the base change~$C',B'\subseteq A'\to S'$ of~$C,B\subseteq A\to S$. Then~$B'[N]$ is a trivial cover of~$S'$: there exists sections~$\sigma_1,\ldots, \sigma_{N^{2\dim(B/S)}}:S'\to B'$ such that~$B'[N]=\bigcup_{1\leq i \leq N^{2\dim(B/S)}} \sigma_i(S')$.

Note that~$\Sigma_N$ is defined in terms of fibers. Let~$\Sigma'_N$ be the analogous set for~$C'$. Then~$\Sigma'_N$ is the pullback of~$\Sigma_N$, and it is enough to prove that~$\Sigma'_N$ describes a closed subvariety of~$C'$. 

Let~$i\in\{1;\ldots;N^{2\dim(B/S)}\}$ be arbitrary.
Note that the fibrewise sum~$C'+\sigma_i(S')$ is the image of the fibre product~$C'\times_{S'}\sigma_i(S')$ by the regular addition map~$A'\times_{S'}A'\to A'$. 
Recall that latter is a proper map, and that~$C'$ and~$S'$ are closed subvarieties of~$A'$. It follows that~$C'\times_{S'}\sigma_i(S')$ and~$C'+\sigma_i(S')$ are closed subvarieties of~$A'\times_{S'}A'$ and~$A'$.

One can check fibrewise that
\[
\Sigma'_N=\bigcap_{1\leq i \leq N^{2\dim(B/S)}}  C'+\sigma_i(S').
\]
This is an intersection of closed subvariety, and thus is itself a closed subvariety. The statement follows.\end{proof}

We have proved Proposition~\ref{prop:truc}. 

In order to prove Theorem~\ref{thm:weak}, it remains to prove
\[
\# E_n/B_{s_n}\leq O(1)\text{ as }n\to +\infty.
\]
Let~$c$ be the number of connected components of~$H_\infty$. We have
\begin{multline}
\# \left(E_n+B_{s_n}\right)/B_{s_n}
=\#\left(\tau(E_n)+\tau(B_{s_n})\right)/\tau(B_{s_n})\\
= \#\left(E(H_n,\tau_n,\kappa_n)+H^+_{\infty}\right)/H^+_{\infty}\\
\leq\#H_\infty/H_\infty^+\cdot \left(\#E(H_n,\tau_n,\kappa_n)+H_{\infty}\right)/H_{\infty}\\
=c\cdot \#E(H_\infty,t_\infty,\kappa_n)/H_{\infty}\\\leq c\cdot ord(t_\infty+H_{\infty}/H_{\infty}).
\end{multline}
This finishes the proof of Theorem~\ref{thm:weak}.
%\subsection{Proof of Theorem~\ref{thm:weak}}
%
%
%For every~$n$, let~$B_n\leq B'_n\leq A_n$ be the biggest abelian subvariety such that~$B'_n+E_n\subseteq V$.
%Substituting~$B_n$ with~$B'_n$, resp.~$G_n$ with~$G_n+B'_n$, we may assume that~$B_n=B'_n$, resp.~$G_n^0=B'_n$.
% 
%Let~$a'_n\in (A_n)_{tors}\cap (B'_n+L\cdot a_n)$ be a torsion point of minimal order.
%
%Without loss of generality, we substitute~$B_n$ with~$B'_n$ and~$a_n$ with~$a'_n$.
%
%Let~$H_n=\tau(B_n(\C))$, let~$f_n=\tau(a_n)$ and~$g_n=\tau(k_n)$ and let~$\mu_n$ be the~$H_n$ invariant probability measure on~$F_n:=H_n+L_e\cdot f_n+g_n$ such that~$\mu_n(H_n+x)=\mu_n(H_n+y)$ for every~$x,y\in F_n$. Passing to a subsequence we may assume~$\mu_n$ has a tight limit~$\mu_\infty$ and~$g_n$ has a limit~$g_\infty$ in~$U(1)^{2g}$. 
%Then (cf. \cite{}), there exists~$H_\infty\leq U(1)^{2g}$ and~$f_\infty \in (U(1)^{2g})_{tors}$ such that~$\mu_\infty$ is the~$H_\infty$-invariant probability measure on~$F_\infty:=H_\infty+L_e\cdot f_\infty+g_\infty$ such that~$\mu_\infty(H_\infty+x)=\mu_\infty(H_\infty+y)$ for every~$x,y\in F_\infty$. Moreover, we have
%\[
%\forall n\gg0, H_n+L_e\cdot f_n\subseteq H_\infty+L_e\cdot f_\infty.
%\]
%Let~$V'$ be the image of~$V(\C)$ by~$A(\C)\to S(\C)\times U(1)^{2g}$ and let~$F'_n=\{s_n\}\times F_\infty$. This are~$\R_{an,\exp}$-definable sets. By~\cite{}, we have
%\[
%\forall n\gg0, \dim(F'_n\cap V') =\dim(F_\infty), 
%\]
%and for such~$n\gg0$, the set~$V'$ contains an open subset of~$F'_n$. [...]
%
%Let~$V'':=\sqcup_{g\in  U(1)^{2g}}(V'-g)\times\{g\}\subseteq A(\C)\to S(\C)\times U(1)^{2g}\times U(1)^{2g}$.
%

%\section{On Hodge locus and quasi abelian subschemes}
%
%Let~$\alpha: A\to S$ be a polarised abelian scheme over a connected algebraic variety over~$\C$ and let~$s\in S(\C)$. Then
%the goup~$\Gamma:=\pi_1(S(\C),s)$ acts on~$L:=H_1(A_s(\C);\Z)$ and on~$\End(L)$.
%The polarisation induces an euclidean structure on~$\End(L)$, which is invariant under the action of~$\Gamma$.
%Since the orthogonal group of~$\End(L)\tens\R$ is compact, and~$GL(L,\Z)$ is discrete in~$GL(L)(\R)$, the image of~$\Gamma$ in~$GL(L,\Z)$ is finite. Let~$\wt{S}\to S$ be the corresponding Galois covering. Then the pullback of~$\wt{A}\to\wt{S}$ is such that, for any~$\wt{s}\in \wt{S}(\C)$, the action of~$\pi_1(\wt{S}(\C),\wt{s})$ on~$\End(H_1(A_s(\C);\Z)$ is trivial.
%
%\begin{proposition}
%Let~$A\to K:=U(1)^{2g}$ be a definable trivialisation, let~$H\leq K$ be a connected real Lie subgroup, and let~$s_n$ be a generic sequence in~$S$. Let~$B_n$ be the inverse image of~$H$ by~$A_n:=A_{s_n}\to A\to K$. 
%Assume that~$B_n$ is an abelian subvariety of~$A_n$ for every~$n$.
%
%\begin{enumerate}
%\item Assume~$S=\wt{S}$. Then the Zariski closure of~$\bigcup B_n$ is an abelian subcheme~$B\leq A\to S$ such that~$B_n=B_{s_n}$ for every~$n$.
%\item In general,  the Zariski closure of~$V:=\ol{\bigcup B_n}^{Zar}$ is the image of an abelian subcheme~$\wt{B}\leq \wt{A}\to \wt{S}$ and is such that the connected component of~$V_s$ is an abelian subvariety and its images by monodromy [...]
%\end{enumerate}
%\end{proposition}
%
%There is a finite partition~$S(\C)=S_1\sqcup\ldots \sqcup S_k$ such that for~$S'\in\{S_1;\ldots; S_k\}$, this~$S'$ is connected and simply connected, and the trivialisation induces a homeomorphism.
%\[
%A_{S'}\to S'\times K.
%\]
%Extracting a subsequence, we may assume~$s_n\in S'$ for every~$n$.
%\begin{lemma}
%For every~$s\in S'$, the inverse image of~$H$ by~$A_{s'}\to A\to K$ is an abelian subvariety~$B_{s'}\leq A_{s'}$.
%\end{lemma}
%\begin{proof}
%Let~$S\to A_g$ be the moduli map. As~$S'$ is connected and simply connected, there is a lift~$S'\to X_g$ such that~$S'\to S\to A_g$ equals~$s'\to X_g\to A_g$. Every~$x\in X_g$ is a Hodge cocharacter~$x:\C^\times \to \End(\Z^{2g})$, where one can identify~$K$ with~$\R^{2g}/\Z^{2g}$. Let~$\mathfrak{h}\leq \R^{2g}$ be the Lie algebra of~$H\leq K$. We consider
%\[
%Y:=\{x\in X| x(\C^\times)\cdot \mathfrak{h}=\mathfrak{h}\}.
%\]
%and~$G=Stab_{GSp(2g)}(\mathfrak{h})$.
%We claim
%\[
%(G,Y)\text{ is a Shimura subdatum of }(GSp(2g),X_g).
%\]
%\begin{proof} .[...].
%\end{proof}
%Consequently, the image~$Z$ of~$Y$ in~$\mathbf{A}_g$ is a special subvariety, particularly an algebraic subvariety. The inverse image~$Z'$ of~$Z$ in~$S$ is thus a Zariski closed subvariety. By ..., the generic sequence~$s_n$ is contained in~$Z'$. We have thus ~$Z'=S$.
%
%[...]
%
%We deduce that the image of~$S'\to X_g$ is contained in~$Y$. QED
%\end{proof}
%\begin{lemma}There is an (quasi)  abelian subscheme~$B\leq A\to S$ such that for every~$s'\in S'$, ...
%\end{lemma}
%
%\begin{proposition}Let~$H_n\leq K$ be a sequence of connected real Lie subgroups, and let~$B_n$ be the inverse  of~$H_n$ by~$A_n:=A_{s_n}\to A\to K$. (Assume~$s_n$ generic) Assume that~$V:=\ol{\bigcup B_n}^{Zar}$ is irreducible.
%
%Then~$V$ is a quasi-abelian subscheme.
%?Assume that~$B_n$ is an abelian subvariety of~$A_n$ for every~$n$.
%\end{proposition}


\section{Galois actions associated with a hybrid orbit of abelian varieties}\label{sec:Galois}
The main results of this section are Theorem~\ref{thm:Serre hybrid} and Theorem~\ref{thm:uniform Kummer1}.
Of independent interest is the reformulation Theorem~\ref{thm:uniform Faltings} of Falting's results
on Tate isogeny conjecture.

The strategy of this article to prove the main theorems is to find subset of Galois
orbits which are amenable to equidistribution results. We use two cases.
\begin{itemize}
\item The method of~\cite{RWeyl} applies to torsion points and uses Serre's Theorem~\ref{thm:Serre eAK}
on Lang's conjecture. 
\item The Galois action on division points of rational points which are not torsion points is the subject of Kummer theory for abelian varieties.
\end{itemize}
In both case we need a uniform results as our abelian variety varies in a hybrid orbit.
For Serre's theorem, this is Theorem~\ref{thm:Serre hybrid}, which relies on~\cite[App. B]{RY:IHES}'s approach to Serre's theorem. For Kummer theory, we follow~\cite{Hindry}, adding uniformity with respect to passing to a finite extension~$E/K$
and the maximal abelian extension~$E^{ab}/K$.
% A fundamental remark is that Proposition~\ref{} works with~$(\rho_A,r)$, whereas Hindry and Ribet considered~$\rho_A$.


\subsection{Tate module}
Let~$A$ be an abelian variety over an extension~$K/\Q$ and let~$\ol{K}$ be an algebraic closure of~$K$. Denote by~$A_{tors}\leq A(\ol{K})$ the subgroup of torsion points.
The~$\wh{\Z}$-Tate module of~$A$ is then
\[
\wh{T}(A):=\Hom(\Q/\Z,A_{tors})\approx\wh{\Z}^{2\dim(A)}.
\]
The action of~$\Gal(\ol{K}/K)$ on~$A_{tors}$ induces a~$\wh{\Z}$-linear representation
\[
\rho_A:\Gal(\ol{K}/K)\to \Aut(\wh{T}(A))\approx GL(2\dim(A),\wh{\Z}).
\]
\subsubsection{Several theorems}We rely on the following.
\begin{theorem}[Mordell-Weil-Néron, {\cite[p. 52 Remark]{Serre:MWT}}, \cite{LN}]\label{thm:MWN}
Let~$A$  be an abelian variety over a finitely generated extension~$K$ over~$\Q$ or~$\F_p$.
Then~$A(K)$ is a finitely generated abelian group.
\end{theorem}

We observe that~$GL(1,\wh{\Z})$ acts naturally on any~$\wh{\Z}$-module, in particular torsion and any profinite abelian groups: $A_{tors}$, $A[N]$, $\wh{T}(A)$, \textit{etc}.

\begin{theorem}[Serre (see~{\cite[33, p.\,34, Th. 2 and 2']{Serre:O4}}, cf.~{\cite[Lemme 12]{Hindry}})]\label{thm:Serre eAK}
Let~$A\neq0$  be an abelian variety over a finitely generated extension~$K/\Q$. 
There exists~$e=e(A,K)\in\Z_{\geq1}$ such that
\begin{equation}
\label{eq:Serre e}
\forall \lambda\in GL(1,\wh{\Z}), \lambda^e\in \rho_A(Gal(\ol{K}/K)).
\end{equation}
\end{theorem}

The following is a reformulation of~\cite[Definition 2.1]{RY:IHES}.
\begin{theorem}[Uniform integral Tate conjecture, Faltings]\label{thm:uniform Faltings}
Let~$A$  be an abelian variety over a finitely generated extension~$K/\Q$ such that~$\End_A:=\End(A/K)=\End(A/\ol{K})$.
%\begin{enumerate}
%\item
%For every~$D\in\Z_{\geq1}$, there exists~$\mu(D)$ such that for every extension~$E/K$ of degree~$[E:K]\leq d$, and every~$\ell\geq\mu(D)$
%\begin{enumerate}
%\item the commutant, in~$\End_{\F_\ell}(A[\ell])$ of~$\F_\ell[\Gal({\ol{K}}/E)]$ is the image of~$\End_A\to \End_{\F_\ell}(A[\ell])$;
%\item the~$\Gal(\ol{K}/E)$-module~$\End_{\F_\ell}(A[\ell])$ is semisimple.
%\end{enumerate}
%If~$\ell\geq \mu(D)$, then for every~$\Gal({\ol{K}}/E)$-invariant $\F_\ell$-vector subspace~$M\leq A[\ell]$ there exists~$\alpha\in \End_A$ such that~$M=\ker(\alpha)=(Id-\alpha)(A[\ell])$.
%\item
%For every~$\ell$, and every finite extension~$E/K$,
%\begin{enumerate}
%\item the subalgebra~$\Z_\ell[\Gal({\ol{K}}/E)]\leq \End(\wh{T}(A))\tens \Q_\ell$ is an order in the commutant of~$\End_A\tens \Q_\ell$ in~$\End(\wh{T}(A))\tens \Q_\ell$.
%\item the~$\Gal(\ol{K}/E)$-module~$\wh{T}(A)\tens\Q_\ell$ is semisimple.
%\end{enumerate}
%Morevoer~$[\Z_\ell[\Gal({\ol{K}}/K)]\cap \ell...:\Z_\ell[\Gal({\ol{K}}/E)\cap \ell...]\leq ?([E:K])\cdot \#\F_\ell[\Gal({\ol{K}}/K)].$
%\item 

Let~$Z$ be the commutant of~$\End(A)$ in~$\End_{\wh{\Z}}(\wh{T}(A))$. Then for every~$D\geq 1$, there exists~$F(A,K,D)$ such that for every extension~$E/K$ of degree~$[E:K]\leq D$,
\begin{equation}\label{eq:TateUnifLinear}
\left[Z:\wh{\Z}\left[\rho_A(Gal(\ol{E}/E))\right]\right]\leq F(A,K,D)<+\infty.
\end{equation}
\end{theorem}
For~$K=E$ a number field,~\eqref{eq:TateUnifLinear} is~\cite[Théorème 2.7]{Deligne}.
This is also true when~$K$ is a finitely generated extension of~$\Q$.

For every prime~$\ell$, let~$F_\ell(A,E):=\left[Z\tens\Z_\ell:\Z_\ell\left[\rho_A(Gal(\ol{E}/E))\right]\right]$. We have
\[
\left[Z:\wh{\Z}\left[\rho_A(Gal(\ol{E}/E))\right]\right]=\prod_\ell F_\ell(A,E)=F(A,E,1).
\]
We will prove that there exists~$\ell(A,K,D)$ such that for every extension~$E/K$ of degree~$[E:K]\leq D$,
and every prime~$\ell\geq \ell(A,K,D)$, we have
\begin{equation}\label{eq:Tateclaim1}
F_\ell(A,E)=1
\end{equation}
We will also prove that, for every prime~$\ell$,
\begin{equation}\label{eq:Tateclaim2}
\sup_{E/K, [E:K]\leq D}F_\ell(A,E)<+\infty,
\end{equation}
where~$E$ ranges through the extensions of~$K$ of degree at most~$D$. Theorem~\ref{thm:uniform Faltings} follows from~\eqref{eq:Tateclaim1} together with \eqref{eq:Tateclaim2}.


Let us prove~\eqref{eq:Tateclaim2}. 
\begin{proof}{Proof of \eqref{eq:Tateclaim2}}Let~$\rho_{A,\ell}:Gal(\ol{E}/K)\to GL(\widehat{T}(A)\tens\Z_\ell)$ denote the natural representation on the~$\Z_\ell$-linear Tate module. We note that~$F_\ell(A,E)$ only depends on the subgroup~$\rho_{A,\ell}(Gal(\ol{E}/E))\leq \rho_{A,\ell}(Gal(\ol{E}/K))$. But the~$\ell$-adic group~$\rho_{A,\ell}(Gal(\ol{E}/K))$ has only finitely many closed subgroups of index at most~$D$.  Moreover
\begin{multline*}
\left[\rho_{A,\ell}(Gal(\ol{E}/E)): \rho_{A,\ell}(Gal(\ol{E}/K))\right]\leq \left[Gal(\ol{E}/E): Gal(\ol{E}/K)\right]
\\=[E:K]\leq D.
\end{multline*}
Therefore, in~\eqref{eq:Tateclaim2} the~$F_\ell(A,E)$ takes only finitely many values as~$E$ ranges through the extensions of~$K$ of degree at most~$D$. Note also that~$F_\ell(A,E)\leq F(A,E,1)<+\infty$. We deduce the finiteness in \eqref{eq:Tateclaim2}.
\end{proof}

\begin{proof}{Proof of \eqref{eq:Tateclaim1}}
For every prime~$\ell$, let~$\rho_{A[\ell]}$ be the natural~$\F_\ell$-linear representation~$Gal(\ol{K}/K)\to GL_{\wh{\Z}}(T(A))\to GL_{\F_\ell}(A[\ell])$.
If~$\ell \nmid F_\ell(A,K)$, then the image of~$Z$ in~$\End(A[\ell])$ is~$\F_\ell[\rho_{A[\ell]}(Gal(\ol{K}/K))]$.

Consider an embedding~$K\to \C$ and the corresponding~$\Z$-structure~$\Lambda:=H^1(A(\C);\Z)\tens\wh{\Z}\simeq T(A)$ with respect to some embedding~$\ol{K}\to \C$. Then~$R:=\End(A,\ol{K})$ embeds in~$\End_\Z(\Lambda)$ and its commutator~$Z'$ satisfies~$Z'\tens \wh{\Z}\simeq Z$. The algebras~$R\tens\Q$ and~$Z\tens\Q$ are semisimple. For~$\ell>\ell(A)$, the morphisms~$R\tens \F_\ell\to \End(A[\ell])$ and~$Z'\tens\F_\ell\simeq Z\tens\F_\ell\to\End(A[\ell])$ are injective; the images are semisimple~$\F_\ell$-algebras; and the images are the commutator of each other in~$\End(A[\ell])$.

Let~$G_\ell:=\rho_{A[\ell]}(Gal(\ol{K}/E))$.
For~$\ell\geq \ell(A,K,[E:K])$, \cite[Def. 2.1]{RY:IHES} implies that~$A[\ell]$ is a semisimple~$G_\ell$-module, and that~$\End(A)\tens\F_\ell$ is the algebra of~$G_\ell$-equivariant endomorphisms of~$\End(A)\tens\F_\ell$. This implies that~$\F_\ell[\rho_{A[\ell]}(Gal(\ol{K}/K))]=Z\tens\F_\ell=\F_\ell[\rho_{A[\ell]}(Gal(\ol{K}/E))]$. 

Recall that~$\rho_{A,\ell}(Gal(\ol{K}/K))\leq Z\tens\Z_\ell$. By Nakayama's lemma, we conclude~$\Z_\ell[\rho_A(Gal(\ol{K}/E))]=Z\tens\Z_\ell$ for~$\ell\geq \ell(A,K,[E:\Q])$.
\end{proof}
We have proved~\eqref{eq:Tateclaim1} and \eqref{eq:Tateclaim2}.
This concludes the proof of Theorem~\ref{thm:uniform Faltings}.
%
%\subsubsection*{Small primes}
%Let~$\ell$ be an arbitrary prime. Let~$K_\ell:=K(A[\ell])$. Let~$\rho_{A,\ell}$ be the action of~$Gal(\ol{K}/K)$ on~$T(A)\tens\Z_\ell$. We have~$[K_\ell:K]\leq \#GL(A[\ell])\leq \ell^{\dim(A)^2}$ and
%\[
%\rho_{A,\ell}(Gal(\ol{K}/K_\ell))=\ker(\rho_{A,\ell}(Gal(\ol{K}/K_\ell))\to GL(A[\ell]))\leq 1+\ell\End_{\Z_\ell}(T(A)\tens\Z_\ell)
%\]
%is a pro~$\ell$-group. For~$\ell=2$ one adapts the 
%
%Therefore
%\[
%[\rho_{A,\ell}(Gal(\ol{K}/K_\ell)):\rho_{A,\ell}(Gal(\ol{K}/E\cdot K_\ell))]
%\]
%is a power of~$\ell$ that divides~$[E\cdot K_\ell:K_\ell]$, which divides~$[E:K]$. Thus, for~$\ell\nmid[E:K]$, we have
%\(
%\rho_{A,\ell}(Gal(\ol{K}/K_\ell))=\rho_{A,\ell}(Gal(\ol{K}/E\cdot K_\ell))
%\)
%and therefore
%\[
%\Z_\ell[\rho_{A,\ell}(Gal(\ol{K}/E\cdot K_\ell)]= \Z_\ell[\rho_{A,\ell}(Gal(\ol{K}/K_\ell))].
%\]
%Assume~$\ell\nmid[E:K]$. Then~$$\rho_{A,\ell}(Gal(\ol{K}/E\cdot K_\ell))\leq \rho_{A,\ell}(Gal(\ol{K}/K_\ell))$$ and
%\[
%X:=\Z_\ell[\rho_{A,\ell}(Gal(\ol{K}/E\cdot K_\ell)]\leq Y:=\Z_\ell[\rho_{A,\ell}(Gal(\ol{K}/K_\ell))].
%\]
%Therefore
%\[
%[Z\tens\Z_\ell:Y]\mid [Z\tens\Z_\ell:X]\mid [Z:\ell Z]\cdot [\Z_\ell[\rho_{A,\ell}(Gal(\ol{K}/K)]].
%\]
%
%Assume~$[E:K]=\ell$ and that~$\ell$ is odd. Let~$g_1,\ldots,g_k\in \rho_{A,\ell}(Gal(\ol{K}/K))$ generate~$\Z_\ell[\rho_{A,\ell}(Gal(\ol{K}/K_\ell)]$ as a~$\Z_\ell$-module. Note that~$g_1^\ell\equiv$ [....]
%
%If~$\ell=2$ one adapts the argument with~$K_\ell=K(A[4])$.
%
%By induction on extensions of prime degree [conclude].
%
%OLD
%Let us prove that for a given~$\ell$, we there is~$F_\ell(A,K,[E:K])$ such that
%\[
%\left[Z\tens\Z_\ell:\Z_\ell[\rho_A(Gal(\ol{K}/E))]\right]<F_\ell(A,K,[E:K]).
%\]
%For~$K=E$ a number field, this is~\cite{Deligne}. 
%
%Let~$G$ be a group and~$\rho:G\to GL(n,\Z_\ell)$ be a homomorphism.
%Define~$G_k:=\ker(G\to GL(n,\Z_\ell)\to GL(n,\Z/(\ell^k))$. We claim that, for~$k\geq 1$,
%\[
%[G_k:G_{k+1}]\neq 1\Rightarrow \ell\mid [G_k:G_{k+1}].
%\]
%Indeed,~$G_k/G_{k+1}$ is isomorphic to a subgroup of
%\[
%1+\ell^k\End(\Z_\ell^n)/\left(1+\ell^k\End(\Z_\ell^n)\right)\simeq (\End(\F_\ell),+).
%\]
%We deduce that~$[E(A[\ell]):K(A[\ell])$...
%Moreover...
%
%For general~$K=E$, let us prove the claim
%\[
%\forall G\to GL(m,\Z_\ell), G\cap 1+\ell^k\End(\Z_\ell^n)
%\]
%... .
%It is enough to prove~$\left[\Z_\ell[\rho_A(Gal(\ol{K}/E))]\right]:\Z_\ell\left[\rho_A(Gal(\ol{EK}/E))]\right]]$. We recall that
%\[
%\frac{1+\ell Z\tens\Z_\ell}{1+\ell^k Z\tens\Z_\ell}
%\]
%as a subquotient of~$Z\tens\Z_ell^\times$, is a~$\ell$-group of cardinal at least~$\ell^{k-1}$. For~$\ell\nmid [E:K]$...~$F_\ell=\ell$.
%
%[TODO reste] 

\subsection{Uniform variant of a result of Serre}
Here is a variant of Theorem~\ref{thm:Serre eAK}, which is uniform on a Hybrid orbit of abelian varieties.
\begin{theorem}[{\cite[loc. cit.]{Serre:O4}} and {\cite[App.~B]{RY:IHES}}]\label{thm:Serre hybrid}
Let~$A$ be an abelian variety over a finitely generated field extension~$K/\Q$ and define~$e(A,K)$ be the smallest~$e\in\Z_{\geq1}$ such that~\eqref{eq:Serre e}.
% be as in Theorem~\ref{Serre}.
 Let~$B\leq A$ be an algebraic subgroup defined over~$\ol{K}$,  let~$K\leq K(B)\leq \ol{K}$ the smallest extension over which~$B$ is defined.
Then
\[\label{eq:isog unif}
%e(A,K)=e(A,K(C))\text{ and }
e(B,K(B))\text{ and }e(A/B,K(B))\text{ divides }e(A,K).
\]

Let~$d\in\Z_{\geq1}$. There exists~$e(A,K,d)\in\Z_{\geq1}$ such that for every abelian variety~$C$ of dimension at most~$d$, with~$CM$ with reflex field~$E$, and defined over the Hilbert class field~$E(C)/E$, and for every algebraic subgroup~$F\leq A\times C$, 
\[
e((A\times C)/F,K\cdot E^\ast(C))\leq e(A,K,d),
\]
where~$Gal(\ol{\Q}/ E^\ast(C))$ is the inverse image of~$GL(1,\wh{\Z})$ by~$Gal(\ol{\Q}/ E(C))\to Aut((A\times C)_{tors})$.
%\[
%\forall \lambda\in GL(1,\wh{\Z})\leq GL(\wh{T}(A)), \lambda^f\in \rho_A(Gal(\ol{K}/K\cdot E(C))).
%\]
\end{theorem}
\begin{proof}[Proof of Theorem~\ref{thm:Serre hybrid} assuming Theorem~\ref{thm:Serre unif principal}]
Let~$M$ be the Mumford-Tate group of~$A\times C$. We recall that there is an inclusion~$GL(1)\to Z(M)$ such that the~$\A_f$-linear representation of~$M(\A_f)$ on the~$\A_f$-Tate module of~$A\times C$ induces the action of~$GL(1,\A_f)$ by homotheties.

Let~$Z(M)(\wh{\Z})$ be as in Theorem~\ref{thm:Serre unif principal}. Let~$e:=e'(A,K,\dim(C))$ be as in Theorem~\ref{thm:Serre unif principal}. Then Theorem~\ref{thm:Serre unif principal} implies in particular that for every~$\lambda \in GL(1,\wh{\Z})$, the image of~$Gal(\ol{K}/K\cdot E(C))\to M(\A_f)$ contains~$\lambda^e$. Equivalently, the image of~$Gal(\ol{K}/K\cdot E(C))\to Aut((A\times C)_{tors})$ contains~$\lambda^e$.

By definition,~$\lambda^e$ is in the image of~$Gal(\ol{K}/K\cdot E^\ast(C))\leq Gal(\ol{K}/K\cdot E(C))\to Aut((A\times C)_{tors})$. 

Note that~$F$ is globally invariant under homotheties. Therefore~$(A\times C)/F$ and the quotient map~$A\times C\to (A\times C)/F$ are defined over~$K\cdot E^\ast(C)$. Let~$\sigma\in Gal(\ol{K}/K\cdot E^{\ast}(C))$ be such that~$\sigma$ act by the homothety~$\lambda^e$ on~$(A\times C)_{tors}$. Then~$\sigma$ act by the homothety~$\lambda^e$ on~$(A\times (C/F))_{tors}\simeq (A\times C)_{tors}/F_{tors}$.

This implies Theorem~\ref{thm:Serre hybrid}.
\end{proof}
%\begin{proof}Let~$L_{A,K}=\rho_A(Gal(\ol{K}/K)\cap GL(1,\wh{\Z})\leq GL(\wh{T}(A))$. Let~$\lambda\in GL(1,\wh{\Z})$ and~$\sigma\in Gal(\ol{K}/K)$ be suc that~$\rho_A(\sigma)=\lambda$.
%Then~$\lambda$, acting on~$A(\ol{K})_{tors}$, stabilises the group~$B(\ol{K})_\tors$. Therefore~$B=\ol{B(\ol{K})_\tors}^{Zar}=\sigma(\ol{B(\ol{K})_\tors}^{Zar})=\sigma(B)$.
%That is,~$\sigma\in Gal(\ol{K}/K(B))$.  
%We deduce
%\[
%L_{A,K}=L_{A,K(B)}:=\rho_A(Gal(\ol{K}/K(B))\cap GL(1,\wh{\Z})\leq GL(\wh{T}(A)).
%\] 
%
%The morphism~$\wh{T}(A)\to \wh{T}(A/B)$ is~$Gal(\ol{K}/K(B))$-equivariant. Therefore~$\sigma$ acts by~$\Lambda\in GL(1,\wh{\Z})\leq GL(\wh{T}(A/B))$ on~$\wh{T}(A/B)$.
%We deduce
%\[
%L_{A,K}\leq L_{A,K(B)}\hookrightarrow L_{A/B,K(B)}.
%\]
% 
%It follows~\eqref{eq:isog unif}.
%
%Let~$d$ and~$C$ be as in the statement. We use~\cite{}...
%\end{proof}

%\begin{corollary}Let~$E(A_{tors}\times C_{tors})$ be the field generated by~$K$ and~$E(C)$ and the coordinates of isogeny class NOf points in~$C(\ol{K})_{tors}\times A(\ol{K})_{tors}$. 
%Let~$A\times C\to B$ be an epimorpism. Then
%\[e(B,K(A_{tors}\times C_{tors}))\leq e(A,K,d).\]
%\end{corollary}

Theorem~\ref{thm:Serre hybrid} is therefore a consequence of the following.
\begin{theorem}\label{thm:Serre unif principal}
 Assume that~$A$ is a product of simple factors, and that the CM factors of~$A$ are principal, and let~$C$ be a principal CM abelian variety. We denote by~$E$ the reflex field of~$C$ and by~$E(C)/E$ the Hilbert class field extension. We recall that~$C$ admits a model over~$E(C)$.


Let~$M$ be the Mumford-Tate group of~$A\times C$ and let~$Z(M)(\widehat{\Z})\leq Z(M)({\A_f})$ be the maximal compact subgroup. Let
\[
\rho_{A\times C}(Gal(\ol{K}/K\cdot E(C))\to M(\A_f)
\]
be the Galois representation induced by the Tate module of~$A\times C$, and let~$U\leq  M(\A_f)$ be its image.

Then there exists~$e':=e'(A,K,\dim(C))$ such that
\begin{equation}\label{eq:thm SUP}
\{z^{e'}| z\in Z(M)(\widehat{\Z})\}\subseteq U.
\end{equation}
\end{theorem}
The following was proved in~\cite[App. B]{RY:IHES}. Let~$M_A$ be the Mumford-Tate group of~$A$, and let~$U_{A/K}$ be the image of~$\rho_A:Gal(\ol{K}/K)\to M_A(\A_f)$. 
For every prime~$p$, let~$Y(p)$ be the image of~$U_{A/K}$ in~$M^{ad}(\Q_p)$, and let~$Y(p)^{ab}$ be its maximal abelian quotient. Then, by~\cite[App. B (110) and (111)]{RY:IHES}
\begin{itemize}
\item for every prime~$p$, we have
\begin{equation}\label{eq:110}
c(p,A,K):=\# Y(p)^{ab}<+\infty
\end{equation}
\item and there exists~$p_0=p_0(A,K)$ such that
\begin{equation}
\label{eq:111}
\forall p\geq p_0, \# Y(p)^{ab}=1.
\end{equation}
\end{itemize}
Moreover~\cite[App. B (113)]{RY:IHES} imply that there exists~$p_1(A,K)$ such that for every~$p\geq p_1(A,K)$, the group~$Y(p)$ is generated by topologically~$p$-nilpotent elements. In particular, for~$p\geq p_1(A,K)$ and every closed subgroup~$H\leq Y(p)$, we have
\begin{equation}\label{eq:113}
[Y(p):H]<p\Rightarrow
[Y(p):H]=1.
\end{equation}

\begin{proof}[Proof of Theorem~\ref{thm:Serre unif principal}]
Let~$Y_{E(C)}(p)\leq Y_E(p)\leq Y_{A/K}(p)$ be the images of~$Gal(\ol{K}/K\cdot E(C))\leq Gal(\ol{K}/K\cdot E)\leq Gal(\ol{K}/K\cdot E)$ in~$M^{ad}(\A_f)=M_A^{ad}(\A_f)$. Then
\[
[Y(p):Y(p)_E]\leq [K\cdot E:K]\leq [E:C]\leq f(\dim(C)):=2^{\dim(C)}\cdot \dim(C)!.
\]
and~$Y_{E(C)}(p)\leq Y_E(p)$ is a normal subgroup such that~$Y_E(p)/Y_{E(C)}(p)$ is abelian. 
%We denote by~$Y_E(p),Y_{E(C)}(p)\leq Y(p)$ the images of~$U_E$ and~$U_{E(C)}$.

Note that, for every~$p$, the~$p$-adic Lie group~$Y(p)$ has only finitely many open subgroups of index at most~$f:=f(\dim(C))$. The proof of~\cite[App. B (111)]{RY:IHES} applies to every open subgroup of~$Y(p)$.
We have thus
\[
\#{Y_E(p)}^{ab}\leq c_0(A,K,f,p):=
\max\{\# H^{ab}| [Y(p):H]\leq f\}<+\infty.
\]
Therefore
\[
[Y_E(p):Y_{E(C)}]\leq c_0:=c_0(A,K,f,p).
\]
Similarly,
\[
\#{Y_{E(C)}(p)}^{ab}\leq c_1(A,K,f,p):=c_0(A,K,f\cdot c_0,p).
\]

Let~$p\geq p_2(A,K,\dim(C)):=\max\{p_0(A,K);p_1(A,K);f(\dim(C))+1\}$. By~\eqref{eq:113}, we deduce from~$[Y(p):Y_E(p)]\leq f(\dim(C))$ that
\[
Y(p)=Y_E(p).
\]
We can thus deduce from~\eqref{eq:111} that
\[
Y_E(p)=Y_{E(C)}(p),
\]
and thus that
\[
\#Y(p)^{ab}=\#Y_E(p)^{ab}=\#Y_{E(C)}(p)^{ab}=1.
\]
Finally, we have
\begin{multline}
\prod_p \#{Y_{E(C)}(p)}^{ab}\\
\leq
c_2(A,K,\dim(C))
:=
\prod_{p\leq p_2(A,K,\dim(C))}  c_1(A,K,f,p)<+\infty.
\end{multline}
Let~$U_{E(C)}$ be the image of~$\rho_{A\times C}(Gal(\ol{K}/K\cdot E(C))\to M(\A_f)$.
We consider and follow the proof of \cite[App. B, \S B.4.2, p.\,326]{RY:IHES}, for~$M$ the Mumford-Tate group of~$A\times C$ and~$U_{E(C)}$ instead of~$U$. 

We deduce that 
\begin{multline}
[\Gamma:(\Gamma\cap G_1)\times (\Gamma\cap G_2)]=\abs{\Gamma/(\Gamma\cap G_1)}
\\=
\abs{\Gamma/(\Gamma\cap G_2)}
\leq c_2(A,K,\dim(C)).
\end{multline}
Then~\cite[App. B, \S B.4.2, p.\,326]{RY:IHES} gives us
\[
\left[U_{E(C)}:\left(U_{E(C)}\cap M^{der}(\widehat{\Z})\right)\cdot
\left(U_{E(C)}\cap Z(M)(\widehat{\Z})\right)\right
]\leq c_2(A,K,\dim(C)).
\]
We consider the map
\[
\pi:Z(M)\to M^{ab}.
\]
Then~$\ker(\pi)=Z(M^{der})$, where~$M^{der}$ only depends on~$A$. Then for every prime~$p$, we have
\[
\#Z(M^{der})(\Q_p)\leq \#Z(M^{der})(\ol{\Q_p})=c_3(A):=\# Z(M^{der})(\ol{\Q}).
\]
We deduce that
\[
\forall z\in Z(M)(\wh{\Z}), \pi(z)=1\Rightarrow z^{c_3(A)}=1.
\]

We recall that~$C$ is assumed to be a principal CM abelian variety. By Appendix~\ref{sec:appendix reciprocity}, Theorem~\ref{thm:Reciprocity principal} (cf. also \cite[Lem.\,2.7 and its proof]{UY}), there exists~$n:=n(A,\dim(C),K)$ such that for every~$\lambda$ in~$M^{ab}(\wh{\Z})$, the element~$\lambda^{n}$
is in the image of~$U\to M^{ab}(\wh{\Z})$.

Let~$U':=\left(U_{E(C)}\cap M^{der}(\widehat{\Z})\right)\cdot \left(U_{E(C)}\cap Z(M)(\widehat{\Z})\right)$.
Then the image of~$U'\to M^{ab}(\wh{\Z})$ contains
\[
\{\lambda^{n\cdot c_2(A,K,\dim(C))!}|\lambda\in M^{ab}(\wh{\Z})\}.
\]
But this is also the image of~$U''':=U\cap \left(U_{E(C)}\cap Z(M)(\widehat{\Z})\right)$. 
By Lemma~\ref{lem:liftkernel}, we deduce that
\[
\forall z\in Z(M)(\widehat{\Z}), \lambda^{n\cdot c_2(A,K,\dim(C))!\cdot c_3(A)}\in U'''.
\]
This proves~\eqref{eq:thm SUP} for~$e'=n\cdot c_2(A,K,\dim(C))!\cdot c_3(A)$.
\end{proof}


%OLD
%
%
%The claim follows from the following Lemma applied to~$f:=\pi:Z(M)(\wh{\Z})\to M^{ab}(\wh{\Z})$ and~$U:=U_{E(C)}\cap Z(M)(\wh{\Z})$ and~$a:=c_3(A)$ and~$b:=e''$.	
%\begin{lemma}\label{lem:fGH}
%Let~$f:G\to H$ be a morphism of abelian groups. Let~$U\leq G$ be a subgroup. Let~$a,b\in \Z_{\geq1}$ be such that
%\[
%\forall h\in H, h^a\in f(U)\text{ and }\forall g\in \ker(f), g^b=0. 
%\]
%Then,
%\[
%\forall g\in G, g^{a\cdot b}\in U.
%\]
%\end{lemma}
%\begin{proof}[Proof of Lemma~\ref{lem:fGH}] The sequence
%\[
%\ker(f)\to G/U\to H/f(U)
%\]
%is exact at~$G/U$. Thus~$G/U$ is an extension of~$f(U)$ by~$\ker(f)/(\ker(f)+U)$, which are of exponent dividing~$b$ and~$a$ respectively. Therefore~$G/U$ is of exponent dividing~$a\cdot b$. The Lemma~\ref{lem:fGH} follows.
%
%\end{proof}
%\begin{proof}
%\end{proof}
%We observe that~$c_4=c_4(K):=[K\cap\ol{\Q}:\Q]<+\infty$.
%We deduce that for every~$\lambda$ in~$M^{ab}(\wh{\Z})$, the element~$\lambda^{e''\cdot c_4}$
%is in the image of~$Gal(\ol{\Q}/((K\cap{\ol{\Q}})\cdot E(C)))\to M^{ab}(\wh{\Z})$.
%
%We deduce that for every~$\lambda$ in~$M^{ab}(\wh{\Z})$, the element~$\lambda^{e''\cdot c_4\cdot c_2}$
%is in the image of~$U_{E(C)}\cap Z(M)(\wh{\Z})\to M^{ab}(\wh{\Z})$.
%
%We apply ... and deduce that
%\[
%\forall z\in Z(M)(\wh{\Z}), z^{e''\cdot c_4\cdot c_2\cdot c_3}\in U_{E(C)}\cap Z(M)(\wh{\Z}).
%\]
%This proves~\eqref{eq:thm SUP} with
%\[
%e'(A,K,\dim(C)):=e''(A,\dim(C))\cdot c_4(K)\cdot c_2(A,K,\dim(C))\cdot c_3(A).\qedhere
%\]
%\end{proof}


\subsection{}

We will later use the following.
\begin{lemma}\label{lem:Lang ne}
For every~$e\in \Z_{\geq1}$, there exists~$n_e\in \Z_{\geq1}$ such that we have the following inclusion of ideals
\[
\sum_{\lambda\in \Z^{\times}} (\lambda^e-1)\cdot \wh{\Z}\geq n_e\cdot \wh{\Z}.
\]
\end{lemma}
\begin{proof}Let~$\ell$ be a prime such that~$\ell-1$ does not divide~$e$. As~$\F_\ell^\times\approx \Z/(\ell-1)\Z$, there exists~$\ol{\lambda}\in \F^\times_\ell$ such that~$\ol{\lambda}^e\neq1$. Let~$\lambda\in \wh{\Z}^\times$ 
such that~$\ol{\lambda}=\lambda\pmod{e}$. The image of~$(\lambda^e-1)\cdot \wh{\Z}$ in~${\F}_\ell$ contains the non-zero invertible element~$\ol{\lambda}^e-1$. This image is thus~${\F}_\ell$.
By Nakayama lemma,~$({\lambda}^e-1)\cdot\Z_\ell=\Z_\ell$.

Let~$\ell$ be any prime and~$\lambda=1+\ell$. Then~$((1+\ell)^e-1)\cdot \Z_\ell$ is a non zero ideal of~$\Z_\ell$, and thus is open in~$\Z_\ell$.

We deduce that, with~$n_e:=\prod_{\ell-1\mid e}((1+\ell)^e-1)$, we have 
\[
\sum_{\lambda\in \Z^{\times}} ((\lambda^e-1)-1)\cdot \wh{\Z}\geq \left(\prod_{\ell-1\mid e}((1+\ell)^e-1)\Z_\ell \right)\times\left(\prod_{\ell-1\nmid e} \Z_\ell\right)\geq n_e\cdot \wh{\Z}.
\] This proves Lemma~\ref{lem:Lang ne}.
\end{proof}

\subsection{Kummer theory for abelian varieties}\label{sec:Kummer}

Let~$\Gamma\leq A(K)$ be a subgroup, and let~$\Gamma_0:=\{Q\in A(\ol{K})|\Z_{\geq1}\cdot Q\cap \Gamma\neq \emptyset\}$. We have an exact sequence
\[
0\to A_{tors}\to \Gamma_0\to \Gamma\tens \Q\to 0.
\]
Observe that~$Gal(\ol{K}/K)$ acts trivially on~$\Gamma$, and thus acts trivially on~$\Gamma\tens \Q$ and~$\Gamma\tens\Q/\Z$.
Therefore, we have~$Gal(\ol{K}/K)$-equivariant sequence
\[
0\to A_{tors}\to \Gamma_0/\Gamma \to \Gamma\tens \Q/\Z\to 0.
\]
We assume~$\Gamma\approx\Z^k$ for some~$k\in\Z_{\geq1}$. Then~$\Hom(\Q/\Z,\Gamma\tens\Q/\Z)\simeq \Gamma\tens\wh{\Z}\approx\wh{\Z}^k$ and
\[
\wh{T}_{A,\Gamma}:=
\Hom(\Q/\Z, \Gamma_0/\Gamma)
\approx 
\wh{\Z}^{2\dim(A)+k}
\]
is a~$Gal(\ol{K}/K)$-equivariant extension
\begin{equation}\label{eq:extension}
0\to \wh{T}_{A}\to \wh{T}_{A,\Gamma}\to \Gamma\tens\wh{\Z}\to 0.
\end{equation}
The representation~$\rho_{A,\Gamma}:Gal(\ol{K}/K)\to Aut(\wh{T}_{A,\Gamma})$ stabilises~$\wh{T}_{A}$ and acts trivially on the quotient~$\wh{T}_{A,\Gamma}/\wh{T}_{A}$. There exists thus a map
\[
m=m_{A,K,\Gamma}:Gal(\ol{K}/K)\to\Hom(\Gamma,\wh{T}(A))
\] such that we can represent
\[
\rho_{A,\Gamma}(\sigma)
=
\begin{pmatrix}
\rho_A(\sigma)&m(\sigma)\\
0& 1
\end{pmatrix}
\in
\begin{pmatrix}
\End(\wh{T}(A))&\Hom(\Gamma,\wh{T}(A))\\
\Hom(\wh{T}(A),\Gamma\tens\wh{\Z})& \End(\Gamma\tens\wh{\Z})
\end{pmatrix}.
\]
We call~$(\wh{T}_{A,\Gamma},\rho_{A,\Gamma})$ the \emph{affine Tate module}. The map~$m$ is the cocycle of the~$Gal(\ol{K}/K)$-module~$\Hom(\Gamma,\wh{T}(A))$ coming
from the extension~\eqref{eq:extension}. Its restriction on~$Gal(\ol{K}/K(A_{tors}))=\ker(\rho_A)$, which acts trivially on~$\Hom(\Gamma,\wh{T}(A))$, is an homomorphism of groups~$\ker(\rho_A)\to (\Hom(\Gamma,\wh{T}(A)),+)$. 


The object of \emph{Kummer theory for abelian varieties} is to study
\[
\rho_{A,\Gamma}(\ker(\rho_A))=
\begin{pmatrix}
1&m(\ker(\rho_A))\\
0& 1
\end{pmatrix}
=
Im(\rho_{A,\Gamma})
\cap
\begin{pmatrix}
1&\Hom(\Gamma,\wh{T}(A))\\
0& 1
\end{pmatrix},
\]
or to study the map
\begin{equation}\label{eq:Kummer m}
m:Gal(\ol{K}/K(A_{tors}))\to \Hom(\Gamma,\wh{T}(A)).
\end{equation}
We refer to \cite{Ribet}, \cite[App. 2]{Hindry} or \cite[{\S\S} 1--2]{Bertrand} for the following.
\begin{theorem}[Ribet]\label{thm:Kummer}
Let~$A'$ be the abelian variety~$\Hom(\Gamma,A)\approx A^k$, and let~$P\in A'(K)$ be the identity~$\Gamma\to A$.
Let~$[P]:=\ol{\Z\cdot P}^{Zar}$ and~$B=[P]^0$. Then
\[
m(Gal(\ol{K}/K(A_{tors})\cdot K(B))\leq \wh{T}(B)\text{ in }\wh{T}(A')\simeq \Hom(\Gamma,\wh{T}(A))
\]
Let
\begin{equation}\label{}
k(A,K,\Gamma):=[\wh{T}(B):m(Gal(\ol{K}/K(A_{tors})\cdot K(B))]<+\infty.
\end{equation}
If that~$K$ is of finite type over~$\Q$,
then
\[
k(A,K,\Gamma)<+\infty.
\]
\end{theorem}

%Here are some immediate functoriality properties.
%\begin{corollary}[SUPPRIMER???]\label{thm:isogeny Kummer}
%For every finite extension~$E/K$ we have
%\[
%k(A,K,\Gamma)\text{ divides }k(A,E,\Gamma)\text{ divides }[E:K]\cdot k(A,E,\Gamma).
%\]
%For every $K$-isogeny~$A\to A'$, we have ???
%\[
%k(A,K,\Gamma)=k(A',K,\Gamma).
%\]
%For every~$K$-morphism~$A\to A/B$, we have:
%\[
%k(A/B,K,\Gamma)\text{ divides }k(A,K,\Gamma).
%\]
%\end{corollary}
%\begin{proof}
%The statement~\eqref{} follows from the definition and the observations
%\[
%\rho_{A,\Gamma}(Gal(\ol{K}/E))\leq 
%\rho_{A,\Gamma}(Gal(\ol{K}/K))
%\]
%and
%\[
%Gal(E:K)\text{ acts transitively on }
%\frac{\rho_{A,\Gamma}(Gal(\ol{K}/E))}
%{\rho_{A,\Gamma}(Gal(\ol{K}/K))}.
%\]
%
%We have
%\[
%k(A,K,\Gamma)=k(A\tens \Gamma^\vee,K,P\cdot\Z)=k([P],K(A_{tors}),P\cdot\Z).
%\]
%It amounts to prove, with~$K'=K(A_{tors})$
%\[
%k([P]+B/B,K',(P+B)\cdot\Z)\text{ divides }
%k([P],K',P\cdot\Z).
%\]
%
%Assume now that...~$B$ is finite. We may assume that~$\phi:A\to A'$ 
%is a non-factorable isogeny of prime exponent~$\ell$. We have
%\[
%k(\phi([P]),K',\phi(P)\cdot\Z)\text{ divides }
%\ell\cdot k([P],K',P\cdot\Z).
%\]
%Using the complementary isogeny,
%\[
%k([P],K',P\cdot\Z).
%\text{ divides }
%\ell\cdot k(\phi([P]),K',\phi(P)\cdot\Z).
%\]
%We conclude
%\[
%k([P],K',P\cdot\Z).
%=
%\ell\cdot k(\phi([P]),K',\phi(P)\cdot\Z).
%\]
%
%Assume now that~$B\leq [P]$ is an abelian subvariety (connected): 
%Then~$\wh{T}([P])$ surjects onto~$\wh{T}([P]/B)$: TODO write conclusion.
%\end{proof}

\subsection{Uniform Kummer theory}

We need the following uniform version.
\begin{theorem}[Uniform Kummer theory on a hybrid orbit]\label{thm:uniform Kummer1}
For every~$d\in\Z_{\geq0}$, there exists~$k(A,K,\Gamma;d)\in\Z_{\geq1}$ such that for every CM abelian variety~$C$ of dimension at most~$d$ %(of reflex field~$E_C$ of degree at most~$d$),
\begin{equation}\label{eq:unifKummer1}
k(A\times C,K(C),\Gamma)\leq k(A,K,\Gamma;d),
\end{equation}
where~$\Q(C)$ is the field of definition of~$C$, and~$K(C)=K\cdot \Q(C)$.
\end{theorem}
In essence, this is a statement about linear independance, over~$K(A_{tors})$, between~$K(C_{tors})$ and~$K(\Gamma_0)$.

Theorem~\ref{thm:uniform Kummer1} will is a consequence of the following.
\begin{theorem}[Uniform Kummer theory]\label{thm:uniform Kummer2}
For every~$d\in\Z_{\geq0}$, there exists~$\kappa(A,K,\Gamma;d)$ such that for every extension~$[E:K]$ of dimension at most~$d$
\begin{equation}\label{eq:unifKummer2}
k(A,E^{ab},\Gamma)\leq  \kappa(A,K,\Gamma;d)<+\infty.
\end{equation}
\end{theorem}
We will prove Theorem~\ref{thm:uniform Kummer2} in~\S\S\ref{sec:proof 4.10}. We first deduce Theorem~\ref{thm:uniform Kummer1}.

\begin{proof}[Proof of Theorem~\ref{thm:uniform Kummer1} assuming Theorem~\ref{thm:uniform Kummer2}]
Let~$E'$ be the reflex field of the CM-type of~$C$ and let~$E=E'\cdot K$. One has~$[E:K]\leq d:=2^{\dim(C)}\cdot \dim(C)!$, and~$K(C,C_{tors})\subseteq E^{ab}$.
Therefore~$\ker(\rho_{A\times C})=K(C,(A\times C)_{tors})=K(C,C_{\tors},A_{tors})
\subseteq E^{ab}(A_{tors})$. We deduce
\[
m_{A,\Gamma}(Gal(\ol{K}/E^{ab}(A_{tors}))
\subseteq 
m_{A\times C,\Gamma}(\ker(\rho_{A\times C}))
\subseteq 
\wh{T}(B)
\] Therefore~\eqref{eq:unifKummer2} implies~\eqref{eq:unifKummer1} with
\[k(A\times C,K(C_{tors}),\Gamma):= \kappa(A,K,\Gamma;2^{\dim(C)}{\dim(C)}!).\qedhere\]
\end{proof}


\subsubsection{Preliminaries} We adapt~\cite[App. 2]{Hindry}.

\begin{proposition}\label{prop:411}
Let~$A$ be an abelian variety over an extension~$K/\Q$ and let~$N\in\Z_{\geq1}$. There is an injective homomorphism
\[
A(K)/N\cdot A(K)\to H^1(Gal(\ol{K}/K); A[N])
\]
such that for every~$P\in A(K)$ and~$Q\in A(\ol{K})$ such that~$N\cdot Q=P$, the image of~$P+NA(K)$ is the cohomology class of the cocycle~$\sigma\mapsto \sigma(Q)-Q:Gal(\ol{K}/K)\to A[N]$. 

Let~$E/K$ be a finite extension. Then the kernel of~$A(K)/N\cdot A(K)\to A(E)/N\cdot A(E)$ is annihilated by~$[E:K]$.
\end{proposition}
\begin{proof}See~~\cite[App. 2, Lem.~F]{Hindry} for the first statement. Let us prove the last statement. We can embed~$A(K)/N\cdot A(K)$ in~$H^1(Gal(\ol{K}/K); A[N])$ and~$A(E)/N\cdot A(E)$ in~$H^1(Gal(\ol{K}/E); A[N])$. The map~$A(K)/A(K)\to A(E)/N\cdot A(E)$ is induced by the restriction map~$H^1(Gal(\ol{K}/K); A[N])\to H^1(Gal(\ol{K}/E); A[N])$. It suffices to show that the kernel of the latter is annihilated by~$[E:K]$. This is~\cite[I\S\, 2.4, Prop. 9]{LNM5}.
\end{proof}
\begin{proposition}\label{prop:412}
 We consider the context of Theorem~\ref{thm:Kummer}. Let~$E/K$ be a finite extension and let~$n:=n_{e(A,K)\cdot [E:K]!}$ in the notations of Lemma~\ref{lem:Lang ne}. Let~$G'$ and~$H$ be the image and kernel of~$G:=Gal(\ol{K}/E)\to Aut(\wh{T}(A))\times G^{ab}$. 
\begin{enumerate}
\item \label{4121}
For every~$N\in\Z_{\geq1}$, we have a short exact sequence
\begin{equation}\label{eq:truccc}
1\to H^1(G'; B[N]) \to H^1(G; B[N])\to H^1(H;B[N])
\end{equation}
and a natural identification~$H^1(H;B[N])\simeq Hom(H,B[N])$.
\item \label{4122}
For every~$N\in\Z_{\geq1}$, the group~$H^1(G'; B[N])$ is annihilated by~$n$.
\item \label{4123}
For every~$Q\in B(E)$ and~$N\in\Z_{\geq1}$, we have
\[
n\cdot Q\not\in N\cdot B(E)\Rightarrow m_Q(H)\not\subseteq N\cdot \wh{T}(B).
\]
where~$m_Q\in Hom(H,B[N])$ is such that the image of~$Q$ by
\[
B(E)\to B(E)/NB(E)\hookrightarrow H^1(G;B[N])\to H^1(H;B[N])\simeq Hom(H,B[N]).
\]
is the map
\[
H\xrightarrow{m_Q} \wh{T}(B)\to \wh{T}(B)\tens \Z/(N)\simeq B[N].
\]
\end{enumerate}
\end{proposition}

\begin{proof}[Proof of Proposition~\ref{prop:412}]Let us prove~\eqref{4121}. Note that~$H$ acts trivially~$B[N]$. This implies $B[N]^H=B[N]$, and allows us to deduce the exact sequence~\eqref{eq:truccc} from~\cite[Ch.\,I, \S2.6 b)]{LNM5}. This trivial action also implies the identification.

Let us prove~\eqref{4122}.
By Lemma~\ref{lem:Lang ne}, it is enough to prove the claim that for every~$\lambda \in\wh{\Z}^\times$, the group~$H^1(H;A[N])$ is annihilated by~$\lambda^{e'}-1$ with~$e':=e(A,K)\cdot [E:K]!$.
\begin{proof}[Proof of the claim]By Theorem~\ref{thm:Serre eAK}, there exists~$\sigma\in Gal(\ol{K}/K)$ such that~$\sigma$ acts by~$\lambda^{e(A,K)}$ on~$\wh{T}(A)$. Therefore~$\sigma':=\sigma^{[E:K]!}\in Gal(\ol{K}/E)$ and acts by~$\lambda^{e'}$.  Therefore~$\rho_A(\sigma')$ is central in~$Aut(\wh{T}(A))\geq \rho_A(Gal(\ol{K}/K))$. But the image of~$\sigma$ in~$Gal(\ol{K}/E)^{ab}$ is obviously central. Therefore the image of~$\sigma'$ in~$Aut(\wh{T}(A))\times G^{ab}\geq G'$ is central. Note that~$\sigma'$ acts by~$\lambda^{e'}$ on~$B[N]$, and thus by~$\lambda^{e'}$ on~$H^1(H;B[N])$. By Sah's~\cite[App. 2, Lem.~D]{Hindry}, the action of~$\sigma'$ on~$H^1(H;B[N])$ is trivial. This proves the claim.
\end{proof}

Let us prove \eqref{4123}. We can embed~$B(E)/N\cdot B(E)$ in~$H^1(G;B[N])$. By Lemma~\ref{lem:liftkernel} for~$\phi:H^1(G;B[N])\to H^1(H;B[N])$, and~$x=Q+N\cdot B(E)\in B(E)/N\cdot B(E)\leq H^1(G;B[N])$, we have
\[
n\cdot Q+N\cdot B(E)\neq N\cdot B(E)\Rightarrow \phi(Q+N\cdot B(E))\neq 0.
\]
We observe that~$\phi(Q+N\cdot B(E))\in Hom(H,B[N])$ is the image of~$m_Q \in Hom(H,\wh{T}(B))$. We deduce
\[
\phi(Q+N\cdot B(E))\neq 0\Leftrightarrow m_Q(H)\not\subseteq N\cdot \wh{T}(B).
\]
The conclusion follows.

Proposition~\ref{prop:412} is proven.
\end{proof}

\begin{corollary}\label{cor:613}
Let~$\Theta\leq A(K)$ be a subgroup. There exists~$c':=c'(\Theta,A(K))$ such that
\[
\forall N\in\Z_{\geq0},\ker\bigl(\Theta/N\Theta\to A(E)/NA(E)\bigr)
\]
is annihilated by~$c'\cdot [E:K]$. 


For all~$N\in\Z_{\geq0}$, the kernel of~$\Theta/N\Theta\to Hom(H,B[N])$ is annihilated by~$c'\cdot [E:K]\cdot n$.
\end{corollary}
\begin{proof} By Proposition~\ref{prop:411}, the kernel~$K_1$ of~$A(K)/NA(K)\to A(E)/NA(E)$ is annihilated by~$[E:K]$. By Proposition~\ref{prop:412} the kernel~$K_2$ of~$A(E)/NA(E)\to  Hom(H,B[N])$ is annihilated by~$n$. The kernel of~$A(K)/NA(K)\to  Hom(H,B[N])$ is an extension~$K_3$ of~$K_2$ by~$K_1$. It follows that~$K_3$ is annihilated by~$n\cdot [E:K]$.

Let~$c'$ be the exponent of~$(A(K)/\Theta)_{tors}$. Theorem~\ref{thm:MWN} implies that the group~$A(K)/\Theta$ is finitely generated, and thus~$c'<+\infty$. Therefore~$x\in A(K), \exists N, N\cdot x\in \Theta\Rightarrow c\cdot x\in \Theta$. Thus~$x\in N\cdot A(K)\Rightarrow c\cdot x\in N\cdot \Theta$. Equivalently, the kernel~$K_4$ of~$\Theta/N\Theta\to A(K)/NA(K)$ is annihilated by~$c'$. The kernel of~$\Theta/N\Theta\to Hom(H,B[N])$ is an extension of~$K_3$ by~$K_4$, and is thus annihilated by~$c'\cdot [E:K]\cdot n$.
\end{proof}
\begin{proposition}\label{prop:preliminaire}
The map
\begin{equation}\label{eq:100}
\alpha\mapsto\alpha(P):\End(B)\mapsto \Theta:=\End(B)\cdot P
\end{equation}
is bijective.

There exists~$c=c(P,B,K)$ such that
\begin{equation}\label{eq:0}
\forall N\in\Z_{\geq1}, \alpha\not\in N\cdot \End(B) \Rightarrow m_{\alpha(P)}(H)\not\subseteq c\cdot [E:K]\cdot N\cdot \wh{T}(B).
\end{equation}
\end{proposition}
\begin{proof}Recall that~$B$ is the algebraic group generated by~$P$. Therefore, for~$\alpha\in \End(B)$ we have~$\alpha(P)=0\Rightarrow \alpha=0$. This proves the injectivity of~\eqref{eq:100}. The surjectivity of~\eqref{eq:100} is immediate.

Let~$c=c(P,B,K):=c'(\Theta,B(K))$ and~$n':=c\cdot [E:K]\cdot n$.
By Corollary~\ref{cor:613}, we have
\begin{equation}\label{eq:1}
n'\cdot \alpha\not\in N\cdot \End(B)\Rightarrow m_{\alpha(P)}(H)\not\subseteq N\cdot \wh{T}(B).
\end{equation}
Substituting~$N$ with~$n'\cdot N$, we obtain
\begin{equation}\label{eq:2}
n'\cdot \alpha\not\in n'\cdot N\cdot \End(B)\Rightarrow m_{\alpha(P)}(H)\not\subseteq n'\cdot N\cdot \wh{T}(B).
\end{equation}


Recall that~$\End(B)$ is a free module of finite rank over~$\Z$. Therefore~$n'\cdot \alpha\not\in n'\cdot N\cdot \End(B)\Leftrightarrow \alpha\not\in N\cdot \End(B)$. We can thus deduce~\eqref{eq:0} from~\eqref{eq:2}.
\end{proof}
\subsubsection{Proof of Theorem~\ref{thm:uniform Kummer2}}\label{sec:proof 4.10}
 We combine Poposition~\ref{prop:preliminaire} with Falting's Theorem in the form~\ref{thm:uniform Faltings} and Theorem~\ref{thm:D1} from Appendix~\ref{app:D}.

Let~$\wh{T}$ be the~$\wh{\Z}$-linear Tate module of~$B$, let~$R=\End(B)$, and let~$\wh{S'}$ be the commutant
of~$R$ in~$\End_{\wh{\Z}}(\wh{T})$. Let~$F(A,K,[E:K])$ be as in Theorem~\ref{thm:uniform Faltings}, and let~$\wh{S}:=\wh{\Z}\cdot Id_{\wh{T}}+F(A,K,[E:K])\cdot \wh{S'}$. According to Theorem~\ref{thm:uniform Faltings},
we have
\[
\wh{S}\subseteq \wh{\Z}\left[\rho_{B}(Gal(\ol{K}/E))\right].
\]
Let~$R^0:=R\tens \Q$. Choose an embedding~$\ol{K}\to \C$, and define~$T_\Z:=H^1(B(\C);\Z)$ and~$T_\Q:=T_\Z\tens\Q$.
Let~$S'$ be the commutant of~$R$ in~$\End_\Z(T_\Z)$ and~$S:=\Z\cdot Id_{T_\Z}+F(A,K,[E:K])\cdot S'$. Then~$\wh{S}\simeq S\tens \wh{\Z}$. Note that~$R$ is an order of the semisimple algebra~$R^0$ and~$S$ is an order of the semisimple algebra~$S^0:=S'\tens\Q$. 

Let~$L$ and~$C(S,R,T)$ be as in Theorem~\ref{thm:D1}. Let~$c(P,B,K)$ be as in Proposition~\ref{prop:preliminaire} and let us define~$n:=c\cdot [E:K]$. Let~$\wh{W}:=m_{P}(H) \leq \wh{T}$. Note that~$m_{\alpha(P)}=\alpha\circ m_{P}(H)$ and therefore
\[
m_{\alpha(P)}(H)\subseteq n \cdot \wh{T}(B)\Leftrightarrow
\alpha(\wh{W})\subseteq n\cdot \wh{T}.
\]

Let~$\ell$ be a prime. We claim that~$W_\ell:=\wh{W}\tens_{\wh{\Z}}\Z_\ell\leq T_\ell:=\wh{T}\tens\Z_\ell$ is open in~$T_\ell$.
\begin{proof} Assume for contradiction that~$U:=W_\ell\tens_{\Z_\ell} \Q_\ell\neq T_\ell\tens_{\Z_\ell} \Q_\ell$. Recall that~$U$ is stable under the action of~$Gal(\ol{E}/E)$. By Faltings' theorem, there exists~$\alpha\in R\tens \Z_\ell$ such that~$\ker(\alpha)=U$. We have~$\alpha\neq 0$ by assumption.
Let~$\alpha_i$ be a sequence in~$R$ converging, in~$R\tens\Z_\ell$, to~$\alpha$. We may assume~$\alpha_i\neq 0$: there exists~$N_i\in\Z_{\geq1}$ such~$\alpha_i\not\in N_i\cdot R$. As~$\alpha_i|_W$ converges to the constant map~$0:W\to T$, ~\eqref{eq:0} implies~$N_i\to 0$ in~$\Q_\ell$ as~$i\to+\infty$. Therefore~$\alpha$ belongs to~$\bigcap N_i\cdot R\tens\Z_\ell=\{0\}$. Therefore~$\alpha=0$ and~$U=\ker(\alpha)=T_\ell \tens_{\Z_\ell} \Q_\ell$, which is a contradiction.
\end{proof}
Let~$W'_\ell\leq T$ be the intersection~$T\cap W_\ell$ in~$T_\ell$. Then~$W'_\ell\leq T$ is of finite index and~$W_\ell\simeq W'_\ell\tens\Z_\ell$.

By Theorem~\ref{thm:D1}, we have
\[
W'_\ell\geq C(S,R,T)\cdot n\cdot T.
\]
Tensoring by~$\Z_\ell$, we get
\[
W_\ell\geq C(S,R,T)\cdot n\cdot T_\ell.
\]

As this is true for every prime~$\ell$, we conclude
\[
\wh{W}:={\prod}_\ell W_\ell\geq {\prod}_\ell C(S,R,T)\cdot n\cdot T_\ell=C(S,R,T)\cdot n\cdot \wh{T}.
\]
This proves Theorem~\ref{thm:uniform Kummer2} with
\begin{equation}\label{eq:Kummer explicite}
\kappa(A,\Gamma;[E:K])=c(P,B,K)\cdot [E:K]\cdot C(S,R,T).
\end{equation}


\subsection{Comparison with a Theorem of Rémond}\label{sec:CD}
The formula~\eqref{eq:Kummer explicite} is an explicit form of Theorem~\ref{thm:uniform Kummer2}. From the proofs of Proposition~\ref{prop:preliminaire} Corollary~\ref{cor:613}, we see that the quantity~$c(P,B,K)$ is the exponent $c'(\End(B)\cdot P,B(K))$ of~$(B(K)/\End(B)\cdot P)_{tors}$. We observe that the quantity
\[
(G(K)\cap (\End(G)\cdot P)_{div})/\End(G)\cdot P
\] 
from~\cite[Theorem~17 (Rémond)]{CD} is equal to~$c'(\End(B)\cdot P,B(K))$, for~$G=A$ and~$B=G_P$.

Therefore, the explicit form~\eqref{eq:Kummer explicite} of Theorem~\ref{thm:uniform Kummer2} implies~\cite[Theorem~17 (Rémond)]{CD} for~$G$ an abelian variety (\cite{CD} also considers $G=A\times GL(1)^N$ for some~$N$).

Our result is stronger:
\begin{itemize}
\item formula~\eqref{eq:Kummer explicite} exhibits an explicit dependency on the field~$E$: the dependency is linear in~$[E:K]$ (cf. Proposition~\ref{prop:411});
\item our results are uniform when~$A$ ranges through a hybrid orbit of abelian varieties (see Theorem~\ref{thm:Gal En}).
\end{itemize}



%Let~$T_\ell:=\wh{T}(B)\tens \Z_\ell\leq V:=\wh{T}(B)\tens \Q_\ell$. Let~$R:=\End(B)\tens \Z_\ell\leq R^0:=R\tens\Q_\ell\leq \End_{\Q_\ell}(V)$. Let~$\wh{S}:=\wh{\Z}[\rho_{B,P}(Gal(\ol{K}/E))]$ and~$S=\wh{S}\tens\Z_\ell\leq S^0:=S\tens\Q_\ell\leq \End_{\Q_\ell}(V)$. 
%
%
%Recall that~$R^0$ and~$S^0$ are semisimple finite~$\Q_\ell$-subalgebras of~$\End_{\Q_\ell}(V)$. There exists a finite extension~$L/\Q_\ell$ and an isomorphism~$V\simeq \bigoplus_{i=1}^k L^{n_i}\tens L^{m_i}$ such that~$R^0\tens \Q_\ell$ and~$S^0\tens\Q_\ell$ corresponds to the subalgebras~$\prod_{i=1}^k Id_{L^{n_i}}\tens\End(O_L^{m_i}) $ and~$\prod_{i=1}^k \End(O_L^{n_i})\tens Id_{L^{m_i}}$ of~$\End(\bigoplus_{i=1}^k L^{n_i}\tens L^{m_i})$. Moreover we can choose an isomorphism such that~$\bigoplus {O_L}^{n_i}$ and~$\oplus {O_L}^{m_i}$ are stable under the action of~$R$ and~$S$ on~$\oplus {L}^{n_i}$ and on~$\oplus {L}^{m_i}$.
%
%There exists~$F_R, F_S\in\Z_{\geq1}$ such that, under these identifications,
%\[
%\Z_\ell+F_R\cdot \bigoplus {O_L}^{n_i} \leq R\tens O_L\leq \bigoplus {O_L}^{n_i}
%\qquad
%\Z_\ell+F_S\cdot \bigoplus {O_L}^{m_i} \leq S\tens O_L\leq \bigoplus {O_L}^{m_i}.
%\]
%
%Recall that~$m_P(H)$ is a~$\wh{\Z}$-submodule of~$\wh{T}(B)$. Then, for~$\alpha\in R$, we have~$m_{\alpha(P)}(H)=\alpha(m_P(H))$.
%
%Let~$\Lambda:=m_P(H)\tens_{\wh{\Z}}\Z_\ell$, which is a $\Z_\ell$-submodule of~$T$. 
%
%By~\eqref{eq:0}, we have	
%\[
%\forall N\in\Z_{\geq1}, 
%\forall \alpha\in R,
%\alpha(\wh{T}(B))\leq N\cdot [E:K]\cdot c\cdot \wh{T}(B)\Rightarrow \alpha\in N\cdot R.
%\]
%
%We claim that
%\[
%\forall N\in\Z_{\geq1}, 
%\forall \alpha\in R\tens O_L,
%\alpha(\wh{T}(B))\leq N\cdot [E:K]\cdot c\cdot \wh{T}(B)\tens O_L\Rightarrow \alpha\in N\cdot R\tens O_L.
%\]
%\begin{proof}
%...
%\end{proof}
%\section{old}
%\subsubsection{Strategy - REDO}
%
%The strategy is to adapt Hindry's presentation~\cite{} of Ribet's arguments~\cite{}, with the introduction 
%of some modifications in order to obtain uniformity. We make also use of the "uniform integral Tate conjecture" Theorem~\ref{} version~\cite{}. We need to make effective the dependence of~$k$ upon Faltings' theorem~$F$. We give a uniform treatment of~$\ell$ using the linear algebra over~$\Z$...
%
%%We mention to the reader~\cite[App. C]{}
%%for uniformity results on Galois representation attached to points varying in a Hybrid orbit of a Shimura variety.
%
%
%%\subsection{Some effective Lemmas}
%%
%%
%%We know that~$\End_A$ is an order in~$\End^0_A:=\End_A\tens \Q$ and that~$\End^0_A$ is a semsimple
%%algebra of finite dimension over~$\Q$.
%%\begin{lemma}\label{lem:disc} 
%%There exists~$d_A$ such that for every order~$O\leq \End_A^0$ containing~$\End_A$, we have
%%\[
%%[O:\End_A]\text{ divides }d_A.
%%\]
%%\end{lemma}
%%\begin{proof}
%%Let~$B(a,b)=Tr(c\mapsto abc:\End^0_A\to\End^0_A$ denote the trace form. Let~$(\End_A)^{\bot}:=\{b\in \End^0_A|B(b,\End_A)\leq \Z\}$. Then
%%\[
%%(\End_A)^{\bot}\geq O^{\bot}\geq O\geq \End_A.
%%\]
%%Therefore~$[O:\End(A)]^2$ divides the \emph{discriminant}~$d_A:=[(\End_A)^{\bot}:\End_A]$.
%%\end{proof}
%%\begin{corollary}
%%Let~$\End_A\to\End(\Z^m)$ be a (faithful) representation. Then, for~$n\in\Z_{\geq1}$
%%\[
%%[\End_A\cap n\End(\Z^m):n\End_A]\text{ divides }d_A. 
%%\]
%%\end{corollary}
%%\begin{proof}
%%Apply Lemma~\ref{lem:disc} to~$O=\End_A\cap n\End(\Z^m)$.
%%\end{proof}
%%\begin{corollary}
%%For~$(n,d_A)=1$, the natural map 
%%\[
%%\End_A/n\End_A\to \End(\Z/(n)^m).
%%\]
%%is injective.
%%\end{corollary}
%%
%%\begin{lemma}
%%For~$\ell\geq \ell_1(A)$, the algebra~$\End_A\tens\F_\ell$ is a finite semisimple algebra. In particular every representation~$\End_A(\tens\F_\ell?)\to M(n,\F_\ell)$ is semi-simple, and its commutant is semisimple?
%%\end{lemma}
%%\begin{proof}References....
%%\end{proof}
%
%\subsection{Proof of Theorem~\ref{thm:uniform Kummer}}
%Passing to a finite extension of~$K$ we do assume that~$\End_A:=\End(A/K)=\End(A/\ol{K})$.
%
%We consider the situation of~Theorem~\ref{}. We define~$H:=\ker(\rho_A)$ and~$W:=m_P(H)\leq \wh{T}(B)$.
%
%Let~$S=...$ (from Faltings theorem). We have then~$\Delta(S)....\Delta(S_0) F()^?$.
%
%The Theorem~\ref{thm:uniform Kummer} is a consequence of Proposition~\ref{} below and Corollary~\ref{} from the appendix.
%
%\begin{proposition}\label{prop:xxx}
%There exists~$n(e(A,K),P,[E:K])\in\Z_{\geq1}$ such that
%\[
%\forall N\in\Z_{\geq1}, \alpha\in \End(B),\alpha (W)\leq N\wh{T}(B)\Rightarrow n\cdot \alpha\in N\End(B).
%\]
%where~$W:=m(Gal(\ol{K}/E^{ab}))$.
%\end{proposition}
%Let~$G:=Gal(\ol{K}/E^{ab})$ and let~$m_{\alpha}\in  \wh{T}(B)^{
%G}$ be the map
%\[
%m_\alpha:=\alpha\circ m|_G:G\to \wh{T}(B).
%\]
%Then~$\alpha (W)\leq N\wh{T}(B)$ if and only if
%\[
%m_\alpha\equiv 0\pmod{N\wh{T}(B)^G}.
%\]
%The statement~\eqref{} is equivalent to the conjonction of~$\alpha\mapsto m_\alpha$ is injective;
%and
%\[
%exponent(\wh{T}(B)^G/\{m_\alpha|\alpha\in \End(B)\})_{tors}|n.
%\]
%
%\begin{lemma}
%Let~$\phi:X\to Y$ be a morphism of abelian groups and~$n\in\Z_{\geq1}$.
%Then	
%\[
%\forall N\in\Z_{\geq0}, \alpha\in X,\phi(\alpha)\in n\cdot N\cdot Y\Rightarrow \alpha\in N\cdot X
%\]
%if and only if~$\phi$ is injective and
%\[
%n\text{ acts by zero on }(Y/\phi(X))_{tors.}\text{ and? if~$X$ not torsion free}\ker(\phi).
%\]
%\end{lemma}
%\begin{proof}We may replace~$Y$ by~$\{y\in Y|\exists N\in\Z_{\geq1}, Ny\in \phi(X)\}$.
%
%If~$n$ anihilate~$X/\phi(Y)$ then~$nY\leq \phi(X)$ and thus~$n\cdot N\cdot Y\leq N\cdot \phi(X)=\phi(N\cdot X)$.
%If moreover~$\phi$ is injective, this implies~\eqref{}.
%
%Assume~\eqref{}. For~$N=0$, we deduce that~$\ker(\phi)$ is annihilated by~$n$... injective.
%For~$N=1$ we deduce~$nY\leq \phi(X)$, hence~\eqref{}. 
%\end{proof}
%\subsubsection{Passing from~$\End(B)$ to~$A(E)$}
%
%
%\begin{proposition}We have an exact sequence
%\[
%NA(E)\to A(E) \xrightarrow{Q\mapsto m_Q:G_E\to \wh{T}\to B[N]} H^1(G_E;B[N])
%\] 
%\end{proposition}
%\begin{proof}See~\cite{}.
%\end{proof}
%\begin{proposition}\label{prop:LNM5}
%We have
%\[
%[E:K]\cdot (A(E)/A(K))_{tors}=0.
%\]
%Equivalently,
%\[
%\forall N\in\Z_{\geq1},\forall Q\in A(E), 
%Q\in NA(K)\Rightarrow Q\in [E:K]\cdot N\cdot A(E).
%\]
%\end{proposition}
%\begin{proof} We have commutative diagram
%\[
%...
%\]
%By~\cite[...]{LNM5} the kernel of the top row is of~$[E:K]$-torsion.
%Therefore the kernel of te bottom row is also of~$[E:K]$-torsion. The statements follow.
%\end{proof}
%\begin{lemma}
%The map
%\[
%\End(B)\mapsto \Theta:=\End(B)\cdot P
%\]
%is injective, and there is~$n(P,K)\in\Z_{\geq1}$ such that
%\[
%\forall N\in\Z_{\geq1},\forall \alpha\in \End(B),
%\alpha\not \in N\End(B)\Rightarrow \alpha(P)\not \in [E:K]\cdot A(E).
%\]
%\end{lemma}
%\begin{proof}Since~$P$ generates~$B$ as an algebraic group, we have~$\alpha(P)=0$ 
%if and only if~$\alpha(B)=0$. This proves the injectivity statement. We deduce
%\[
%\forall N\in\Z_{\geq1}, \forall \alpha\in \End(B), \alpha\in N\End(B)\Leftrightarrow \alpha(P)\in N\Theta.
%\]
%
%By Theorem~\ref{}, the group~$A(K)$ is finitely generated.
%Therefore~$(A(K)/\Theta)_{tors}$ is a finitely generated abelian group:
%there exists~$n(P,K)\in\Z_{\geq1}$ such that~$(A(K)/\Theta)_{tors}$ is~$n(P,K)$-torsion.
%
%Using Proposition~\ref{prop:LNM5}, we deduce that~$(A(E)/\Theta)_{tors}$ is~$([E:K]\cdot n(P,K))$-torsion.
%
%The Lemma follows.
%\end{proof}
%
%\subsubsection{Passing from~$G_E$ to~$H$}
%
%\begin{proposition}\label{prop:Sah}
%Let~$\rho:G\to GL(V)$ be a representation on a finitely generated profinite abelian group~$V$ (a quotient of~$\wh{\Z}^g$), and let~$r:G\to G^{ab}\to I$ be a morphism to an abelian group~$I$. Let~$G'$ be the image of~$(\rho,r)$. Let~$\sigma\in G$ and~$\lambda\in\wh{\Z}^\times$ be such that~$\rho(\sigma)$ acts by~$\lambda$ on~$V$.
%
%Then~$\sigma$ acts by~$G$-automorphism on~$V$ and~$\lambda$ acts trivially on~$H^1(G;V)$.
%\end{proposition}
%\begin{proof} As~$\rho(\sigma)$ is in the image of~$GL(1,\wh{\Z})\to Z(GL(V))$ and~$r(\sigma)\in Z(I)=I$, the element~$(\rho,r)(\sigma)\in G'$ is central in~$G'$. It acts~$G$-equivariantly on~$V$ and thus acts on~$H^1(G';V)$ by functoriality. 
%The conclusion follows from~\cite[Lemma~\ref{}]{Hin}.
%%It plainly acts on~$H^1(G';V)$ by~$\lambda=\rho(\sigma)\pmod{\ell}\in \F_\ell$. But~$\lambda\neq 1$ by assumption, and~$\lambda$ acts trivially on~$H^1(G';V)$ by~\cite[Lemma]{}.... Therefore~$H^1(G';V)=0$.  
%\end{proof}
%
%\begin{corollary}\label{cor:Sa ne} In Proposition~\ref{prop:Sah}, assume that for every~$\lambda\in \Z^{\times}$ there is~$\sigma\in G$ such that~$\sigma$ acts by~$\lambda^e$ on~$V$.
%Then, for every subquotient~$G$-module~$V''\twoheadleftarrow{} V'\hookrightarrow V$, 
%\[
%\text{$H^1(G;V'')$ is a~$n_e$-torsion abelian group:~$n_e\cdot H^1(G';V')=0\leq H^1(G';V')$.}
%\]
%\end{corollary}
%\begin{corollary}
%Let~$G'=(\rho,r)(G)$ and~$H=\ker(\rho)\cap \ker(r)= \ker(G\to G')$. Consider the corresponding exact sequence
%\[
%0\to H^1(G';V'')\to H^1(G;V'') \xrightarrow{\mathrm{res}_{G,H}} H^1(H;V'').
%\]
%Then 
%%and let~$M=\mathrm{res}(H^1(G;V))$. There exists~$\matrm{sec}:M\to H^1(G;V)$ such that~$\matrm{sec}\circ \mathrm{res}_{G,H}$ is the multiplication-by-$n_e$ map~$H^1(G;V)\toH^1(G;V)$. In particular
%\[
%\forall \gamma\in H^1(G;V''), n_e\cdot \gamma\neq 0\Rightarrow \mathrm{res}_{G,H}(\gamma)\neq 0.
%\]
%and
%\[
%\forall N\in \Z_{\geq1}, \forall \gamma\in H^1(G;V''), n_e\cdot \gamma\not \in N\cdot H^1(G;V'')\Rightarrow \mathrm{res}_{G,H}(\gamma)\not\in N\cdot H^1(H;V'').
%\]
%
%\end{corollary}
%\begin{proof}
%By Corollary~\ref{cor:Sah ne}, $H^1(G';V)$ is~$n_e$-torsion. But~$n_e\cdot \gamma\neq 0$. Therefore~$\gamma\not\in \ker(\mathrm{res}_{G,H})$.
%Replacing~$V''\twoheadleftarrow{} V'\hookrightarrow V$ with~$V''\tens\Z/(N)\twoheadleftarrow{}V''\twoheadleftarrow{} V'\hookrightarrow V$, yields~\eqref{}.
% +...
%\end{proof}
%
%\begin{proposition}Let~$G'$ and~$H$ be the image and kernel of~$(\rho_A,r):G_E\to GL(\wh{T}(A))\times G^{ab}_E$.
%Let~$e(A,K)\in\Z_{\geq1}$ be as in Theorem~\ref{thm:Serre eAK} and~$e:=e(A,K)\cdot[E:K]$ and~$n_e$ as in Lemma~\ref{lem:Lang ne}.
%Let~$N\in\Z_{\geq1}$.
%In the exact sequence
%\[
%0\to
%H^1(G';B[N])
%\to
%H^1(G_E;B[N])
%\to
%H^1(H;B[N])
%\]
%the group \(H^1(G';B[N])\) is of~$n_e$-torsion.
%
%For every~$Q\in A(E)$, we have
%\[
%Q\not\in NA(E)\Rightarrow m_{Q}\not\equiv 0\pmod{N\wh{T}(B)}.
%\]
%\end{proposition}
%\begin{proof}
%Sah lemma plus...
%\end{proof}
%\subsubsection{}
%We can now finish the proof of Proposition~\ref{prop:xxx}.
%\begin{proof}[Proof of Proposition~\ref{prop:xxx}]...
%We prove the contraposite implication. We need to prove that~$\alpha (W)\not\equiv 0\pmod{N\wh{T}}$.
%Equivalently, that the map~$m_{\alpha(P)}|_H\not\equiv 0\pmod{N\wh{T}}$.
%
%By... ... ...
%\end{proof}

%
%\section{OLD}
%
%
%\subsubsection{Effectivity lemmas}
%
%\subsubsection{}
%Let~$G=Gal(\ol{K}/E)$, let~$V'=\wh{T}(B)\leq V=\wh{T}(A')$, let~$\rho_A:G\to GL(V)$ be given by~\ref{}. 
%Let~$r:G\to G/[G,G] = Gal(E^{ab})$ be the quotient by the derived subgroup. 
%
%By Theorem~\ref{thm:Serre eAK} we may apply Lemma~\ref{lem:Lang ne} with~$e=e(A,K)\cdot [E:K]$.
%
%For every~$N\in \Z_{\geq1}$ the kernel of
%\[
%Q\in A(E) \mapsto m_Q\tens\Z/(N)\in H^1(G;B[N])
%\]
%is~$N\cdot A(E)$. Passing to the limit we have an injection
%\[
%A(E)\tens \wh{\Z}\to m_Q\in H^1(G;\wh{T}(B))
%\]
%such that
%\[
%\forall N\in\Z_{\geq 1}, Q\in A(E)\smallsetminus NA(E) \Rightarrow N\cdot m_Q|_H\not \in N\cdot H^1(H;P) ou image. 
%\]
%
%We may replace~$P$ by~$\#\pi_0([P])$. Then~$P\in B$ and~$P$ generates~$B$. Therefore~$\forall \alpha\in \End(B),\alpha(P)=0\Leftrightarrow \alpha(B)=0$.
%We have a bijection
%\[
%\End(B) \xrightarrow{\sim} \Theta:=\End(B)\cdot P\leq A(K)\leq A(E)
%\] 
%The maps
%\[
%\Theta\hookrightarrow A(K)\to A(E)\tens\wh{\Z} \to H^1(G;V')\to H^1(H;V')
%\] 
%are~$\End(B)$-equivariant.
%
%
%As~$X\cdot P$ generates ...
%
%\subsubsection{Proof of Theorem~\ref{thm:uniform Kummer} for large~$\ell$}
%The application~$m_{A,P}:L\sigma\mapsto m(\sigma)$ is a cocycle~$Gal(\ol{K}/K)\to \wh{T}(A'):=\Hom(\Gamma,\wh{T}(A))$.
%The map~$Q\mapsto m_{A,P}:A(K)\to H^1(G_K;\wh{T}(A'))\to H^1(G_K;\wh{T}(A')\tens\F_\ell)$ induces an injection
%\[
%A(K)/\ell A(K)\hookrightarrow H^1(G_K;\wh{T}(A')\tens\F_\ell).
%\]
%Let~$E_C$ be the reflex field of~$C$ and let~$E=E_C\cdot K$. Then~$[E:K]\leq [E:\Q]\leq d':=2^d\cdot d!$
%is bounded in terms of~$d$. Let~$T_C$ be the Mumford-Tate group and let~$r:G:=Gal(\ol{K}/E)\to Gal(\ol{Q}/E_C)\to T_E(\A_f)/T(\Q)$ be the reciprocity map. Let~$\rho:G\to Gal(\ol{K}/K)\to GL(\wh{T}(A))$ be the restriction of~$\rho_A$. We have~$H:=\ker(r)\leq Gal(\ol{K}/K(C_{tors}))\leq \ker(r)$.
%
%Let~$e(A;K;d')$ be as in Theorem~\ref{}. Let~$m(d')$ be as in Theorem~\ref{}.
%We apply Proposition~\ref{}: for~$\ell\geq e(A;K;d')$, we have~$H^1((\rho,r)(Gal(\ol{K}/E);A'[\ell])=0$, and
%\[
%A(E)/\ell A(E)\hookrightarrow H^1(G_K;\wh{T}(A')\tens\F_\ell)\hookrightarrow H^1(H;\wh{T}(A')\tens\F_\ell).
%\]
%Note that~$H$ acts trivially on~$\wh{T}(A')\tens\F_\ell$ and therefore~$H^1(H;\wh{T}(A')\tens\F_\ell)=\Hom(H;\wh{T}(A')\tens\F_\ell)$.
%
%
%We use the following (\cite[Prop.\,9, p\,)]{LMN}).
%\begin{lemma}Let~$\ell\nmid[G:H]$. Then
%\[
%H^1(G;V)\tens\Z_\ell\to H^1(H;V)\tens\Z_\ell
%\]
%is an isomorphism.
%\end{lemma}
%
%\begin{proposition}
%For~$\ell\nmid [E:K]$, we have
%\[
%A(K)/\ell A(K)\hookrightarrow A(E)/\ell A(E).
%\]
%\end{proposition}
%\begin{proof}The map~$A(K)/\ell A(K)\to A(E)/\ell A(E) \to H^1(G_K;\wh{T}(A')\tens\F_\ell)$ is identical to the map~$A(K)/\ell A(K)\hookrightarrow H^1(G_K;\wh{T}(A')\tens\F_\ell)\to H^1(G_E;\wh{T}(A')\tens\F_\ell)$. By assumption we may apply Lemma~\ref{}. This is an injective map. This implies that~\eqref{} is injective.
%\end{proof}
%
%\begin{proposition}Let~$\Theta:=\End_A\cdot P\leq A(K)$. There is~$\ell_2(A,P)\in\Z_{\geq1}$ such that for~$\ell\nmid \ell_2(A,P)$, we have
%\[
%P\Z/\ell\Z P\hookrightarrow \Theta/\ell\Theta\hookrightarrow A(K)/\ell A(K).
%\]
%\end{proposition}
%\begin{proof}This follows from the fact that~$A(K)$ is a finitely generated abelian group (see Theorem~\ref{}).
%\end{proof}
%\begin{corollary}For~$\ell\geq\max\{d';\ell_2(A,P);e(A,K;d')\}$, we have
%\[
%\Theta/\ell\Theta\hookrightarrow \Hom(\Gamma,\wh{T}(A)\tens\F_\ell).
%\]
%In particular~$m_P:H\to \Hom(\Gamma,\wh{T}(A)\tens\F_\ell)$ is not zero.
%\end{corollary}
%\begin{proof}
%\end{proof}
%
%For~$\ell\geq \mu(d')$, let~$\alpha$ be as in Theorem~\ref{}. We have~$0=\alpha(m_P)=m_{\alpha(P)}$.
%By Corollary~\ref{}, we have~$\alpha(P)\in \ell\Theta$. [TODO...] By ... we have~$\alpha=0$ or~$\alpha(P)=0$.
%
%In the first case, by~\eqref{}, we have ($\ell\geq ...\pi_0$)
%\[
%A[\ell]=\ker(\alpha:A[\ell]\to A[\ell])=M\leq B[\ell]\leq A[\ell].
%\]
%We deduce~$B=A$...
%
%In the second case we have~$\ol{[P]}^{Zar}\leq \ker(\alpha)$ and
%\[
%M=\ker(\alpha)\cap A[\ell]=B[\ell].
%\]
%This proves
%\[
%\forall \ell\geq \ell_0(A,P,K,d)=... M\tens \F_\ell=B[\ell].
%\]
%Using Nakayama lemma we deduce
%\[
%M\tens\Z_\ell=\wh{T}(B)\tens \Z_\ell.
%\]
%
%\begin{proposition}[QUESTION, prolly no]We have
%\[
%\End_A/\ell\End_A\hookrightarrow \Theta/\ell\Theta.
%\]
%\end{proposition}
%
%
%\subsubsection{Proof of Theorem~\ref{thm:uniform Kummer} for small~$\ell$}
%
%
%\subsection{OLD DISCARD}
%
%Let~$A$ be an abelian variety over a finitely generated extension~$K$ of~$\Q$ and let~$\ol{K}$ be an algebraic closure of~$K$. Let~$A_{tors}\leq A(\ol{K})$ and let~$G:=Gal(\ol{K}/K(A_{tors}))$ be the kernel of the action~$Gal(\ol{K}/K)\to Aut(A_{tors})$.
%
%We consider the~$\wh{Z}$ and~$\A_f$-linear Tate modules
%\[
%\wh{T}:=\wh{T}(A):=\varprojlim A_{tors}[N]\simeq \Hom(\Q/\Z,A_{tors})\approx\wh{\Z}^{2\dim(A)}
%\qquad  V:=V_f(A):=\wh{T}(A)\tens \Q\simeq \Hom(\Q,A_{tors})\approx A_f^{2\dim(A)}.
%\]
%We then have a short exact sequence.
%\[
%\wh{T}\hookrightarrow V\to A_{tors}.
%\]
%We denote by~$\wh{T}_C$ the inverse image of~$C$.
%
%
%Let~$C\leq A_{tors}$ be a finite subgroup and let~$c: A\to A/C$ be the quotient map. We have induced map.
%\[
%\wh{T}\hookrightarrow \wh{T}(A/C)\qquad
%V\xrightarrow{\sim} V_f(A/C).
%\]
%Then~$V\hookrightarrow V_f(A/C)$ is an isomorphism whic identifies~$\wh{T}_C$ with~$\wh{T}(A/C)$.
%
%\begin{lemma}\label{lem:pairing}
%Let~$\Gamma$ be a subgroup of~$A(K)$ and let~$\Gamma_0:=\{y\in A(\ol{K})|\exists N\in\Z_{\geq1}, N\cdot y\in \Gamma\}$.
%
%We have a pairing
%\[
%(,)_A:G\times \Gamma_0\to A_{tors}
%\]
%defined by
%\[
%(\sigma,y):=\sigma(y)-y.
%\]
%which factors through
%\[
%G\times \Gamma_0/A_{tors}
%\]
%\end{lemma}
%
%\begin{proof} We need to prove
%\begin{subequations}
%\begin{equation}\label{eq:lem0}
%\forall \sigma\in G,\forall y\in \Gamma_0, (\sigma,y)\in A_{tors}.
%\end{equation}
%\begin{equation}\label{eq:lem1}
%\forall \sigma\in G,\forall y,y'\in \Gamma_0, (\sigma,y+y')=(\sigma,y)+(\sigma,y').
%\end{equation}
%\begin{equation}\label{eq:lem2}
%\forall \sigma,\forall\sigma'\in G,y\in \Gamma_0, (\sigma+\sigma',y)=(\sigma,y)+(\sigma',y).
%\end{equation}
%\begin{equation}\label{eq:lem3}
%\forall \sigma\in G,\forall y\in A_{tors}, (\sigma,y)=0.
%\end{equation}
%\end{subequations}
%
%\begin{proof}[Proof of~\eqref{eq:lem0}] Let~$y\in\Gamma_0$. There exists~$N\in\Z_{\geq1}$ such that~$N\cdot y\in \Gamma\subseteq A(K)$.
%Therefore, for~$\sigma\in G\subseteq Gal(\ol{K}/K)$, we have~$N(\sigma(y)-y))=\sigma(Ny)-Ny=0$. As~$N\neq0$, we have~$\sigma(y)-y)\in A[N]\subseteq A_{tors}$.
%\end{proof}
%\begin{proof}[Proof of~\eqref{eq:lem1}] We have~$\sigma(y+y')-(y+y')=
%\sigma(y)+\sigma(y')-y-y'=(\sigma(y)-y)+(\sigma(y')-y')$.
%\end{proof}
%\begin{proof}[Proof of~\eqref{eq:lem3} and~\eqref{eq:lem2}] If~$\sigma\in G$, then, for~$y\in A_{tors}$, we have~$\sigma(y)-y=0$.
%
%By~\eqref{eq:lem0} we may substitute~$y$ with~$\sigma'(y)-y$. We have
%\[
%0=\sigma(\sigma'(y)-y)-(\sigma'(y)-y)=\sigma(\sigma'(y))-\sigma(y)-\sigma(y) +y +(y-y).
%\]
%We deduce
%\[
%\sigma(y)-y+\sigma'(y)-y=
%\sigma(\sigma'(y))-\sigma(y)-\sigma(y) -y.\qedhere
%\]
%\end{proof}
%This proves Lemma~\ref{lem:pairing}.
%\end{proof}
%
%\begin{proposition}Let~$[N]:V/\wh{T}\xrightarrow{\sim} V/N\cdot \wh{T}$ be the map induced by multiplication by~$N$. Then
%\[
%(\sigma,y):=\varprojlim_{N} [N](\sigma,\frac{1}{N}\cdot y)'_A
%\]
%is a well-defined pairing
%\[
%(,):G\times \Gamma_0/A_{tors}\to V.
%\]
%\end{proposition}
%We check that~$(\sigma,y)$ is well defined. We need to prove that, for~$N,M\in\Z_{\geq1}$, we have
%\[
%[NM](\sigma,\frac{1}{NM}\cdot y)'_A\equiv [N](\sigma,\frac{1}{N}\cdot y)'_A\pmod{N\wh{T}}.
%\]
%\begin{proof}
%Substituting~$y'=\frac{1}{MN}\cdot y$, let us first prove
%\[
%[M](\sigma,\cdot y')'_A\equiv (\sigma,\cdot M\cdot y)'_A\pmod{\wh{T}}.
%\]
%This is a consequence that
%\[
%V/\wh{T}\xrightarrow{\sim} V/N\cdot \wh{T}\to V/\wh{T}
%\]
%is the multiplication~$y'\mapsto My': V/\wh{T}\to V/\wh{T}$.
%\end{proof}
%
%Assume that~$E:=\End(A/\ol{K})=\End(A/K)$ and let~$E_0:=E\tens \Q$. Then~$E$ acts on~$A_{tors}$.
%We assume that~$\Gamma$ is a~$E$-module. We have then an action of~$E$ on~$\Gamma$,~$\Gamma_0$ and~$\wh{T}$
%and an action of~$E_0$ on~$\Gamma_0/A_{tors}\simeq \Gamma\tens\Q$ and on~$V$.
%
%The pairing~\eqref{} is right~$E_0$-linear, that is
%\[
%\forall \psi\in E_0, \forall \sigma\in G,\forall y \in \Gamma_0/A_{tors}, (\sigma,\psi(y))=\psi((\sigma,y)).
%\]
%\begin{proof}If~$\psi\in E$ and~$z\in \Gamma_0$, we have
%\[
%(\sigma,\psi(z))_A=\sigma(\psi(z))-\psi(z)=\psi(\sigma(z)-z).
%\]
%Therefore~$(\sigma,z)_A$ is right~$E$-linear. We deduce that~$[N](\sigma,\frac{1}{N}y)'_A$ is right~$E$ linear,
%and thus that~$(,)$ is right~$E$ linear. Since it is additive, it is~$\Q$-linear. It is thus~$E_0$-linear.
%\end{proof}
%
%We have a map~$\Hom(A,A/C)\to \Hom(V_f(A),V_f(A/C))$. Upon the identification~$V:=V_f(A)\simeq V_f(A/C)$, 
%we have the following
%\[
%\Hom(A,A/C)\simeq \{\psi\in E_0|\psi(\wh{T})\subseteq \wh{T}_C\}.
%\]
%We consider
%\[
%\Gamma_C:= \Hom(A,A/C)\cdot \Gamma\leq \Gamma_0/A_{tors}.
%\]
%
%Let us now assume that~$\Gamma=E\cdot x\simeq E$ for some~$x$. We have
%\[
%\Gamma_0/A_{tors}\simeq \Gamma\tens\Q\simeq E_0.
%\]
%Let~$[y]=\psi(x)$ with~$\psi\in E_0$. We have
%\[
%(\sigma,[y])=\psi((\sigma,x)).
%\]
%We deduce
%\[
%G\cdot y\pmod \wh{T}=... 
%\]
%\newpage

%\section{}
%
%Let~$A$ be an abelian variety over~$\C$ and let~$P\in A(\C)$. Let~$A_P:=(\ol{\Z\cdot P}^{Zar})^0$ be the abelian subvariety which is the neutral component of the algebraic subgroup generated by~$P$. Without loss of generality, we assume that~$P\in (\ol{\Z\cdot P}^{Zar})^0$.
%
%For an abelian variety~$A'$, let
%\[
%\Theta_{A'}:=\{\phi(P)|\phi\in \Hom(A_P,A')\}
%\]
%and
%\[
%\Sigma_{A'}:=
%\Theta_{A'}+A'_{tors}.
%\]
%and, for each~$d\in\Z_{\geq1}$,
%\[
%\Theta_{A',d}:=
%\{Q'\in A'|d\cdot Q'\in \Theta\}
%\]
%and
%\[
%\Xi_{A'}:=\{Q'\in A'|\exists n\in\Z_{\geq1}, n\cdot Q'\in \Theta\}=\bigcup_{d\in\Z_{\geq1}}\Theta_{A',d}.
%\] 
%To each~$Q\in\Xi$ we associate
%\[
%d_Q:=\min\{d\in\Z_{\geq1}|d\cdot Q\in \Sigma_{A'}\},\text{ and }
%d'_Q:=\min\{d\in\Z_{\geq1}|d\cdot Q\in \Theta_{A'}\}.
%\]
%By definition of~$A_P$ and~$\Theta_{A'}$, the map
%\[
%\varphi\mapsto \varphi(P):\Hom(A_P,A')\to \Theta_{A'}
%\]
%is a bijection. Therefore~$\Theta_{A'}$ is a free abelian group, and thus is torsion-free.
% Therefore~$\Theta_{A'}\cap A'_{tors}=\{0\}$. That is,~$\Theta_{A'}$ and~$A'_{tors}$
%are in direct sum in~$\Sigma_{A'}$.
%
%To~$Q$ we can associate the unique~$\phi_Q,\phi_Q'\in\Hom(A_P,A')$ and~$T_Q\in A_{tors}$ such that
%\[
%d_Q\cdot Q=\phi_Q(P)+T_Q\text{ and }d'_Q=\phi'_Q(P).
%\]
%Fix~$\kappa,e\in\Z_{\geq1}$. 
%The subgroup~$L_e:=\{\lambda^e|\lambda\in \widehat{\Z}\}\leq \GL(1,\wh{\Z})$ and its action on~$A'_{tors}$ 
%are defined in~\ref{}. We define
%\begin{align}
%F_Q&:=\phi_Q(A_P[d_Q])\\
%E(Q;\kappa)&:=Q+\kappa\cdot F_Q,\\
%\text{ and }
%E'(d_Q\cdot Q;e)&:=d_Q\cdot Q+L_e\cdot T.
%\end{align}
%We will prove the following.
%\begin{proposition}
%There exists an abelian sub-$S$-scheme~$\mathcal{B}_1\leq \mathcal{A}$ and~$d_1\in\Z_{\geq1}$ such that
%\[
%\mathcal{C}=\mathcal{C}+\mathcal{B}_1
%\]
%and 
%\[
%[d_1](\mathcal{C})=\mathcal{C}_1+\mathcal{B}_1
%\]
%where~$\mathcal{C}_1=\bigcup_{n\in\Z_{\geq0}}d_1\cdot \mathcal{C}_{s_n}\cap \Sigma_{A_n}$.
%
%There exists an abelian sub-$S$-scheme~$\mathcal{B}_2\leq \mathcal{A}$ and~$d_2\in\Z_{\geq1}$ such that
%\[
%\mathcal{C}_1=\mathcal{C}_1+\mathcal{B}_2
%\]
%and 
%\[
%[d_2](\mathcal{C}_1)=\mathcal{C}_2+\mathcal{B}_2
%\]
%where~$\mathcal{C}_2=\bigcup_{n\in\Z_{\geq0}}d_2\cdot {\mathcal{C}_1}_{s_n}\cap \Theta_{A_n}$.
%\end{proposition}
%
%Let~$K/\Q$ be an extension of finite type such that~$A$ is definable over~$K$ and~$P\in A(K)$.
% 
%The proof of Theorem~\ref{} will combine the following. 
%
%Firstly, there exists~$e,\kappa$, depending only on ... such that
%\[
%\forall s\in A, \forall Q\in \Xi_{A_s}, Gal(\ol{K}/K), E(Q;\kappa)\subseteq Gal(\ol{K}/K)\cdot (A_s,Q)
%\]
%and
%\[
%\forall s\in A, \forall Q\in \Sigma_{A_s}, Gal(\ol{K}/K), E'(Q;\kappa)\subseteq Gal(\ol{K}/K)\cdot (A_s,Q).
%\]
%
%Secondly, there exists an abelian $S$-subscheme~$\mathcal{B}\leq \mathcal{A}$ such that
%\[
%\forall s\in S, \mathcal{C}_s=\mathcal{C}_s+\mathcal{B}_s
%\]
%and, for a generic sequence~$s_0,\ldots$ in~$S$,
%\[
%\# E(Q_n,s_n) +B_{s_n}/B_{s_n}=O(1)
%\]
%
%Finally, 
%\[
%'''\]
%
%
%We recall that~$\Hom(A_P,A')$ is a free~$\Z$-module of finite type.
%One deduces that~$\Theta\cap A'_{tors}=\{0\}$.
%
%Consider
%\[
%Q\in \Xi:=\{Q'\in A'|\exists n\in\Z_{\geq1}, n\cdot Q'\in \Theta\}. 
%\]
%Then 
%\[
%d:=\min\{n\in\Z_{\geq1}|n\cdot Q_i\in \Sigma\}<+\infty,
%\]
%and there are unique~$\varphi:A_P\to A'$ and~$T\in A'_{tors}$ such that
%\[
%d\cdot Q = \varphi(P)+T.
%\]
%
%For~$\kappa,e\in\Z_{\geq1}$, we define
%\begin{align}
%F_Q&:=\phi_Q(\kappa\cdot A_P[d_Q])=\phi_Q(\kappa\cdot A_P[d_Q/\gcd(\kappa,d_Q]),\\
%E(Q;\kappa)&:=Q+\kappa\cdot F_Q,\\
%\text{ and }
%E'(Q;e)&:=Q+L_e\cdot T.
%\end{align}
%where the subgroup~$L_e\leq \GL(1,\wh{\Z})$ and its action on~$A'_{tors}$ 
%are defined in~\ref{}.
%
%\begin{proposition}
%Let~$Q\in\Xi$ and let~$B\leq A'$ be an abelian subvariety. 
%
%There exists~$T'\in A'_{tors}$ such that~$T'+B=T+B$ and the order of~$T'$
%is the order of~$T+B$ in~$A'/B$.
%
%There exists~$Q'$ such that
%\[
%Q\in E(Q',\kappa)+B\subseteq E(Q,\kappa)
%\]
%and
%\[
%d_Q...
%\]
%\end{proposition}
%\begin{proof}
%Let~$k$ be the order of~$T+B\in A'/B$. The map~$A[k]\to (A'/B)[k]$ is of the form~$(\Z/k\Z)^{2\dim (A')}\to(\Z/k\Z)^{2\dim (A')-2\dim(B)}$. It is surjective. Let~$T'\in A[k]$ be an antecedent of~$T+B$: we have~$T'\in T+B$ and therefore~$T'+B=T+B$.
%As~$T'\in A[k]$, the order of~$T'$ is at most~$k$. On the other hand, the order of~$T+B$ divides the order of~$T'$. We deduce that~$T'$ is of order~$k$.
%
%Let~$\pi:A'\to A'/B$ denote the quotient map.
%We have
%\[
%d\pi(Q)=\pi\circ\phi(P).
%\]
%Therefore~$d':=d_{\pi(Q)}$ divides~$d$. Let~$d''=d/d'$ for some~$d'\in\Z_{\geq1}$.
%Let us consider~$E(Q+B,\kappa)$. We have
%\begin{align}
%E(Q+B,\kappa)&=Q+B+\pi\circ\phi(A_P[d'])\cdot \kappa\\
%&=\pi(Q)+\pi\circ\phi(d''\cdot A_P[d])\cdot \kappa\\
%&\subseteq \pi(E(Q,\kappa)).
%\end{align}
%
%\end{proof}
%
%Let~$\mathcal{A}_\Gamma\to \mathbf{A}_\Gamma$ be a universal abelian scheme as in~\ref{}.
%We consider, for every~$i\in\Z_{\geq0}$
%\[
%(A_i,Q_i,\lambda_i)\in H_\Gamma(A,P). 
%\]
%and define the subsets
%\[
%\Theta_i:=\{\phi(P)|\phi\in \Hom(A,B_i)
%\]
%and
%\[
%\Sigma_i:=
%\Theta_i+(A_i)_{tors}.
%\]
%Then~$(A_i,Q_i,\lambda_i)\in H_\Gamma(A,P)$ implies that
%\[
%\]
%By definition of~$A_P$ and~$\Theta_i$, the map
%\[
%\varphi\mapsto \varphi(P):\Hom(A_P,A_i)\to \Theta_i
%\]
%is a bijection. We recall that~$\Hom(A_P,A_i)$ is a free~$\Z$-module of finite type.
%One deduces that~$\Theta_i\cap (A_i)_{tors}=\{0\}$. Therefore, there are unique~$\varphi_i:A_P\to A_i$ and~$T_i\in (A_i)_{tors}$ such that
%\[
%d_i\cdot Q_i = \varphi_i(P)+T_i.
%\]



\section{Proof of Theorem~\ref{thm:Main}}\label{sec:proof}

We consider the setting of Theorem~\ref{thm:Main} and the introduction. 
\subsection{}\label{sec:51}
Let~$A$ be an abelian variety over~$\C$ and let~$P\in A(\C)$. Let~$A_P:=(\ol{\Z\cdot P}^{Zar})^0$ be the abelian subvariety which is the neutral component of the algebraic subgroup generated by~$P$. Replacing~$P$ by the multiple~$n\cdot P$, where~$n=\#\pi_0(\ol{\Z\cdot P}^{Zar})$, we may assume~$\ol{\Z\cdot P}^{Zar}=A_P$.

Let~$S\subseteq \mathbf{A}_\Gamma$ be a special subvariety, and denote by~$\mathcal{A}:=\mathcal{A}_\Gamma|_S\to S$ the restriction to~$S$ of the abelian scheme~$\mathcal{A}_\Gamma\to\mathbf{A}_\Gamma$. 

Let~$\eta$ be a generic point of~$S$, and let~$A_\eta$ be the fiber at~$\eta$, which is an abelian variety over~$\C(S)$. We denote by~$A_{\ol{\eta}}$ the corresponding abelian variety over an algebraic and algebraically closed extension~$\ol{\C(S)}/\C(S)$. Passing to an étale cover~$\mathcal{A}_{\Gamma'}\to\mathcal{A}_{\Gamma}$ for a finite index congruence subgroup,
we have the following property
\begin{equation}
\End(A_\eta/\C(S))=\End(A_{\ol{\eta}}/\ol{\C(S)}).
\end{equation}
Equivalently, property~\eqref{eq:hypmonodromy} holds true for~$\mathcal{A}\to S$.

Let~$s_n=(A_n,\lambda_n)$ a generic sequence in~$S$, viewed as a sequence of polarised abelian varieties with~$\Gamma$-level structure. Let~$Q_n$ be a sequence of points in~$\mathcal{A}$ such that~$Q_n\in A_n$ for all~$n$. We consider~$\mathcal{V}:=\ol{\{Q_n\}}^{Zar(\mathcal{A})}$ and let~$\mathcal{C}$ be a geometric component of~$\mathcal{V}$ of non-zero dimension. Then~$\{n\in\Z_{\geq0}|Q_n\in \mathcal{C}\}$ is infinite.
Substituting~$s_n$ with the corresponding infinite subsequence, we may assume~$\mathcal{V}=\mathcal{C}$.

\subsection{Some finite sets}
For an abelian variety~$A'$, let
\[
\Theta_{A'}:=\{\phi(P)|\phi\in \Hom(A_P,A')\}
\]
and
\[
\Sigma_{A'}:=
\Theta_{A'}+A'_{tors}.
\]
and, for each~$d\in\Z_{\geq1}$,
\[
\Theta_{A',d}:=
\{Q'\in A'|d\cdot Q'\in \Theta_{A'}\}
\]
and
\[
\Xi_{A'}:=\{Q'\in A'|\exists n\in\Z_{\geq1}, n\cdot Q'\in \Theta_{A'}\}=\bigcup_{d\in\Z_{\geq1}}\Theta_{A',d}.
\] 
To each~$Q\in\Xi_{A'}$ we associate
\[
d_Q:=\min\{d\in\Z_{\geq1}|d\cdot Q\in \Sigma_{A'}\},\text{ and }
d'_Q:=\min\{d\in\Z_{\geq1}|d\cdot Q\in \Theta_{A'}\}.
\]
By definition of~$A_P$ and~$\Theta_{A'}$, the map
\[
\varphi\mapsto \varphi(P):\Hom(A_P,A')\to \Theta_{A'}
\]
is a bijection. Therefore~$\Theta_{A'}$ is a free abelian group, and thus is torsion-free.
 Therefore~$\Theta_{A'}\cap A'_{tors}=\{0\}$. That is,~$\Theta_{A'}$ and~$A'_{tors}$
are in direct sum in~$\Sigma_{A'}$.

To~$Q$ we can associate the unique~$\phi_Q,\phi_Q'\in\Hom(A_P,A')$ and~$T_Q\in A_{tors}$ such that
\[
d_Q\cdot Q=\phi_Q(P)+T_Q\text{ and }d'_Q\cdot Q=\phi'_Q(P).
\]
Fix~$\kappa,e\in\Z_{\geq1}$. 
Let~$L_e:=\{\lambda^e|\lambda\in \widehat{\Z}\}\leq \GL(1,\wh{\Z})$ act on~$A'_{tors}$. Let
\begin{align}\label{eq:def E Eprime}
F_Q&:=\phi_Q(A_P[d_Q])\\
\label{def:E Q kappa}
E(Q;\kappa)&:=Q+\kappa\cdot F_Q,\\
\label{def:E Q e}
\text{ and }
E'(d_Q\cdot Q;e)&:=d_Q\cdot Q+L_e\cdot T_Q.
\end{align}

\subsection{Some field extensions} Let~$K/\Q$ be a finitely generated extension such that~$A$ and~$S$ and~$\mathcal{C}$ are definable over~$K$ and~$P\in A(K)$ and~$\End(A/K)\simeq \End(A/\ol{K})$. Let~$n\in\Z_{\geq0}$ be arbitrary and let~$A':=A_{s_n}$. Let~$\phi\in\Hom(A,A')$ be such that~$C:=A'/\phi(A)$ has CM. 

We choose a model of~$A$ over~$K_1:=K$ and consider the action~$\rho_A:\Gal(\ol{K}/K)\to \Aut(A(\ol{K})_{\tors})$. 
Let~$\ol{K}/K_2/K_2'/K_1$ be defined by~$\Gal(\ol{K}/K_2)=\ker(\rho_A)$
and~$\rho(Gal(\ol{K}/K_2'))=\rho(Gal(\ol{K}/K))\cap GL(1,\wh{\Z})$ where~$GL(1,\wh{\Z})$ acts on~$A(\ol{K})_{\tors}$ by homotheties. 

As~$\ker(\phi)$ is stable under the action of~$GL(1,\wh{\Z})$ the quotient~$A\to A/\ker(\phi)$ is defined over~$K_2'$.


We consider the reflex morphism~$r:\Gal(E_*^{ab}/E_*)\to I_E:=\frac{(E\tens \A_f)^\times}{E^\times}$ and the subgroup~$I_\Q:=\frac{\A_f^\times}{\Q^\times}\leq I_E$. We define~$E_*^{ab}/E'_*/E_*$ by~$\Gal(E^{ab}/E'_*)=r^{-1}(I_\Q)$. We note that~$C$ is defined over~$E'$.

%We consider the action~$\rho_C$ of~$\Gal(E_*^{ab}/E_*)$ on~$\wh{T}(C)\tens \Q$ given by the main theorem of complex multiplication, and let~$E_*^{ab}/E'/E_*$ be such that~$\rho_{C}(\Gal(E^{ab}/E'))=\rho_C(\Gal(E^{ab}/E))\cap GL(1,\A_f)$. 

Let~$K_3:=E_*^{ab}\cdot K_2$ and~$K_3':=E'_*\cdot K_2'$.

There exists an isogeny
\[
\psi:(A/\ker(\phi))\times C\to A'.
\]
Note that~$(A/\ker(\phi))\times C$ is defined over~$K_3'$ and~$\ker(\psi)$ is stable under~$\Gal(\ol{K}/K_3')$. 
The quotient~$A'':=((A/\ker(\phi))\times C)/\ker(\psi)$ is thus a~$K'_3$-algebraic variety and the quotient map~$(A/\ker(\phi))\times C\to A''$ is defined over~$K'_3$.

We deduce that the image~$[A']=[A'']$ of~$s_n$ by~$S\to \mathbb{A}_g$ belongs to~$\mathbb{A}_g(K_3')$.
Let~$d_\Gamma$ be the degree of the map~$\mathbb{A}_\Gamma\to\mathbb{A}_g$. Then
\[
[K_3(s_n):K_3]\leq [K_3'(s_n):K_3']\leq 
[K(s_n):K([A'])]\leq d_\Gamma<+\infty.
\]

We claim that there exists~$c(\dim(A'))$ such that the isomorphism
\begin{equation}\label{eq:ol psi}
\ol{\psi}:A''\to A'
\end{equation}
induced by~$\psi$ is defined over an extension~$K'_4/K_3'(s_n)$ such that~$[K_4:K_3'(s_n)]\leq c(\dim(A'))$. Let~$K_4=K'_4\cdot K_3$. We have~$[K_4:K_3]\leq [K'_4:K_3']\leq d_\Gamma\cdot c(\dim(A'))$.
\begin{proof}[Proof of the claim]
 Both~$A''$ and~$A'$ are defined over~$L=K_3'(s_n)$. We remark that~$\End(A/K)\simeq \End(A/\ol{K})$ implies~$\End(A''/L)\simeq \End(A''/\ol{L})$, and implies that~$Gal(\ol{L}/L)$ acts trivially on~$Aut(A'')$. Therefore~$H^1(L;Aut(A''))\simeq Hom(Gal(\ol{L}/L);Aut(A''))$. Consider~$Isom(A'',A')$ as a~$Aut(A'')$-torsor. Its cohomology class is a morphism~$\phi:Gal(\ol{L}/L)\to Aut(A'')$. But there is a faithfull representation~$Aut(A'')\hookrightarrow GL(H^1(A(\C);\Z))\simeq GL(2g,\Z)$, and the image of~$\phi$ is finite. By Jordan's theorem,~$[Gal(\ol{L}/L):\ker(\phi)]\leq c(2g)$ for some~$c(2g)\in\Z_{\geq1}$. Therefore~$Isom(A'',A')$ has a~$L'$-rational point~$\psi'$ for some~$L'/L$ with~$[L':L]\leq c(2g)$. Therefore~$\ol{\psi}\in\psi'\cdot Aut(A'')(\ol{L'})=\psi'\cdot Aut(A'')(L)$ is defined over~$L'$.
\end{proof}

\subsubsection{Kummer-Ribet theory}
Let us prove the following.
\begin{theorem}\label{thm:Gal En} Let~$d_\Gamma\in\Z_{\geq1}$ be the degree of the map~$S\to \mathbb{A}_g$ and let~$\kappa'=k(A,K,P\cdot \Z;g)$ be as in Theorem~\ref{thm:uniform Kummer1}.
Let~$A'=A_{s_n}$ and let~$Q'\in\Xi_{A'}$. Let~$\Gal(\ol{K}/K)$ act on~$\mathcal{A}(\ol{K})$.
Then
\[
E(Q',\kappa)\subseteq \Gal(\ol{K}/K) \cdot Q'\cap A'(\ol{K})
\]
with~$\kappa'=(d_\Gamma\cdot c(\dim(A')))!\cdot \kappa$, and~$E(Q';\kappa)$ as in~\eqref{def:E Q kappa}, \eqref{eq:ol psi}.
\end{theorem}
\begin{proof}
By~Theorem~\ref{thm:uniform Kummer1},~$W:=m_P(\Gal(\ol{K}/K_3))\leq \wh{T}_{A_P}$ satisfies
\[
\kappa'\cdot\wh{T}_{A_P}\leq  W.
\]
Therefore, for every~$Q\in A(\ol{K})$ and~$p,q\in\Z_{\geq1}$ such that~$q\cdot Q=p\cdot P$ and~$\gcd(p,q)=1$, we have
\[
Q+\kappa'\cdot A_P[q]=Q+A_P[q/\gcd(q,\kappa')]\subseteq \Gal(\ol{K}/K_3)\cdot Q.
\]
Consider~$Q''\in A''(\ol{K})$,~$T\in A''(\ol{K})_{tors}$ and~$\phi'\in \Hom(A_P,A'')$ such that
\[
Q''=\phi'(Q)+T
\]
and~$\phi'$ is "primitive" in~$\Hom(A,A'')$, i.e.:
$
\forall m\in\Z_{\geq2}, \phi'\not\in m\cdot \Hom(A,A'').
$
Then
\[
\phi'(Q)+\kappa'\cdot \phi'(A_P[q])\subseteq \Gal(\ol{K}/K_3)\cdot \phi'(Q).
\]
As~$T\in A''(K_3)$ is fixed by~$\Gal(\ol{K}/K_3)$ we have also
\[
Q''+\kappa'\cdot \phi'(A_P[q])\subseteq \Gal(\ol{K}/K_3)\cdot Q''.
\]
Note that the action of~$\Gal(\ol{K}/K_3)$ is induced by a morphism~$\mu:\Gal(\ol{K}/K_3)\to \wh{T}_{A''}$
such that~$\kappa'\cdot \phi'(\wh{T}(A_P))$ is contained in the image of~$\mu$. Since~$[K_4:K_3]\leq d_\Gamma\cdot c(\dim(A'))$,
we deduce
\[
(d_\Gamma\cdot c(\dim(A')))!\cdot \kappa'\cdot \phi'(\wh{T}(A_P))\subseteq \mu(\Gal(\ol{K}/K_4)).
\]
This implies
\[
Q'+\kappa\cdot \phi'(A_P[q])\subseteq \Gal(\ol{K}/K_4)\cdot Q'.
\]
We recall that~$\ol{\psi}:A''\simeq A'$ is defined over~$K_4$. Therefore, for the action of~$\Gal(\ol{K}/K_4)$ on~$Q':=\ol{\psi}(Q'')\in A'(\ol{K})$ is induced by the action of~$\Gal(\ol{K}/K)$ on~$\mathcal{A}(\ol{K})$, we have
\[
Q'+\kappa\cdot \phi'(A_P[q])\subseteq \Gal(\ol{K}/K_4)\cdot 
Q'.\qedhere
\]
\end{proof}

\subsubsection{Lang-Serre theory} Let us prove the following.

\begin{theorem}\label{thm:Gal E'n} Let~$d_\Gamma\in\Z_{\geq1}$ be the degree of the map~$S\to \mathbb{A}_g$ and let~$e':=e(A,K,g)$ be as in Theorem~\ref{thm:Serre hybrid}.
Let~$A'=A_{s_n}$ and let~$Q'\in\Sigma_{A'}$. Let~$\Gal(\ol{K}/K)$ act on~$\mathcal{A}(\ol{K})$.
Then
\[
E'(Q';e)\subseteq \Gal(\ol{K}/K) \cdot Q'\cap A(\ol{K})
\]
with~$e=(d_\Gamma\cdot  c(\dim(A'')))!\cdot e'$, and~$E'(Q',e)$ as in~\eqref{def:E Q e}.
\end{theorem}
\begin{proof}Let~$A''$ and~$\ol{\psi}:A''\to A'$ be as in~\eqref{eq:ol psi}. Let~$H$ be the Hilbert class field of the reflex field of~$C$, and let~$E(C)/H$ be a minimal extension over which~$C$ is defined. 
By Theorem~\ref{thm:Serre hybrid}, we have
\[
\forall T\in A''_{tors},
L_{e'}\cdot T\subseteq Gal(\ol{K}/K'_3)\cdot T.
\]
We note that~$[K'_4\cdot E(C):K'_3]\leq d_\Gamma \cdot c(\dim(A''))$ and deduce
\[
\forall T\in A''_{tors},
L_{e}\cdot T\subseteq Gal(\ol{K}/K'_4)\cdot T.
\]
Note that~$\ol{\psi}:A''\to A'$ is defined over~$K'_4$. We deduce
\[
\forall T\in A'_{tors},
L_{e}\cdot T\subseteq Gal(\ol{K}/K'_4)\cdot T.
\]
Let us write~$Q'=R+T$ with~$R\in \Theta_{A'}$ and~$T\in A'_{tors}$. 
We have~$R=\phi(P)$ for some~$\phi\in \Hom(A_P,A')$. We claim that~$\phi$
is defined over~$K'_4$.
\begin{proof}Since~$\ol{\psi}$ is defined over~$K'_4$, it is enough to prove that every~$\phi\in \Hom(A_P,A'')$ is defined over~$K'_4$. From~$\End(A/K)\simeq \End(A/\ol{K})$ we deduce that every abelian subvariety of~$A$ is defined over~$K$. Therefore every abelian subvariety of~$A_P$, of~$A''$ and of~$A_P\oplus A''$ is defined over~$K'_3$. This holds in particular for the graph of~$\phi$. Therefore~$\phi$ is defined over~$K_3'\subseteq K'_4$.
\end{proof}
Therefore
\begin{multline}
Gal(\ol{K}/K'_4\cdot E(C))\cdot (R+T)\\
=R+Gal(\ol{K}/K_4'\cdot E(C))\cdot T\\
\supseteq
R+L_e\cdot T=E'(Q';e).\qedhere
\end{multline}
\end{proof}


%\subsubsection{Lang-Serre theory} By \S.~\ref{}, there exists~$e:=e(A,g)$ such that
%\[
%\forall T\in (A\times C)_{tors}, \forall \lambda \in L_e, \lambda\cdot T\in Gal(\ol{K}/K_3')\cdot T.
%\]
%More precisely, a subgoup~$H:=Gal(\ol{K}/L)\leq Gal(\ol{K}/K_3')$ acts on~$L_e\cdot T$ by an epimorphism~$H\to L_e$.
%It follows that the image of~$H\cap Gal(\ol{K}/K_3'K(s_n))$ contains~$L_{e\cdot [K(s_n):K]}\geq L_{e\cdot d_{\Gamma}!}$.
%
%Note that the morphism
%\[
%A\times C\to :(A/\ker(\phi))\times C\to A'
%\] is surjective and defined over~$K_4':=K_3'K(s_n)$.
%\begin{proof}
%\end{proof}
%Therefore every~$T'\in A'_{tors}$ has a preimage~$T\in (A\times C)_{tors}$ and moreover
%\[
%Gal(\ol{K}/K'_4)\cdot T'
%\]
%is the image of~$Gal(\ol{K}/K'_4)\cdot T$. We deduce that
%\[
%L_e\cdot T'\subseteq Gal(\ol{K}/K'_4)\cdot T'.
%\]
%Consider~$Q_n=\phi_n(P)+T_n\in \Sigma_{A'}$. 
%Then, the morphism~$\phi_{Q_n}:A\to A'$ is defined over~$K_4'$.
%\begin{proof}
%\end{proof}
%Therefore
%\[
%Gal(\ol{K}/K'_4)\cdot Q_n=
%Gal(\ol{K}/K'_4)\cdot (\phi_n(P)+T_n)=
%\phi_n(P)+Gal(\ol{K}/K'_4)\cdot T_n.
%\]
%We deduce
%\[
%E'(Q_n,e)=\phi_n(P)+L_e\cdot T_n\subseteq Gal(\ol{K}/K'_4)\cdot Q_n.
%\]

\subsection{Using equidistribution}
We use the notations from~\S\ref{sec:51}
We consider sequence $(E_n)$ defined as follows. Let $\kappa$ be as in Theorem \ref{thm:Gal En} and 
$$
E_n = E(Q_n, \kappa).
$$
Recall that $\mathcal{C}$ is defined over $K$ and 
$$
Q_n \in \mathcal{C}(\overline{K}).
$$
Therefore 
$$
Gal(\overline{K}/K) \cdot Q_n \subset \mathcal{C}.
$$
By Theorem~\ref{thm:Gal En},
$$
Q_n \in E_n \subset Gal(\overline{K}/K) \cdot Q_n \subset \mathcal{C}.
$$
By taking the Zariski closure in $\mathcal{A}$, we see that
$$
\mathcal{C}=(\bigcup Q_n)^{Zar} \subset (\bigcup E_n)^{Zar} \subset \mathcal{C}.
$$ 

By Theorem \ref{thm:weak}, there exists an abelian subscheme $\mathcal{B}$ of $\mathcal{A}$
such that
$$
\mathcal{C} = \mathcal{C} + \mathcal{B}
$$
and there exists a $d$ such that for every $n$, the image of $E_n$ by the quotient map
$A_n \to A_n/B_{s_n}$ is of cardinality at most $d$.




Let $B'$ be an abelian subvariety of the generic fibre $A_{\eta}$ which is a supplement of $B_{\eta}$ that is,
$A_{\eta}= B' + B_{\eta}$ and there exists $n_2$ such that $B' \cap B_{\eta} \subset A[n_2]$.

Let $\pi \in End(A_{\eta})$ be such that $B'=Im(\pi)$.
Recall that we assumed that $End(A_{\eta})=End(\mathcal{A})$, we can consider the abelian subscheme 
$\mathcal{B'}=Im(\pi)$. For every $s \in S$, $B'_s + B_s = A_s$ and $B'_s \cap B_s \subset A_s[n_2]$.

 Let $F_n = F_Q:=\phi_Q(A_P[d_Q])$ as in \eqref{eq:def E Eprime} with $Q=Q_n$. Let 
$\pi : A_n \to A_n/B_{s_n}$ be the quotient map.
Recall that $E_n = F_n + Q_n$. Therefore,
$$
\# \pi(F_n) = \# \pi(E_n)\leq d
$$   
and thus $\pi(F_n) \subset A_n/B_{s_n}[d!]$.

We now apply Theorem \ref{thm:E1} with $A_1=A$, $A'_1=A_P$, $A_2=A_n$, $A'_2=B_n$, $A''_2=A_n/B_n$,
$Q=Q_n$, $\phi=\phi_{Q_n}$ and $d=d_{Q_n}$ and $d' = gcd(d!,d_{Q_n})$.
The $n_2$ from that theorem is the same as above.

There exists $Q'_n = Q' \in Q_n + B_{s_n} \subseteq \mathcal{C}$ such that
$$
d' n_2 Q' = \phi'_n (P) + T_n
$$
with $\phi' \in Hom(A, A_n)$ and $T_n \in A_{n, tors}$.

Let $\mathcal{C}'$ be the closure of the set of $Q'_n$. We observe that 
$\mathcal{C}' \subset \mathcal{C} \subset \mathcal{C}' + \mathcal{B}$.

Therefore
\[\mathcal{C}' + \mathcal{B} = \mathcal{C} + \mathcal{B}= \mathcal{C}.\]


Note also that the closure of the set $\{d!\cdot n_2 Q'_n\}$ is $ \mathcal{C}'':=d!\cdot n_2 \mathcal{C}'$.

\subsubsection{}

Let $E_n'= E'(d!\cdot n_2 Q'_n, e)$ with $e$ as in Theorem \ref{thm:Gal E'n}. 
As before, we prove
\[
\mathcal{C}''=(\bigcup Q'_n)^{Zar} \subset (\bigcup E'_n)^{Zar} \subset \mathcal{C}''.
\]
By Theorem \ref{thm:weak}, there exists an abelian subscheme $\mathcal{B'}$ of $\mathcal{A}$
such that
$$
\mathcal{C}'' = \mathcal{C}'' + \mathcal{B}'
$$
and there exists a $d'$ such that for every $n$, the image of $E'_n$ by the quotient map
$\pi_n:A_n \to A_n/B'_{s_n}$ is of cardinality at most $d'$. Let us write~$d!\cdot n_2 Q'_n=\phi_n(P)+T_n$
with~$T_n\in {(A'_n)}_{tors}$ and~$\phi\in \Hom(A_P,A')$. Then the image of~$E'_n$
is
\[
\pi_n(E'_n)=\pi_n(\phi_n(P))+L_e\cdot \pi_n(T_n)
\]
and we have
\[
\#\pi_n(E'_n)=\omega_n^{1+o(1)}
\]
where~$\omega_n$ is the order of~$\pi_n(T_n)\in (A_n/B'_{s_n})_{tors}$. We deduce from~$\#\pi_n(E'_n)\leq d'$
that, for some~$d''\in\Z_{\geq1}$, we have~$\forall n\in\Z_{\geq1},\omega_n\leq d''$.

By Lemma~\ref{lem:C3} there exists~$T''_n\in {(A_n)}_{tors}$ of order~$\omega_n$ such that
$T''_n+B'_{s_n}= T_n+B'_{s_n}$. We have thus
\[
d''!\cdot \mathcal{C}''=\ol{\bigcup \phi_n(d''!\cdot P)+ d''!\cdot T_n+B'_{s_n}}^{Zar}=\ol{\bigcup \phi_n(d''!\cdot P)+ B'_{s_n}}^{Zar}
\]
and
\[
d''!\cdot \mathcal{C}''=d''!\cdot \mathcal{C}''+ \mathcal{B'}=
\ol{\bigcup \phi_n(P)}^{Zar}+ \mathcal{B}'.
\]
We deduce
\[
d''!n_2d'\cdot 
\mathcal{C}'=
\ol{\bigcup \phi_n(P)}^{Zar}+ \mathcal{B}'.
\]
and
\[
d''!\cdot n_2\cdot d!\cdot 
\mathcal{C}=
\ol{\bigcup \phi_n(P)}^{Zar}+ \mathcal{B}'+ \mathcal{B}.
\]

This proves Theorem~\ref{thm:Main}.
%\subsection{OLD:Some finite subsets of Galois orbits} Let~$K/\Q$ be a finitely generated extension such that~$A,P$ and~$\mathcal{C}$ are definable over~$K$. 
%
%Let~$d_\Gamma\leq [GSp(2,\Z):\Gamma]$ be the degree of the map~$\mathbf{A}_\Gamma\to \mathbf{A}_{GSp(2,\Z)}$.
%For every~$n\in\Z_{\geq0}$, let~$\phi_n:A\to A_n$ be such that~$C_n:=A_n/\phi_n(A)$ has CM. Then~$\phi_n(A)$ is definable over~$K(A_{tors})$.  
%the abelian variety~$A_n$ is definable over~$K(A_{tors})$
%
%\[
%K'_n=K((A_n)_{tors}))\quad K_n:=K'_n(s_n).
%\]
%The extension~$K_n/K$ is algebraic and usually~$[K_n:K]=+\infty$. Let~$E'_n/K$ be the smallest extension on which~$A_n$ is definable, and let~$E''_n$ be the reflex field of a~$CM$ quotient of~$A_n$ of maximal dimension.
%, and let~$E_n=E_n'(s_n)...\cdot E_n''$. Then~$[E'_n:K]\leq d_\Gamma$ and~$[E''_n:\Q]\leq 2^{g}\cdot g!$. Therefore
%\[
%[E_n:K]\leq d:=d_\Gamma\cdot 2^{g}\cdot g!.
%\]
%\[
%K'_n:=K((A_n)_{tors})\leq K(A_{tors})\cdot {E_n''}^{ab}(s_n)
%\]
%\[
%K_n=K_n'(s_n)\qquad
%[K_n:K'_n]
%\]
%We note that~$K_n/E_n({A_n}{tors})$ is an abelian extension. Let
%\[
%\kappa:=\kappa(A,\Z\cdot P; d)
%\]
%be as in~Theorem~\ref{thm:uniform Kummer2}. We consider the sequence 
%\[
%E_n:=E(Q_n;\kappa)
%\]
%of finite subsets of~$\mathcal{A}$. We have~$E_n\subseteq A_n$ for every~$n\in\Z_{\geq0}$.
%\begin{proposition} For every~$n\in\Z_{\geq1}$, we have
%\[
%E_n\subseteq Gal(\ol{K}/K)\cdot Q_n.
%\]
%\end{proposition}
%We will actually prove
%\begin{equation}\label{eq:En in Galois}
%E_n\subseteq Gal(\ol{K}/K_n)\cdot Q_n.
%\end{equation}
%\subsubsection{Proof of~\eqref{eq:En in Galois}}
%Let~$n$ be arbitrary, and let~$Q=Q_n$ and~$A'=A_n$ and~$\phi=\phi_{Q_n}$. Let~$\wh{T}$ be te Tate module of~$A_P$ and let~$W$ be the image of~$Gal(\ol{K}/K_n(A_{tors})$ by the map~\eqref{eq:Kummer m} and~$W'$ be the image of~$W$ by~$\wh{T}\to A_P[d_Q]$. Then~$\kappa\cdot \wh{T}\leq W$ by Theorem~\ref{thm:uniform Kummer2}. Let~$K(A')/K$ be the smallest extension on which~$A'$ is defined.
%~$Gal(\ol{K}/K(A'))$
%Moreover
%\[
%Gal(...)\cdot Q_n= Q_n+\phi_{Q_n}(W_{d_Q})
%\]
%...
%
%\subsection{Using equidistribution}...
%
%
%
%
%
%
%
%
%
%Let~$S\subseteq \mathbf{A}_\Gamma$ be a special subvariety, and denote by~$\mathcal{A}_\Gamma|_S\to S$ the restriction to~$S$ of the abelian scheme~$\mathcal{A}_\Gamma\to\mathbf{A}_\Gamma$. We consider an abelian scheme~$\mathcal{A}\to S$ such that we have the following. There exists~$k\in\Z_{\geq1}$ such that, denoting by~$\mathcal{A}^k_\Gamma|_S$ the~$k$-fold power of $\mathcal{A}^k_\Gamma|_S$ as an abelian~$S$-scheme, there exists a surjective morphism of abelian $S$-schemes
%\[
%\mathcal{A}^k_\Gamma|_S\to\mathcal{A}.
%\]
%In other words,~$\mathcal{A}\to S$ is a quotient (equivalently a subquotient) of a power of~$\mathcal{A}_\Gamma|_S\to S$.
%
%Let~$\eta$ be a generic point of~$S$, and let~$A_\eta$ be the fiber at~$\eta$, which is an abelian variety over~$K(S)$. We denote by~$A_{\ol{\eta}}$ the corresponding abelian variety over an algebraic and algebraically closed extension~$\ol{K(S)}/K(S)$. We consider the following properties
%\begin{equation}
%\End(A_\eta/K)=\End(A_\eta/\ol{K(S)})
%\end{equation}
%and, with~$\ol{K}$ the algebraic closure of~$K$ in~$\ol{K(S)}$,
%\begin{multline}
%\text{Every abelian subvariety~$0\neq B\leq A$}\\
%\text{is not definable over the subfield~$\ol{K}\leq \ol{K(S)}$,}
%\end{multline}
%and
%\begin{equation}\label{eq:finite degree moduli map}
%\text{the moduli map~$m_\mathcal{A}:S\to \mathbb{A}_{\dim(A_\eta)}$ has finite fibers.}
%\end{equation}
%
%
%Let~$s_n$ be a generic sequence in~$S$ and let~$B_n$ be the abelian variety~$\mathcal{A}_{s_n}$ above~$s_n$ in~$\mathcal{A}$. Let~$Q_n$ be a sequence of points in~$\mathcal{A}$ such that~$Q_n\in A_n$ for all~$n$.
%
%We claim that: if~\eqref{eq:finite degree moduli map}, then
%\[
%\forall s\in S, [K(s):K(A_s)]\leq \deg(m_\mathcal{A})<+\infty.
%\]
%\begin{proof}Let us observe that~$K(s_n)\geq K(B_n)\geq K(m_\mathcal{A}(s))$. The orbit of~$s$ under~$Aut\left(\ol{K(s)}/K(m_\mathcal{A}(s))\right)$ is contained in the inverse image of~$m_\mathcal{A}(s)$ by~$m_\mathcal{A}$, which is a finite set of cardinality at most~$\deg(m_\mathcal{A}(s))$. It follows that~$[K(s):K(m_\mathcal{A}(s))]$. The claim follows.
%\end{proof}
%We assume that~$Q_n\in H_\mathcal{A}(A,P)$ for all~$n$. We define
%\[
%V=\ol{\{Q_n\}}^{Zar(\mathcal{A})}
%\]
%and consider a geometric component~$C\subseteq V$. We assume that~$\dim(C)>0$.
%Then~$\{n\in\Z_{\geq1}|Gal(\ol{K}/K)\cdot Q_n\cap C\}\neq \emptyset$ is infinite, and~$C_\eta\neq \emptyset$.
%Let~$B_\eta\leq A_\eta$ be the maximal abelian $K(S)$-subvariety such that~$C_\eta(\ol{K(S)})=C_\eta(\ol{K(S)})+B_\eta(\ol{K(S)})$. We claim that
%\begin{itemize}
%\item this~$B_\eta$ is the generic fibre of an abelian subscheme~$\mathcal{B}\leq\mathcal{A}$
%\item and that
%\begin{equation}\label{asd}
%\forall s\in S(\C), C_s=C_s+B_s.
%\end{equation}\end{itemize}
%\begin{proof}...
%\end{proof}
%
%Let~$\ol{C}\subseteq \mathcal{A}/\mathcal{B}$ denote the image of~$C$ by the quotient map~$\mathcal{A}\to\mathcal{A}/\mathcal{B}$. Let~$\ol{Q}_i$ be the image of~$Q_i$ by the map~$A_{s_i}\to\ol{B}_i:= A_{s_i}/B_{s_i}$. By~\eqref{eq:asd},
%te variety~$C$ is the inverse image of~$\ol{C}$.
%For~$d\in \Z$, we denote by~$d\cdot Q_n$ denote the multiple of~$Q_n$ in the abelian variety~$A_n$. We observe that
%\begin{equation}
%\text{$Q\in B_n[d]$ for some~$d_n\in \Z_{\geq1}$, chosen minimal,}
%\end{equation}
%or there exists~$d_n\in \Z_{\geq1}$, and~$\phi\in \Hom(A_P,B_n)$ such that
%\begin{equation}
%d_n\cdot Q_n\in \phi_n(P)+B_{tors}, 
%\end{equation}
%where~$d_n$ is chosen to be minimal.
%
%We also assume that~$s_n$ is a strict sequence in~$S$.
%
%Let~$V\subseteq\ol{\{Q_n\}}^{Zar}$ be an irreducible component of the closure of the sequence~$Q_n$ in the algebraic variety~$\mathcal{A}$ for the Zariski topology.
%
%\begin{theorem}(passer a revetement fini)
%There exists~$d\in\Z_{\geq1}$ an abelian sub~$S$-scheme~$\mathcal{B}\leq \mathcal{A}$ such that~$V=V+\mathcal{B}$ 
%and that the image~$W$ of~$V$ in~$\mathcal{A}/\mathcal{B}$ satisfies
%\[
%\ol{d\cdot W\cap M_{\mathcal{A}/\mathcal{B}}(A,P)}=W.
%\]
%\end{theorem}
%\subsection{Kummer orbits in a hybrid orbit}
%
%\begin{lemma}
%Let~$d,e\in\Z_{\geq1}$, let~$G:=(\Z/d\Z)^e$ and let~$F\leq G$ be a subgroup.
%Assume that for every~$1<D\leq d$, we have
%\[
%G[D]\not\leq F.
%\]
%Then there exists~$x\in G/F$ of order~$d$. In particular
%\[
%\#G/F\geq d.
%\]
%\end{lemma}
%\begin{proof}We can factor~$G=\prod_{p^{m_p}\mid\mid d}(\Z/d\Z)^e$ and argue prime-by-prime.
%That is, we may assume~$d=p^m$ for some prime~$p$ and~$m\in\Z_{\geq1}$. Let~$g_1,\ldots,g_e$
% generate~$G$. One remarks that~$G[p]$ is generated by~$g_1^{p^{m-1}},\ldots,g_e^{p^{m-1}}$.
%Assume by contradiction that none of the~$x_i:=g_i\pmod{F}$ is of order~$p^m$. Then, for each~$1\leq i\leq e$, the order of~$x_i$ divides~$p^{m-1}$. That is,~$g_i^{p^{m-1}}\in F$. We deduce that~$G[p]\leq F$, contradicting the assumption.
%\end{proof}
%
%\begin{lemma}Let~$A,B$ be abelian varieties over~$\C$ and let~$\phi\in\Lambda:=\Hom(A,B)$.
%For~$N\in \Z_{\geq1}$, the following are equivalent:
%\begin{itemize}
%\item 
%We have~$\phi\in N\Lambda$.
%\item There exists~$\psi:A/A[N]\to B$ such that~$\phi=\psi\circ \pi_N$, where~$\pi_N$ is the quotient map~$A/A[N]$.
%\item We have~$A[N]\leq \ker(\phi)$.
%\end{itemize}
%For every~$\phi\in\Lambda\smallsetminus\{0\}$, there is a unique~$n(\phi)\in\Z_{\geq1}$ such that
%\[
%\forall N\in\Z_{\geq1}, A[N]\leq \ker(\phi)\Leftrightarrow 
%\phi \in N\Lambda\Leftrightarrow n(\phi)\in N\Z.
%\]
%\end{lemma}
%
%\begin{corollary}For every~$d\in \Z_{\geq1}$ and~$\phi\in\Lambda$, the group~$\phi(A[d])$ contains an element of order~$d/\gcd(d,n(\phi))$. In particular,~$\#\phi(A[d])\geq d/\gcd(d,n(\phi))$.
%\end{corollary}
%
%\begin{theorem}Fix~$(A,P)$ a pointed abelian variety over a finitely generated extension~$K/\Q$, and let~$\widehat{T}$ be the Tate module of the abelian variety~$A_P:=(\ol{P\cdot \Z}^{Zar})^0$. Let~$E/K$ be a finite extension and let~$K'=E^{ab}(A(\ol{K})_{tors})$. Then there exists~$\kappa\in\Z_{\geq1}$ depending on~$(A,P,K,[E:K],g)$, and a subgroup~$\kappa\cdot\wh{T}\leq \widehat{F}\leq \wh{T}$ such that for every~$B\in Hyb_{\mathbb{A}_g}(A)$, and every non torsion~$Q\in B$ such that
%\[
%d\cdot Q\in \phi(P)+B_{tors}
%\]
%for some~$d\in\Z_{\geq1}$ and non-zero~$\phi\in Hom(A_P,B)$ such that~$gcd(n(\phi),d)=1$, we have 
%\[
%Gal(\ol{K}/K')\cdot Q=Q+\phi(F)
%\]
%where~$F$ is the image of~$\wh{F}$ by~$\wh{T}\to A_P[d]$, and
%\[
%\#\phi(F)\geq d/\gcd(d,\kappa).
%\]
%\end{theorem}
%%%%%%%%%%%%%%%%%%%%%%%%%%%%%%%5
\appendix
\section{}
\begin{lemma}\label{lem:liftkernel}
Let~$\varphi:X\to Y$ be a homomorphism of abelian groups. Assume that~$\ker(\phi)$ is annihilated by~$n\in\Z_{\geq1}$. Then
\[
\varphi(x)\mapsto n\cdot x:\varphi(X)\to X 
\]
is a well defined homomorphism, and
\[
\forall x\in X, n\cdot x\neq 0\Rightarrow \phi(x)\neq 0.
\]
\end{lemma}
\begin{proof}[Proof of Lemma~\ref{lem:liftkernel}]If~$\varphi(x)=\varphi(x')$, then~$\varphi(x-x')=0$ and thus~$n(x-x')=nx-nx'=0$. For~$\phi(x),\phi(x')\in \varphi(X)$, we have~$\varphi(x)-\varphi(x')=\varphi(x-x')\mapsto n\cdot (x-x')=n\cdot x- n\cdot x'$. The claim follows immediately.
\end{proof}
%
%\begin{lemma}Let~$\phi:X\to Y$ be a homomorphism of abelian groups. For every~$N\in\Z_{\geq0}$,
%let~$\phi_N:X/N\cdot X\to Y/N\cdot Y$ be the induced morphisms. Let~$n\in\Z_{\geq 1}$.
%
%
%Assume that~$\coker(\phi)_{tors}$ is annihilated by~$n$. Then
%\[
%\forall N\in\Z_{\geq1}, \forall x\in X, \phi(x)\in n\cdot N\cdot Y\Rightarrow \phi(x)\in N\cdot \phi(X).
%\]
%Let~$n\in\Z_{\geq 1}$ be such that
%\[
%\forall N\in\Z_{\geq1}, \text{ $\ker(\phi_N)$ is annihilated by~$n$.}
%\] 
%Then~$\ker(\phi)$  is annihilated by~$n$.
%\end{lemma}


\section{Some linear algebra and arithmetic stability}\label{App:B}
The aim of this appendix is to prove Theorem~\ref{thm:D1}. This result is used in~\S\ref{sec:proof 4.10} for the proof of Theorem~\ref{thm:uniform Kummer2}.

\subsection{Split case}
\begin{theorem}\label{thm:split linear}
Let~$L/\Q_\ell$ be a finite extension, let~$k,n_1,\ldots,n_k\in\Z_{\geq1}$ and let~$V_1,\ldots,V_k$ be finite dimensional $L$-vector spaces.
Let~$W\leq V= \bigoplus_{i=1}^k L^{n_i}\tens V_i$ be an~$O_L$-submodule.  We consider the action
\[
E:=\prod_{i=1}^k \End(O_L^{n_i})\tens Id_{V_i}\leq \prod_{i=1}^k\End(L^{n_i})\tens \End(V_i)\leq \End(V).
\]
Let~$\lambda\in L$ be such that
\[
\lambda \cdot E\cdot W\leq W.
\]
Then there are~$O_L$-modules~$\Lambda_i\leq V_i$ such that
\[
\lambda\cdot\bigoplus_{i=1}^k O_L^{n_i}\tens_{O_L} \Lambda_i\leq W
\text{ and }
\lambda\cdot W\leq \bigoplus_{i=1}^k O_L^{n_i}\tens_{O_L}\Lambda_i.
\]
\end{theorem}

\begin{proof}[Proof of theorem~\ref{thm:split linear}]
Let~$e_{i,1},\ldots,e_{i,n_i}$ be the standard basis of~${L}^{n_i}$, let
\[
\eps_{i,(m,m')}:= e_{i,m}\times {}^te_{i,m'}\in \End(L^{n_i}):(\lambda_n)_{1\leq n\leq n_i}\mapsto (\delta_{m,n} \lambda_{m'})_{1\leq n\leq n_i}
\]
be the standard elementary matrices and let~
\[
\pi_{i,m}=\eps_{i,(m,m)}\in \End(L^{n_i}):(\lambda_n)_{1\leq n\leq n_i}\mapsto (\delta_{m,n} \lambda_n)_{1\leq n\leq n_i}
\]
be the~$m$-th standard projector.

We can uniquely decompose any~$v\in V$ as
\[
v=\bigoplus_{i=1}^k \sum^{n_i}_{m=1} e_{i,m}\tens v_{i,m}.
\]
Let~$\Lambda_{i,m}:=\{v\in V_i|e_{i,m}\tens v\in W\}$. If~$v\in \Lambda_{i,m}$,
then
\[
e_{i,m'}\tens \lambda \cdot v = \lambda \eps_{i,(m',m)}(e_{i,m}\tens v)\in \lambda\cdot E(W)\leq W.
\]
Therefore~$\lambda\cdot \Lambda_{i,m}\leq \Lambda_{i,m'}$. We deduce
\[
\Lambda_i:=\sum_{i=1}^{m}\Lambda_{i,m'}\geq \bigcap_{i=1}^{n_i} \Lambda_{i,m}\geq \lambda\cdot \Lambda_i.
\]
By construction
\begin{equation}\label{eq:356}
\bigoplus_{i=1}^k O_L^{n_i}\tens \lambda\cdot \Lambda_i
\leq 
\bigoplus_{i=1}^k \sum^{n_i}_{m=1} e_{i,m}\tens \Lambda_{i,m}
\leq W.
\end{equation}
For~$w\in W$, we have
\(
e_i\tens \lambda w_{i,m}=\lambda \pi_{i,m}(w)\in W.
\)
Therefore~$\lambda\cdot w_{i,m}\in \Lambda_{i,m}$ and
\[
\lambda\cdot w=\sum_{1\leq i\leq k,1\leq m\leq n_i}e_{i,m}\tens \lambda\cdot w_{i,m}\in 
\bigoplus_{i=1}^k \sum^{n_i}_{m=1} e_{i,m}\tens \Lambda_{i,m}\leq \bigoplus_{i=1}^k O_L^{n_i}\tens \Lambda_i.
\]
As~$w\in W$ is arbitrary, we conclude
\begin{equation}\label{eq:123}
\lambda\cdot W\leq \bigoplus_{i=1}^k O_L^{n_i}\tens \Lambda_i.
\end{equation}
The Theorem follows from~\eqref{eq:356} and~\eqref{eq:123}.
\end{proof}
\begin{corollary}Let us choose a~$L$-homogeneous norm on each~$V_i$ and equip~$V$ with the corresponding norm. 
Let~$e_i\in (O^\times_L)^{n_i}$ for~$i=1,\ldots,k$. Then
\[
\abs{\lambda}^2\cdot \norm{W}\leq \max_{1\leq i\leq k}\norm{(e_i\tens V_i)\cap W}
\]
\end{corollary}
\begin{proof}For~$e\in O_L^{n_i}$, we have~$\lambda\cdot (e\tens \Lambda_i)\leq (e\tens V_i)\cap W$. Therefore
\[
\abs{\lambda}\cdot \norm{e}\cdot \norm{O_L^{n_i}\tens\Lambda_i}=
\abs{\lambda}\cdot \norm{e\tens \Lambda_i}\leq \norm{(e\tens V_i)\cap W}.
\]
For~$e\in (O_L^\times)^{n_i}$ we have~$\norm{e}=1$.
We deduce
\[
\abs{\lambda}^2 \cdot \norm*{W}\leq \abs{\lambda}\cdot \norm*{\bigoplus_{i=1}^k O_L^{n_i}\tens \Lambda_i}=
\abs{\lambda}\cdot \max_{1\leq i \leq k}\norm*{O_L^{n_i}\tens \Lambda_i}\leq 
\max_{1\leq i \leq k}\norm*{(e\tens V_i)\cap W}.
\]
\end{proof}
\subsubsection{Applications}
Let us mention the case~$\lambda=1$. 
\begin{corollary}
A~$O_L$-submodule~$W\leq \bigoplus_{i=1}^k L^{n_i}\tens V_i$ is of the form~$W=\bigoplus_{i=1}^k O_L^{n_i}\tens_{O_L} \Lambda_i$ if and only if it is~$\prod_{i=1}^kGL(n_i,O_{L})$-invariant.
\end{corollary}
This contains criterions for a lattice to be a direct sum, or a tensor product. For instance the following.
\begin{corollary}Let~$W\leq V:= V_1\tens V_2$ be a lattice. Then~$W$ is of the form~$W_1\tens W_2$ if and only if the maximal compact subgroup~$GL(W)\leq GL(V)$ contains a maximal subgroup of~$GL(V_1)\tens Id_{V_2}$. The latter is necessarily~$GL(W)\cap GL(V_1)\tens Id_{V_2}= GL(W_1)\tens Id_{V_2}$.

The maximal compact subgroups~$GL(W)\leq GL(V)$ containing~$GL(W_1)\tens Id_{V_2}$ are the~$GL(W_1\tens W_2)$.
They form an orbit under~$Id_{V_1}\tens GL(V_2)$.
\end{corollary}

%\subsection{More general case}
%
%Let~$S$ be an order in a semisimple algebra~$S_0$ over~$\Q_\ell$.  Let~$\tau:S_0\times S_0\to \Q_\ell$ denote the Trace form, that is~$\tau(s_1,s_2):=Tr(s\mapsto s\cdot s_1\cdot s_2:S_0\to S_0)$. Let~$S^\bot:=\{s\in S_0|\tau(s,S)\leq \Z_\ell)$ be the corresponding Legendre dual. We have~$S\leq S^\bot$ because~$S$ is an order. The  discriminant of~$S$ is~$d_S:=[S^\bot:S]$.
%Let~$\delta_S=\ell^k$ denote the exponent of~$S^\bot/S$: one has~$\delta_S\cdot S^\bot\leq S$. One has the trivial bound~$\delta_S\mid d_S$.
%
%Let~$L/\Q_\ell$ be a finite extension, and let~$S\tens O_L\leq S'$ be a maximal order of~$S\tens L$. With respect to the trace form~$\tau\tens O_L$,
%\[
%S\tens_{\Z_\ell} O_L\leq S'\leq S'^\bot\leq (S\tens_{\Z_\ell} O_L)^\bot= S^\bot \tens_{\Z_\ell} O_L.
%\]
%Therefore,~$\ell^k\cdot S'\leq \ell^k\cdot S'^\bot\leq \ell^k\cdot S^\bot \tens_{\Z_\ell} O_L\leq S\tens_{\Z_\ell} O_L$. Assume that~$S$ is split over~$L$: one has~$S\tens L\approx \prod \End(L^{n_i})$. Then every maximal order is conjugated to~$\prod \End(O_L^{n_i})$ and satisfy~$S'^\bot=S'$.
%
%
%
%\begin{theorem}Let~$S\leq \End(\Q_\ell^n)$ be an order in a semisimple subalgebra~$S_0$,
%let~$R$ be the commutant of~$S$ in~$\End(\Z_\ell^n)$ and~$R_0:=R\tens\Q$. Let~$V\leq\End(\Q_\ell^n)$
%the~$\Q_\ell$-linear subalgebra generated by~$R\cup S_0$.
%
%Let~$W\leq V$ be a lattice such that~$\ell^{k'}\cdot S\cdot W\leq W$ and~$\ell^k\cdot R\cdot W\leq W$.
%
%Then
%\[
%\norm{W\cap R_0S}\leq \norm{W\cap R_0}\cdot \ell^{-?}.
%\]
%
%
%Let~$L$ be a splitting field of~$S$, and let~$S'$ be a maximal order of~$S\tens L$
%containing~$S$. Let~$\lambda\in L$ be such that~$\lambda \cdot S'\leq S\tens O_L$.
%
%Let~$V:=\Q_\ell^n\simeq \bigoplus M_i\tens V_i$ be a decompososition into irreducible,
%and~$W_i\leq M_i$ ,,,
%Let~$W\leq \Q_\ell^n$ be a~$S$-module.
%
%
%\end{theorem}
%
%\subsubsection{Application}
%\begin{proposition}\label{prop:tensor linear}
%Let~$T:=\Z^r\leq V:=\Q^r$, let~$S$ be an order in a semisimple subalgebra~$S_0\leq \End(\Q^r)$ such that~$S\cdot T\leq T$, and let~$R_0=\End_S(V)$. There exists~$\Delta(S)$ such that for every~$S$-stable group~$W\leq T$, and defining
%\[
%W^\vee:=\{\phi\in \Hom(V,\Q)|\phi(W)\leq \Z\},\quad 
%W':=\{\alpha\in R_0|\alpha(W)\leq T\},
%\]
%we have
%\begin{equation}
%\forall n\in \Q, n\cdot W'\leq \End(T)
%\Rightarrow n\cdot \Delta(Z) \cdot W^\vee\leq T^\vee.
%\end{equation}
%\end{proposition}
%\begin{proof}
%We consider
%\[
%W^\ast:=\{\alpha\in \End_\Q(V)|\alpha(W)\leq T\}\simeq \Hom(W,T)\simeq W^\vee\tens T.
%\]
%Then~$N\cdot W^\vee\leq T^\vee \Leftrightarrow N\cdot W^\ast\leq \End(T)$. We have
%\[
%W'=W^\ast\cap \End_Z(V)=W^\ast\cap R^0.
%\]
%Note that~$W^\star:=\Hom(W,T)\cap R^0\tens S^0$ is right~$S$-invariant. 
%Recall that~$R^0\tens S^0...$. By Theorem~\ref{}, we have
%\[
%\norm{W^\star}\leq cte(S)\cdot \norm{W'}.
%\]
%Let~$V=\oplus V_i$ be the decomposition into $S^0$-isotypic components,
%that is into~$R^0\tens S^0$-irreducible components. Let~$V^\vee=\oplus V_i^\vee$ dually.
%We have~$R^0\tens S^0=\bigoplus V^\vee_i\tens V_i$,
%and
%\[
%W_i\geq \leq x \bigoplus W\cap V_i.
%\]
%Let~$L=\prod L_i$ be the center of~$R^0\tens S^0$. Let~$\tau_i:L_i\to \Q$ be the (normalised)
%trace. We have
%\[
%R^0\tens S^0\cap --- vs \End(V_i) vs V_i^\vee.
%\]
%We consider
%\[
%W^\star:=\{\alpha\in \End_Z(V)|\alpha(W)\leq T\}\simeq \Hom_Z(W,T).
%\]
%
%\end{proof}
%
%\begin{corollary}
%\begin{equation}
%\forall n\in \Q, n\cdot W'\leq \End(T)
%\Rightarrow n\cdot \Delta(Z) \cdot T\leq W.
%\end{equation}
%\begin{proof} This follows from~\eqref{} and te duality property
%\[
%n\cdot \Delta(Z) \cdot T\leq W\Leftrightarrow 
%n\cdot \Delta(Z) \cdot W^\vee\leq T^\vee.\qedhere
%\]
%\end{proof}
%
%\end{corollary}
%%\begin{proof}[Proof of Proposition~\ref{prop:tensor linear}]
%%Let~$d_Z$ be the discriminant of~$Z$. [...] Then
%%\[
%%\forall z\in Z_0, z\cdot W\geq d_Z\cdot H_f(z)\cdot W.
%%\]
%%Therefore, for~$(\alpha,z)\in W'\times Z$ and~$m\in\Q$ such that~$n\alpha\in \End(T)$, we have
%%\[
%%n\cdot d_Z\cdot \alpha\cdot z\in \End(T)...
%%\]
%%\end{proof}
%
%
%
%\subsection{Arithmetic stability for general reductive groups} 
%Previous \S\S\ref{} can be considered as the case of reductive groups of Dynkyn type~$A_n$ of the following problematic.
%
%Let~$H\to G$ be a embedding of reductive groups over~$\Q_\ell$.
%A question is to determine the maximal compact subgroups~$K'\leq G(\Q_\ell)$ containing a given open compact subgroup~$K\leq H(\Q_\ell)$. 
%
%
%
%\subsubsection{Archimedean case}
%At the archimedean place,~$\R=\Q_\infty$, this is a classical question: which Cartan involutions~$\Theta':G(\R)\to G(\R)$ extend a given to a Cartan involution to~$\Theta:H(\R)\to H(\R)$. Firstly, there always exists such~$\Theta'$, by Mostow.
%Secondly, any other~$\Theta''$ is of the form~$\Theta''=g\circ \Theta$ where~$g$ is an automorphism of~$G$ centralising~$H$. Thirdly,~$g$ is in the neutral component~$G^{ad}(\R)^+\simeq Aut(G)^+$. Thus~$g=AD(z)$ with~$z\in Z_{G(\R)^+}(H)$. Lastly, the stabiliser of~$\Theta$ in~$Z_{G(\R)^+}(H)$ is~$...$. We can summarise with
%\[
%X_H\times X_{Z_{G(\R)^+}(H)}\to X_G.
%\]
%
%\subsubsection{ultrametric case in general}
%In the ultrametric case, the solution is less precise in general. This is answered by~\cite{}.
%
%\subsubsection{ultrametric case with good reduction}The situation of the split case~\S\ref{} with lambda = 1 was the "nicest" (cf. corollaries~\ref{}). When in presence of good reduction (hyperspecial goups in quasisplit reductive groups), we have the analog~\cite[\S7?]{APZ2}.

\subsection{Another case}

\begin{theorem}\label{thm:C}
Let~$L_1,\ldots,L_f$ be finite extensions of~$\Q$ or of~$\Q_\ell$, and~$L:=\bigoplus_{i=1}^f L_i$. Let~$Z\leq O_L:=\bigoplus_{i=1}^f O_{L_i}$ be a subring and~$F\in\Z_{\geq1}$ be such that~$F\cdot O_L\leq Z$. 

Let~$O_V:=\bigoplus_{i=1}^f O_{L_i}^{m_i}\leq V:=\bigoplus_{i=1}^f {L_i}^{m_i}$, and~$R:=\End_{O_L}(O_V)=\bigoplus_{i=1}^f \End_{O_{L_i}}(O_{L_i}^{m_i})$, and~$R^0:=\End_{L}(V)\simeq R\tens_{O_L}L$.

Let~$\Lambda\leq O_V$ be a~$Z$-module and~$n\in\Z_{\geq1}$ be such that
\begin{equation}\label{eq:14b}
\forall N\in \Z, \forall \alpha\in R^0, \alpha(\Lambda)\leq N\cdot  O_V\Rightarrow \alpha\in \frac{1}{n}\cdot N\cdot R. 
\end{equation}

Then
\begin{equation}\label{eq:14a}
\Lambda\geq n\cdot F\cdot O_V.
\end{equation}
\end{theorem}
\begin{proof} Let~$\pi_1=(1,0,\ldots,0),\ldots,\pi_f=(0,\ldots,0,1)$ be the units of the~$L_i\leq L$. Let~$\Lambda_i:=O_{L_i}\cdot \pi_i(\Lambda)$. As~$\Lambda$ is a~$Z$-module and~$F\cdot O_{L_i}\cdot \pi_i\in Z$, we have~$F\cdot \Lambda_i=F\cdot O_{L_i}\cdot \pi_i(\Lambda)\leq \Lambda$. Thus
\[
\Lambda\geq F\cdot \bigoplus_{j=1}^f \Lambda_j.
\]
It is sufficient to prove~$\bigoplus_{j=1}^f \Lambda_j\geq  n\cdot O_V.$
We may thus substitute~$F$ with~$1$ if we replace~$\Lambda$ by~$\bigoplus_{j=1}^f \Lambda_j$.

Let us view~$\Lambda^\vee:=\Hom_{O_L}(\Lambda,O_L)$ as a~$O_L$-submodule of~$V^\vee:=\Hom_{L}(\Lambda,L)$. Then~\eqref{eq:14a} is equivalent to
\begin{equation}\label{eq:16}
\Lambda^{\vee}\leq \frac{1}{n}\cdot O^\vee_V.
\end{equation}
Let~$v:=((1,0,\ldots),\ldots,(1,0,\ldots))\in V$. For~$\phi=(\phi_1,\ldots,\phi_f)\in \Lambda^\vee$, the assumption~\eqref{eq:14b} with~$\alpha:=\phi\tens v\in R_0$ and~$N=1$ gives
\[
(\phi_1\tens(1,0,\ldots),\ldots,\phi_f\tens(1,0,\ldots))
=((\phi_1,0,\ldots),\ldots,(\phi_f,0,\ldots))
  \in \frac{1}{n}\cdot R.
\]
This implies~$\phi\in \frac{1}{n}\cdot O_V^\vee$ and proves~\eqref{eq:16}.
\end{proof}
\subsection{General case}\label{app:D}

Let~$K=\Q$ or~$\Q_\ell$ and~$O_K=\Z$ or~$\Z_\ell$ respectively.
Let~$V$ be a $K$-vector space of finite dimension. 
Let~$R^0,S^0\leq \End_K(V)$ be two semisimple~$K$-algebras such that~$S^0$ is the commutant of~$R^0$. Let~$R\leq R^0$ and~$S\leq S^0$ be $O_K$-orders. 

\begin{theorem}\label{thm:D1}
Let~$W,T\leq V$ be a~$O_K$-lattices.
There exists~$C(S,R,T)$, such that
\begin{itemize}
\item if~$S\cdot W\leq W$
\item and 
\begin{equation}\label{eq:D1hyp}
\exists n\in \Z_{\geq1}, \forall N\in\Z_{\geq1},
\forall \alpha \in R^0, 
\alpha(W)\leq N\cdot T\Rightarrow n\cdot \alpha \in N\cdot R.
\end{equation}
\end{itemize}
then 
\[
W\geq C(S,R,T)\cdot n\cdot T.
\]
\end{theorem}


We know that~$R^0$ is the commutant of~$S^0$, and that, for some finite extension~$L/K$, 
one can decompose~$V_L:=V\tens_K L$ as
\[
V_L\simeq \bigoplus_{i=1}^k L^{n_i}\tens L^{m_i}
\]
in such a way that~$R_L:=R\tens L$ and~$S_L:=S\tens L$ act as
\[
\prod_{i=1}^k \End(L^{n_i})\tens Id_{L^{m_i}}\text{ and }
\prod_{i=1}^k  Id_{L^{n_i}} \tens \End(L^{m_i}).
\]
and that~$R$ and~$S$, acting on~$ \bigoplus_{i=1}^k L^{n_i}$ and on~$\bigoplus_{i=1}^k L^{m_i}$,
preserve the integral structures~$ \bigoplus_{i=1}^k O_L^{n_i}$ and on~$\bigoplus_{i=1}^k O_L^{m_i}$.

Let~$R^\bot\leq R^0$ denote the orthogonal of~$R\leq R^0$ with respect to the~$K$-bilinear trace form~$R_0\times R_0\to K$.
Recall that~$R\leq R^\bot$ and that the discriminant, say~$d_R$, of~$R$ is~$[R^\bot:R]$. In particular,~$R\geq d_R\cdot R^\bot$. In~$R_L$ we have
\[
(R\tens O_L)^\bot=R^\bot\tens O_L,
\]
with respect to the~$L$-bilinear trace form~$R_L\times R_L\to L$. We also have
\begin{multline}
d_R\cdot R^\bot\tens O_L
\leq R\tens O_L\leq \prod^k_{i=1}\End(O_L^{n_i})=\left(\prod^k_{i=1}\End(O_L^{n_i})\right)^\bot
\\
\leq (R\tens O_L)^\bot = R^\bot\tens O_L.
\end{multline}
We deduce
\begin{equation}\label{eq:b}
d_R\cdot \prod^k_{i=1}\End(O_L^{n_i})\leq  R\tens O_L\leq \prod^k_{i=1}\End(O_L^{n_i}).
\end{equation}
Similarly,
\[
d_S\cdot \prod^k_{i=1}\End(O_L^{m_i})\leq  S\tens O_L\leq \prod^k_{i=1}\End(O_L^{m_i}).
\]

Let~$W$ and~$T$ satisfy the assumption in Theorem~\ref{thm:D1}. Let~$W'=W\tens O_L$ and~$S'=S\tens O_L$
and~$E= \prod^k_{i=1}Id_{n_i}\tens \End(O_L^{m_i})$ and~$\lambda=d_S$. We have
\[
\lambda\cdot E\cdot W'\leq S'\cdot W' \leq W'.
\]
We may thus apply Theorem~\ref{thm:split linear}. We deduce
\begin{equation}\label{eq:a}
d_S\cdot\bigoplus_{i=1}^k O_L^{n_i}\tens_{O_L}\Lambda_i\leq
W'\leq\frac{1}{d_S}\cdot  \bigoplus_{i=1}^k O_L^{n_i}\tens_{O_L}\Lambda_i
\end{equation}
There exists~$c(T)\in\Z_{\geq1}$ such that
\begin{equation}\label{eq:c}
c(T)\cdot T\tens O_L\leq \bigoplus_{i=1}^k O_L^{n_i}\tens O_L^{m_i}.
\end{equation}

We observe that, for all~$m\in\Z_{\geq1}$,
\begin{align}
\bigoplus_{i=1}^k \Lambda_i
\geq m\cdot \bigoplus_{i=1}^k O_L^{m_i}
&\Rightarrow
\bigoplus_{i=1}^k O_L^{n_i}\tens_{O_L}\Lambda_i
\geq m\cdot \bigoplus_{i=1}^k O_L^{n_i}\tens O_L^{m_i}
\\
&\Rightarrow
W'\geq m\cdot d_S\cdot  \bigoplus_{i=1}^k O_L^{n_i}\tens O_L^{m_i}
\\
&\Rightarrow
W\tens O_L=W'\geq m\cdot c(T)\cdot d_S\cdot T\tens O_L
\\
&\Leftrightarrow 
W\geq m\cdot c(T)\cdot d_S\cdot  T.
\end{align}

Let~$\phi(W,T,R,n)$, where~$n\in\Z_{\geq1}$, denote the formula~\eqref{eq:D1hyp}.

Let~$\Omega:=\{\beta\in \End(V)|\beta(W)\subseteq T\}\simeq W^\vee\tens T$. This is a~$O_K$-lattice of~$\End(V)\simeq V^\vee\tens V$. Then we have
\[
\phi(W,T,R,n)\Leftrightarrow \Omega\cap R^0\subseteq \frac{1}{n}\cdot R.
\]
Let~$R':=R\tens O_L$, and~$T':=T\tens O_L$ and~$\Omega'=W'\tens_{O_L}T'\simeq \Omega\tens O_L\leq \End(V_L)$.
The following is immediate:
\begin{equation}\label{eq:alpha}
\phi(W,T,R,n)\Rightarrow \phi(W',T',R',n).
\end{equation}
Let~$W'':=
\bigoplus_{i=1}^k O_L^{n_i}\tens_{O_L}\Lambda_i$ and~$T'':=\bigoplus_{i=1}^k O_L^{n_i}\tens O_L^{m_i}$
and~$R'':=\prod^k_{i=1}\End(O_L^{n_i})$.
Using~\eqref{eq:a} and~\eqref{eq:b} and~\eqref{eq:c}, we get
\begin{equation}\label{eq:beta}
\phi(W',T',R',n)\Rightarrow \phi(W'',T'',R'',n\cdot d_S\cdot d_R\cdot c(T)).
\end{equation}
Let~$W''':=\bigoplus_{i=1}^k \Lambda_i$ and~$T''':=\bigoplus_{i=1}^k O_L^{m_i}$. 
We also note that
\begin{equation}\label{eq:gamma}
\phi(W'',T'',R'',n\cdot d_S\cdot d_R\cdot c(T))\Leftrightarrow 
\phi(W''',T''',R'',n\cdot d_S\cdot d_R\cdot c(T)).
\end{equation}
Let~$Z=\bigoplus_{i=1}^k  O_L\leq R''=\bigoplus_{i=1}^k \End(O_L^{m_i})$ be the centre of~$R''$.
We note that the~$\Lambda_i$ are~$O_L$-submodules of~$O_L^{m_i}$. Therefore~$W''':=\Lambda_i$ is a~$Z$-submodule
of~$\bigoplus_{i=1}^k O_L^{m_i}$.

We are ready to finish the proof of Theorem~\ref{thm:D1}. By~\eqref{eq:D1hyp}, the formula~$\phi(W,T,R,n)$ holds true.
By~\eqref{eq:alpha},~\eqref{eq:beta},~\eqref{eq:gamma} and the assumption~\eqref{eq:D1hyp}, we have~$\phi(W''',T''',R'',n\cdot d_S\cdot d_R\cdot c(T))$. This implies~\eqref{eq:14b}. Moreover,~$W'''$ is a~$Z$-module with~$F=1$. We can thus apply Theorem~\ref{thm:C}. We deduce
\[
W'''\geq n\cdot d_S\cdot d_R\cdot c(T)\cdot T'''.
\]
From~\eqref{eq:a} and~\eqref{eq:c}, we conclude
\[
W\geq n\cdot {d_S}^2\cdot d_R\cdot c(T)^2\cdot  T.
\]
This proves the conclusion of Theorem~\ref{thm:D1}, with
\begin{equation}
C(S,R,T):={d_S}^2\cdot d_R\cdot c(T)^2.
\end{equation}


\section{On abelian varieties}

Let~$A_P,B\leq A$ be abelian varieties over~$\C$, and let~$P\in A_P$ be such that~$\Z\cdot P$ is Zariski dense in~$A_P$. Let~$\pi:A\to A':=A/B$ be the quotient map.

We recall that there exists~$B'\leq A$ such that~\(
A=B'+B\) and \(B\cap B'\) is finite.
There exists thus~$n\in\Z_{\geq1}$ be such that
\begin{equation}\label{eq:n}
B\cap B'\leq A[n]. 
\end{equation}

\begin{theorem}\label{thm:E1}
 Consider~$Q\in A$ and~$T\in {A}_{tors}$ and~$\phi\in\Hom(A_P,A)$ such that~$d\cdot Q=\phi(P)+T$ for some~$d\in\Z_{\geq1}$. Let~$d'|d$ be such that
\[
\pi\circ \phi(A_P[d])\leq A'[d'].
\]
Then there exists~$Q'\in Q+B$ and~$T'\in {A}_{tors}$ and~$\phi'\in \Hom(A_P,A)$ such that
\[
n\cdot d'\cdot Q'=\phi'(P)+T'.
\]
\end{theorem}
\subsection{Auxiliary Definitions and Results}
We define
\[
\Theta:=\{\phi(P)|\phi\in \Hom(A_P,A)\}\text{ and }
\Theta':=\{\phi(P)|\phi\in \Hom(A_P,A')\}.
\]
For~$d\in\Z_{\geq1}$, let
\[
\Theta_d:=\{Q\in A|d\cdot Q\in \Theta\}\qquad\text{and}\qquad
\Theta'_d:=\{Q\in A'|d\cdot Q\in \Theta'\}.
\]
We claim that the map
\[
\Hom(A_P,A)\to \Theta\quad\text{ and }
\quad \Hom(A_P,A')\to \Theta'
\]
given by~$\phi\mapsto\phi(P)$ are bijective.
\begin{proof}The map are subjective by definition of~$\Theta$ and~$\Theta''$. Observe that~$\phi\mapsto \phi(P)$ is an additive morphism. The injectivity follows from
\[
\phi(P)=0\Rightarrow P\in \ker(\phi)\Rightarrow \ol{P\cdot \Z}^{Zar}\leq \ker(\phi)\Rightarrow \phi(A_P)=0\Rightarrow \phi=0.
\qedhere\]
\end{proof}
Recall that for an algebraic subgroup~$G$ of an abelian variety~$A$, and a morphism~$\phi:A\to B$ of abelian varieties, we have
\[
G\leq \ker(\phi)\Leftrightarrow \exists \psi\in \Hom(A/G,B), \phi=\psi\circ \pi,
\]
where~$\pi:A\to A/G$ is the quotient map. In particular, for~$d\in\Z_{\geq1}$ and~$G=A[d]$, 
then
\begin{equation}\label{eq:lift kernel isogeny}
\forall \phi\in\Hom(A,A), A[d]\leq \ker(\phi)\Leftrightarrow \phi\in d\cdot \Hom(A,A).
\end{equation}


\begin{lemma}\label{lem:C2}
Let~$n$ be as in~\eqref{eq:n}. Let~$\pi_{\star}$ be either of the maps
\[
\phi\mapsto \pi\circ \phi:\Hom(A_P,A)\to\Hom(A_P,A')\text{ and }\phi(P)\mapsto \pi\circ\phi(P):\Theta\to \Theta'.
\]
Then for any~$x\in \coker(\pi_{\star})$, we have~$n\cdot x=0$.
\end{lemma}
\begin{proof}
Let~$B'\leq A$ be as in~\eqref{eq:n}, and let~$p=\pi|_{B'}: B' \to A'$. We recall that~$\ker(p) = B \cap B'\leq B'[n]$ by~\eqref{eq:n}.  It follows that there exists a homomorphism~$\sigma: A' \to B'$ such that $p \circ \sigma=n Id_{A'}$ is the multiplication-by-$n$ map on~$A'$. 

For~$\phi'\in \Hom(A_P,A')$, we have 
\[
n\cdot \phi'=\pi\circ \sigma\circ \phi',
\]
and thus~$\phi'=\pi_\star(\phi)$ for~$\phi=\sigma\circ \phi'$. For~$\phi'(P)\in \Theta'$, we have~$n\cdot \phi'(P)=\pi( \phi(P))$. The lemma follows.
\end{proof}
%\begin{proof}
%Let~$p := \pi|_{B'}: B' \to A'$. We recall that~$\ker(p) = B \cap B'\leq B'[n]$ by~\eqref{eq:n}. Since~$A' \cong B'/\ker(p)$, there exists by~\eqref{eq:lift kernel isogeny} a homomorphism~$\sigma: A' \to B'$ such that $p \circ \sigma$ is the multiplication-by-$n$ map on~$A'$. Specifically, for any~$\phi'\in \Theta'$, we have~$n\cdot \phi'=p\circ\sigma\circ \phi'$. Thus, for any~$\phi' \in \Hom(A_P, A')$, the map~$n \cdot \phi': A_P \to A'$ factors through~$B'$. That is, there exists~$\psi: A_P \to B'$ such that~$p \circ \psi = n \cdot \phi'$.
%
%Then~$p:=\pi|_{A'}:B'\to A'$ satisfies~$\ker(p)\leq A[n]\cap B'=B'[n]$. There exists thus~$\sigma:A'\to B'$ such that~$\sigma\circ p:\in\End(B')$ and~$p\circ \sigma\in\End(A')$ are the multilication-by-$n$. For~$\phi'\in \Theta'$, we have~$n\cdot \phi'=p\circ\sigma\circ \phi'$. It follows that the map~$\widetilde{\phi'}:A_P\xrightarrow{\phi'} B'\xrightarrow{\sigma} B'\hookleftarrow A$ satisfies
%\[
%\pi\circ \widetilde{\phi'}=n\cdot \phi'
%\]
%and thus~$\pi\circ \widetilde{\phi'}(P)=n\cdot \phi'(P)$. 
%This means that for any~$\phi' \in \Hom(A_P, A')$, its coset in~$\Hom(A_P, A') / \operatorname{Im}(\pi_{\ast})$ is annihilated by~$n$, and similarly the coset of~$\phi'(P)$ in $\Theta'/\pi_{\star}(\Theta)$. This proves the lemma.
%\end{proof}


\begin{lemma}\label{lem:C3}
For every~$d\in\Z_{\geq1}$
\begin{enumerate}
\item the map~$A[d]\to A'[d]$ is surjective \label{enum:quotient et torsion}
\item \label{enum:2}
and
\[
\forall Q'\in\Theta_d', \exists \widetilde{Q'}\in \Theta_{n\cdot d}, \pi(\widetilde{Q'})=Q'. 
\]
\end{enumerate}
\end{lemma}
\begin{proof}It is a fact that the exact sequence~$0\to B\to A\to A'\to 0$ is split in the category 
of real Lie groups, and a fortiori as a sequence of abstract groups. The first assertion follows.

Let us prove the second assertion. Let~$Q'\in \Theta'_{d}$ and choose~$\phi'\in \Hom(A_P,A')$ such that~$d Q'=\phi'(P)$. Let~$\sigma$ be as in the proof of Lemma~\ref{lem:C2} and~$\phi:=\sigma\circ \phi'$, and let~$Q$ be such that~$nQ=\sigma(Q')$. 
We have
\[
\pi(n(Q))=\pi(\sigma(Q'))=nQ'.
\]
Therefore~$\pi(Q)-Q'=T'$ for some~$T'\in A'[n]$. By~\eqref{enum:quotient et torsion}, we have~$T'=\pi(T)$ for some~$T\in A[n]$. We have~$Q+T\in \Theta_{nd}$ because
\[
nd (Q+T)=d\sigma(Q')=\sigma(dQ')= \sigma\circ \phi'(P)=\phi(P).
\]
This proves the lemma with~$\wt{Q'}=Q+T$.
\end{proof}
\subsection{Proof of Theorem~\ref{thm:E1}}

From~$\pi\circ \phi(A_P[d])\leq A'[d']$ follows~$A_P[d/d']\leq \ker(\pi\circ \phi)$. This implies that there exists~$\phi':A_P\to A'$ such that~$\pi \circ \phi=(d/d')\cdot \phi'$. Therefore
\[
(d/d')\cdot \phi'(P)=\pi\circ \phi(P)=d\cdot \pi(Q)-\pi(T)=(d/d')\cdot d'\cdot \pi(Q)-\pi(T).
\]
Thus~$(d/d')\cdot (\phi'(P)-d'\cdot \pi(Q))=\pi(T)\in A'_{tors}$, and~$d'\cdot \pi(Q)\in \phi'(P)+A'_{tors}$. There exists thus~$T'\in A'_{tors}$ such that~$d'\cdot (\pi(Q)+T')=\phi'(P)$. Since~$Q':=\pi(Q)+T'\in \Theta'_{d'}$, Lemma~\ref{lem:C3} implies that there exists~$\wt{Q'}\in \Theta_{nd'}$ such that~$\pi(\wt{Q'})=\pi(Q)+T'$. There exists~$\wt{T'}\in A_{tors}$ such that~$\pi(\wt{T'})=T'$. 
Let~$\wt{Q}:=\wt{Q'}-\wt{T'}$. Then~$\pi(\wt{Q})=\pi(Q)$, that is~$\wt{Q}\in Q+B$. Since~$\wt{Q'}\in \Theta_{nd'}$, there exists~$\wt{\phi}\in \Hom(A_P,A)$ such that~$nd'\wt{Q'}=\wt{\phi}(P)$. Let~$\wt{T}=-nd'\wt{T'}\in A_{tors}$. We conclude~$nd'\wt{Q}=\wt{\phi}(P)+\wt{T}$. This proves Theorem~\ref{thm:E1}.


\section{On the Deligne-Shimura Reciprocity law}\label{sec:appendix reciprocity}
\subsection{Uniformity properties of the reflex norm}

Let~$(G,X)$ be a Shimura datum and let~$(H,X_H)$ be a subdatum. We condider~$T:=Z(H)^0$ the connected centre of~$H$ and~$C:=H^{ab}$ the maximal abelian quotient of~$H$.

Let~$x\in X_H$ and consider the special Shimura datum~$(C,x^{ab})$. Let~$\mu:GL(1)_\C\to C_\C$ be the compositum~$\C^\times \to \mathbf{S}(\C)\xrightarrow{x^{ab}}C(\C)$. Let~$Gal(\ol{\Q}/\Q)$ act on~$Y(C_{\ol{\Q}})$ and let~$E:=E(C,x^{ab})$ be such that~$Gal(\ol{\Q}/E)$ is the stabiliser of~$\mu$.

Then~$T_E:=Res_{E/\Q}GL(1)$ has a natural basis~$e_{\sigma(\mu)}$ indexed by the~$\sigma(\mu)\in Gal(\ol{\Q}/\Q)\cdot \mu$. Recall that~$Y(C)\tens\Q$ is $\Q$-linearly generated by the conjugates~$\sigma(\mu)\in Gal(\ol{\Q}/\Q)\cdot \mu$. The natural~$\Z$-linear map
\[
Y(T_E)\simeq \bigoplus_{\sigma(\mu)\in Gal(\ol{\Q}/\Q)\cdot \mu} \Z \cdot e_{\sigma(\mu)}
\to 
\sum_{\sigma(\mu)\in Gal(\ol{\Q}/\Q)\cdot \mu} \Z\cdot \sigma(\mu)\hookrightarrow Y(C)
\]
corresponds to a morphism of $\Q$-algebraic tori~$T_E\to C$. Let $(e^*_{\sigma(\mu)})_{\sigma(\mu)\in Gal(\ol{\Q}/\Q)\cdot \mu}$ be the basis of~$X(T_E)$ dual to the basis~$(e_{\sigma(\mu)})_{\sigma(\mu)\in Gal(\ol{\Q}/\Q)\cdot \mu}$.
By~\cite[Lem. 2.6 and \S2,1]{UY}, there exists~$N\in\Z_{\geq1}$ depending only on~$(G,X)$ and a subset
\[
F\subseteq \bigoplus_{\sigma(\mu)\in Gal(\ol{\Q}/\Q)\cdot \mu} \{-N;\ldots;N\}\cdot e^*_{\sigma(\mu)}\subseteq X(T_E),
\]
 such that the image of
\[
X(C)\to X(T_E)
\]
is generated by~$F$.
\begin{theorem}\label{thm:F1}
 Let~$(G,X)$ be a Shimura datum. 
There exists~$e(G,X)$ such that for every Shimura subdatum~$(H,X_H)$ of~$(G,X)$ and Hodge generic~$x\in X_H$, the reciprocity norm morphism
\[
N:Res_{E(H^{ab},\{x^{ab}\})/\Q}GL(1)\to H^{ab}
\]
satifies
\[
\forall g\in H^{ab}(\wh{\Z}), g^{e(G,X)}\in  N\left(\wh{O_{E(H^{ab},\{x^{ab}\})}}^\times\right),
\]
where~$H^{ab}(\wh{\Z})\leq H^{ab}(\A_f)$ is the maximal compact subgroup.
\end{theorem}
\begin{proof} This follows from the preceding discussion and Proposition~\ref{prop:F3} below.
\end{proof}
\subsection{}
Let~$\pi:T\to S$ be a surjective morphism of~$\Q$-algebraic groups. Recall that there exists~$S'\leq T$ a subtorus such that~$\alpha:S'\to T\to S$ is an isogeny. There exists thus an~$e\in\Z_{\geq1}$ such that~$\ker(\alpha)=S\cap \ker(\pi)\leq T[e]$.

% Let~$e\in\Z_{\geq1}$ be suc that~$N\cap S'\leq T[e]$.

Let~$T(\wh{\Z})\leq T(\A_f)$ and~$S(\wh{\Z})\leq S(\A_f)$ be the maximal compact subgroups. 
\begin{proposition}\label{prop:F2}
We have
\[
\forall s\in S(\wh{\Z}), s^e\in \pi(T(\wh{\Z})).
\]
\end{proposition}
\begin{proof}Note that~$\alpha:S'\to T\to S$ is an isogeny, and that~$\ker(\alpha)\leq S'[e]$. Recall that there exists an isogeny (the dual isogeny)~$\alpha':S\to S'$ such that~$\alpha\circ \alpha':S\to S$ is the morphism~$s\mapsto s^e$. 

Note that~$\alpha':S(\A_f)\to S'(\A_f)$ is a continuous map of separated topological spaces. Therefore~$\alpha'(S(\wh{\Z}))$ is a compact subgroup of~$S'(\A_f)$. We have thus~$\alpha'(S(\wh{\Z}))\leq S'(\wh{\Z})\leq T(\wh{\Z})$. 

We conclude
\[
\pi(T(\wh{\Z}))\geq \alpha(S'(\wh{\Z}))\geq \alpha(\alpha'(S(\wh{\Z})))=\{s^e|s\in S(\wh{\Z})\}.\qedhere
\]
\end{proof}

\begin{proposition}\label{prop:F3}
Consider~$r,N\in\Z_{\geq1}$.
There exists~$e(r,N)$ such that for every transtive action~$Gal(\ol{\Q}/\Q)$ on a set~$E$ of cardinality~$\#\leq r$ and every~$\Z[Gal(\ol{\Q}/\Q)]$-submodule~$Y\leq \Z^E$ generated by a subset of~$\{-N;\ldots;N\}^E$, the morphism of~$\Q$-algebraic tori
\[
T\to S
\] 
corresponding to~$Y\to \Z^E$ satisfies
\[
\forall t\in S(\wh{\Z}), s^e\in \pi(T(\wh{\Z})).
\]
\end{proposition}
\begin{proof}
Choose a bijection~$E\simeq \{1;\ldots; r'\}$ and let~$G$ be the image of~$Gal(\ol{\Q}/\Q)\to S_E\simeq S_{r'}$. For each~$G$ and each subset~$F\subseteq\{-N;\ldots;N\}^E$, let
\[
\alpha(G,F):GL(1)_{\ol{\Q}}^r\to S(G,F)
\]
be the associated morphism of~$\ol{\Q}$-algebraic tori. We claim that there exists a~$G$-invariant sub-torus~$S'(r,F)\leq GL(1)_{\ol{\Q}}$ such that~$\beta(G,F):S'(G,F)\to GL(1)_{\ol{\Q}}^r\to S(G,F)$ is an isogeny.
\begin{proof}Recall that the action of~$G$ on~$\Q^r$ is semisimple. Thre is exists thus~$V(G,F)$ be a~$G$-invariant $\Q$-linear subspace such that~$\Q^r=V(G,F)+\sum_{f\in F} f\cdot Q$ and~$\{0\}=V(G,F)\cap \left(\sum_{f\in F} f\cdot Q\right)$. Let~$GL(1)_{\ol{\Q}}^r\to T(G,F)$ be the corresponding morphism of tori. Then the torus~$S'(r,F)=\ker(\alpha')^0$ is as claimed.
\end{proof}
 Let~$e(G,F)\in\Z_{\geq1}$ be such that~$\forall g\in \ker(\beta(G,F)), g^{e(G,F)}=0$. Since there are finitely many possibilities for~$G$ and~$F$, 
\begin{multline}
e(r,N):={gcd}
\Bigl(
\Bigl\{e(G,F)|r'\in\{1;\ldots;r\},\\
G\leq S_{r'}\text{ a transitive subgroup},\\
F\leq \{1;\ldots;N\}^{r'}
\Bigr\}
\Bigr)
<+\infty.
\end{multline}
Then Proposition~\ref{prop:F3} follows from Proposition~\ref{prop:F2}.
\end{proof}
\subsection{Application to reciprocity and abelian varieties}
Recall that the morphism
\[
r:Gal(\ol{\Q}/E)\to M^{ab}(\A_f)/M^{ab}(\Q)
\]
is given by the Deligne-Shimura reciprocity law. 
If we denote Artin morphism from class field theory by
\[
\alpha:Gal(\ol{\Q}/E)^{ab}\to T_E(\A_f)/T_E(\Q)
\]
where~$T_E:=Res_{E/\Q}GL(1)$, then~$r$ is given by
\[
Gal(\ol{\Q}/E)\to Gal(\ol{\Q}/E)^{ab}\xrightarrow{\alpha} T_E(\A_f)/T_E(\Q)
\xrightarrow{N}
 M^{ab}(\A_f)/M^{ab}(\Q)
\]
where the last morphism is induced by~$N:T_E\to M^{ab}$.

Let~$A$ be an abelian variety over a field of finite type~$K/\Q$ and let~$C$ be abelian variety with CM. Up to isogeny,~$C$ is a principal CM abelian variety, and has a model over the Hilbert class field extension, say~$E(C)/E_\Phi$ of the reflex field of the CM type~$\Phi$ of~$C$.

Let~$M$ be the Mumford-Tate group of~$A\times C$ and~$x:\mathbb{S}\to M_\R$ be the Hodge cocharacter of~$ A\times C$. The action of~$Gal(\ol{\Q}/K\cdot E(C))$ on~$(A\times C)_{\tors}$ induces a representation
\[
Gal(\ol{\Q}/E(C)\cdot K)\to M(\wh{\Z}).
\]
The map~$Gal(\ol{\Q}/E(C)\cdot K)\to M^{ab}(\wh{\Z})/M^{ab}(\Z)\leq M^{ab}({\A_f})/M(\Q)$ is then given by
\[
Gal(\ol{\Q}/E(C)\cdot K)\to Gal(\ol{\Q}/E(M^{ab},\{x^{ab}\}))\xrightarrow{r} M^{ab}({\A_f})/M^{ab}(\Q).
\]

Recall that Galois group~$Gal(E^{ab}/E(C))$ of the hilbert class field extension~$E(C)/E$ is the image of
\[
\wh{O_{E}}^\times/O_E^{\times}\leq T_E(\A_f)/T_E(\Q)\simeq Gal(E^{ab}/E)^{ab}.
\]
Therefore, the image of~$Gal(\ol{\Q}/E(C))$ in~$M^{ab}({\A_f})/M^{ab}(\Q)$
\begin{itemize}
\item is contained in~$M^{ab}(\wh{\Z})/M^{ab}(\Z)$;
\item and by Theorem~\ref{thm:F1} for~$H=M$, there exists~$e$ such that for
all~$\lambda\in M^{ab}(\wh{\Z})/M^{ab}(\Z)$,
the element~$\lambda^e$ is contained in the image of~$Gal(\ol{\Q}/E(C))$ in $M^{ab}(\wh{\Z})/M^{ab}(\Z)$.
\end{itemize}
Note that~$e=e(GSp(2g),H_g)$ depends only on~$g:=\dim(A\times C)$ is bounded when~$A$ is fixed and the dimension of~$C$ is bounded.
 
Note that there exists~$c:\Z_{\geq1}\to \Z_{\geq1}$ such that the cardinality of~$M^{ab}(\Z)$ is at most~$c(\dim(M^{ab}))$. Moreover~
\[
[K\cdot E(C):E(C)]\text{ divides } 
c_K:=[K\cap \ol{\Q}:\Q]<+\infty.
\]
Therefore the image of
\[
Gal(\ol{\Q}/K\cdot E(C))\to M^{ab}(\wh{\Z})/M^{ab}(\Z)
\] 
contains
\[
\{\lambda^{c_K\cdot e}|\lambda \in M^{ab}(\wh{\Z})/M^{ab}(\Z)\}.
\]

By Lemma~\ref{lem:liftkernel}, the image of
\[
Gal(\ol{\Q}/K\cdot E(C))\to M^{ab}(\wh{\Z})
\]
contains
\[
\{\lambda^{c(\dim(M^{ab})!\cdot c_K\cdot e}|\lambda \in M^{ab}(\wh{\Z})\}.
\]
We have proved the following.
\begin{theorem}\label{thm:Reciprocity principal}
Let~$A$ be an abelian variety over a field of finite type~$K/\Q$ and let~$C$ be principal abelian variety with CM. Let~$\Phi$ be the CM type of~$C$, let~$E_\Phi$ be the reflex field of~$\Phi$ and let~$E(C)/E_\Phi$ be the Hilbert class field extension.
  
There exists~$f:=f(\dim(A\times C),[K\cap \ol{\Q}:\Q])\in\Z_{\geq1}$ such that the image of the morphism
\[
Gal(\ol{K}/K\cdot E(C))\to M(\wh{\Z}) \to M^{ab}(\wh{\Z})
\]
induced by the action on~$(A\times C)_{tors}$ contains
\[
\{\lambda^{f}|\lambda \in M^{ab}(\wh{\Z})\}.
\]
\end{theorem}
%\section{}
%\subsection{OLD}
%
%The following summarises some properties of abelian varieties.
%\begin{theorem}
%Let
%\[
%0\to A_1'\xrightarrow{\iota_1} A_1\to A_1''\xrightarrow{\pi_1} 0\qquad\text{and}\qquad
%0\to A_2'\xrightarrow{\iota_2} A_2\to A_2''\xrightarrow{\pi_2} 0
%\]
%be two exact sequences of abelian varieties over~$\C$.
%\begin{enumerate}
%\item There exist abelian subvarieties~$A_1'''$ of~$A_1$ and~$A_2'''$ of~$A_2$ such that
%\[
%A_1=A_1'+A_1'''\qquad\text{and}\qquad
%A_2=A_2'+A_2'''
%\]
%and, for some~$n_1,n_2\in\Z_{\geq1}$,
%\[
%A_1'\cap A_1'''\leq A_1[n_1]\qquad\text{and}\qquad
%A_2'\cap A_2'''\leq A_2[n_2].
%\]
%\item For each~$i\in\{1;2\}$ and every~$d\in\Z_{\geq1}$, the induced sequence
%\[
%0\to A_i'[d]\xrightarrow{\iota_i} A_i[d]\to A_i''[d]\xrightarrow{\pi_i} 0\quad
%\]
%is exact and split.
%\item The cokernels of
%\[
%\varphi\mapsto \varphi\circ\iota_1:\Hom(A_1,A_2)\to\Hom(A'_1,A_2)
%\]
%and
%\[
%\varphi\mapsto \pi_2\circ\varphi:\Hom(A_1,A_2)\to\Hom(A_1,A''_2)
%\]
%are annulled by~$n_1$ and by~$n_2$ respectively.
%\item Let~$P\in A_1$ be such that~$\ol{P\cdot \Z}^{Zar}=A_1'$. Let
%\[
%\Theta:=\{\phi(P)|\phi\in \Hom(A'_1,A_2)\}\text{ and }
%\Theta'':=\{\phi(P)|\phi\in \Hom(A'_1,A''_2)\}.
%\]
%For~$d\in\Z_{\geq1}$, let
%\[
%\Theta_d:=\{Q\in A_2|d\cdot Q\in \Theta\}\qquad\text{and}\qquad
%\Theta''_d:=\{Q\in A''_2|d\cdot Q\in \Theta''\}.
%\]
%Then the cokernel of the map
%\[
%Q\mapsto \pi_2(Q):\Theta_d\to \Theta''_d
%\]
%is annulled by~$n_2$.
%\end{enumerate}
%\end{theorem}
%\begin{proof}
%For statement (1), the existence of such abelian subvarieties $A_1'''$ and $A_2'''$ and integers $n_1, n_2$ is a direct consequence of \textbf{Poincaré's Complete Reducibility Theorem}. See~\cite[Chapter III, \S 18, Theorem 1]{Mumford_AV}.
%
%
%\textit{Proof of (2):}
%Let the given exact sequence of abelian varieties over $\C$ be $0\to A_i'\xrightarrow{\iota_i} A_i\xrightarrow{\pi_i} A_i''\to 0$.
%Every complex abelian variety is, by definition, a compact connected complex Lie group. Forgetting the complex structure, it is also a compact connected real Lie group. Furthermore, any compact connected abelian Lie group (real or complex) is topologically isomorphic to a torus, i.e., a product of circles $T^n = (S^1)^n$.
%
%A crucial result in the theory of Lie groups states that any short exact sequence of compact, connected, commutative Lie groups always splits.
%More formally, let $0 \to K \xrightarrow{\alpha} G \xrightarrow{\beta} H \to 0$ be a short exact sequence of such Lie groups. Then there exists a continuous homomorphism (a Lie group homomorphism) $s: H \to G$ such that $\beta \circ s = \mathrm{id}_H$. This implies that $G$ is isomorphic to the direct product $K \times H$ as Lie groups.
%
%Applying this result to our sequences of abelian varieties over $\C$: since $A_i'$, $A_i$, and $A_i''$ are compact connected complex Lie groups (and thus also compact connected real Lie groups), the sequences
%\[
%0\to A_i'\xrightarrow{\iota_i} A_i\xrightarrow{\pi_i} A_i''\to 0
%\]
%are split in the category of Lie groups. This means there exist Lie group homomorphisms $s_i: A_i'' \to A_i$ such that $\pi_i \circ s_i = \mathrm{id}_{A_i''}$.
%This splitting implies that $A_i$ is isomorphic to the direct product $A_i' \times A_i''$ as Lie groups. Since they are isomorphic as Lie groups, they are also isomorphic as abstract abelian groups.
%
%Finally, we consider the induced sequence on $d$-torsion subgroups. For any abelian group $G$ that is a direct sum of subgroups $G_1$ and $G_2$, i.e., $G \cong G_1 \oplus G_2$, its $d$-torsion subgroup $G[d] = \{g \in G \mid d \cdot g = 0\}$ is naturally isomorphic to the direct sum of the $d$-torsion subgroups of its factors: $G[d] \cong G_1[d] \oplus G_2[d]$.
%Given $A_i \cong A_i' \oplus A_i''$, it follows that $A_i[d] \cong A_i'[d] \oplus A_i''[d]$.
%The map $\iota_i: A_i' \to A_i$ restricted to $d$-torsion elements is precisely the inclusion $A_i'[d] \to A_i'[d] \oplus A_i''[d]$.
%The map $\pi_i: A_i \to A_i''$ restricted to $d$-torsion elements is precisely the projection $A_i'[d] \oplus A_i''[d] \to A_i''[d]$.
%The section $s_i: A_i'' \to A_i$ similarly restricts to a section $s_i: A_i''[d] \to A_i[d]$ such that $\pi_i \circ s_i = \mathrm{id}_{A_i''[d]}$.
%Therefore, the induced sequence
%\[
%0\to A_i'[d]\xrightarrow{\iota_i} A_i[d]\to A_i''[d]\xrightarrow{\pi_i} 0
%\]
%is a split exact sequence of finite abelian groups.
%\end{proof}
%\begin{lemma}
%Let
%\[
%0\to B'\to B\to B''\to 0
%\]
%be a short exact sequence of abelian varieties over~$\C$.
%
%Then for every~$d\in\Z_{\geq1}$ we have an induced short exact sequence of torsion subgroups
%\[
%0\to B'[d]\to B[d]\to B''[d]\to 0.
%\]
%
%Let~$\Theta\leq B$ be the image of an injective morphism~$\Z^k\to B$. Let~$\Theta''$ be the image of~$\Theta$ in~$B''$. For each~$d\in\Z_{\geq1}$, let~$\Theta_d$ and~$\Theta''_d$ be the inverse image of~$\Theta$ and~$\Theta''$ by~$P\mapsto d\cdot P:B\to B$ and~$P\mapsto d\cdot P:B''\to B''$. AssumeThen that~$\Theta''$ is torsion free. Then? we have an induced surjection
%\[
%\Theta_d\to \Theta''_d.
%\]
%\end{lemma}
%\begin{proof} At the level of real Lie groups, the sequence~\[
%0\to B'\to B\to B''\to 0
%\]
%is split, and isomorphic to
%\[0\to (\R/\Z)^{\dim(B')}\to 
%\to (\R/\Z)^{\dim(B)}\to 0\mapsto (\R/\Z)^{\dim(B'')}.\]
%The sequence~$0\to B'[d]\to B[d]\to B''[d]\to 0$ becomes
%\[0\to (\Z/d\Z)^{\dim(B')}\to 
%\to (\Z/d\Z)^{\dim(B')}\to 0\mapsto (\Z/d\Z)^{\dim(B'')}.\]
% 
%On torsion subgroups, the map~$\Theta_d\to \Theta''_d$, induces a map
%\[
%(\Theta_d)_{tors}\to (\Theta''_d)_{tors}
%\]
%which is also the surjection~$B[d]\to B''[d]$. 
%It is thus enough to prove that the induced map
%\[
%\Theta/(\Theta_d)_{tors}\to \Theta''/(\Theta''_d)_{tors}
%\]
%By construction we have a surjection~$\Theta\to \Theta''$ and therefore
%\[
%\Theta\tens\Q ...\to \Theta''\tens\Q
%\]
%PB Theta free mod B?
%\end{proof}
%



\begin{thebibliography}{99}

\bibitem{AGHM} F. Andreatta, E. Goren, B. Howard, K. Madapusi Pera, {\it Faltings heights of abelian varieties with complex multiplication}, Annals of Mathematics, 187 (2018), 391–531. \url{https://doi.org/10.4007/annals.2018.187.2.3}


\bibitem{BRU} G. Baldi, R. Richard, E. Ullmo, {\it Manin-Mumford in arithmetic pencils.}  \url{https://arxiv.org/2105.12027v2}

\bibitem{Bertrand} D. Bertrand, {\it Galois Representations and Transcendental Numbers
},
New advances in transcendence theory (Durham, 1986), 37–55.
Cambridge University Press, Cambridge, 1988

\bibitem{CD} S. Checcoli and G. A. Dill,
{\it New evidence for Rémond’s generalisation of Lehmer’s
conjecture}, June 24, 2025 \url{https://doi.org/10.48550/arXiv.2506.18776}.


\bibitem{CU} L. Clozel, E. Ullmo {\it Equidistribution de sous-variétés spéciales}, Annals of Mathematics, 161 (2005), 1571–1588. 

\bibitem{Deligne} P. Deligne, {\it Preuve des conjectures de Tate et de Shafarevitch.} Sém. Bourbaki 1983/84, Exp. 616, Astérisque 121-122, 25–41 (1985). \url{https://www.numdam.org/item/SB_1983-1984__26__25_0/}

\bibitem{Duke} W. Duke, {\it Hyperbolic distribution problems and half-integral weight Maass forms}, Invent. math. 92 (1988), 73–90. \url{https://doi.org/10.1007/BF01393215}



\bibitem{Fal91} G. Faltings, {\it Diophantine Approximation on Abelian Varieties}, Annals of Mathematics, 133 (1991), 549–576. \url{https://doi.org/10.2307/2944319}.

\bibitem{Gao} Z. Gao, {\it Towards the André–Oort conjecture for mixed Shimura varieties: The Ax–Lindemann theorem and lower bounds for Galois orbits of special points}, J. reine angew. Math. (Crelles Journal), 732 (2017), 85–146. \url{https://doi.org/10.1515/crelle-2014-0127}

\bibitem{Hindry} M. Hindry, {\it Autour d'une conjecture de Serge Lang},
Invent. math. 94, 575-603 (1988)

\bibitem{Ilya} I. Khayutin, {\it Joint equidistribution of CM points}, Annals of Mathematics, 189 (2019), 145–276. \url{https://doi.org/10.4007/annals.2019.189.1.4}

\bibitem{LN} S. Lang, A. Néron, {\it Rational Points of Abelian Varieties Over Function Fields.} Amer. J. Math. 81, 95–118 (1959). \url{https://doi.org/10.2307/2372851}

\bibitem{McQ} M. McQuillan, {\it Division points on semi-abelian varieties.} Invent. math. 120, 143–159 (1995). \url{https://doi.org/10.1007/BF01241125}

\bibitem{PS} Y. Peterzil, S. Starchenko {\it Definability of restricted theta functions and families of abelian varieties}, Duke Math. J., Tome 162, no 4 (2013), p. 731-765.
\url{https://doi.org/10.1215/00127094-2080018}

\bibitem{PT} J. Pila, J. Tsimerman, {\it Ax-Lindemann for $\mathcal{A}_g$}, Annals of Mathematics, 179 (2014), 659–681. \url{https://doi.org/10.4007/annals.2014.179.2.4}



\bibitem{Pink} R. Pink, {\it A combination of the conjectures of Mordell-Lang and André-Oort.} Progr. Math. 235, 251–282 (2005).

\bibitem{Ribet} K. A. Ribet, {\it Kummer Theory on Extensions of Abelian Varieties by Tori},
Duke Mathematical Journal Vol. 46, No. 4 December 1979.

\bibitem{RWeyl} R. Richard,	
{\it Manin-Mumford par le critère de Weyl}
Journal of Number Theory
Volume 239, October 2022, Pages 137-150
\url{https://doi.org/10.1016/j.jnt.2021.11.007}

\bibitem{RU} R. Richard, E. Ullmo {\it Équidistribution de sous-variétés faiblement spéciales et
o-minimalité: André-Oort géométrique}, Ann. Inst. Fourier, Grenoble, Tome 74, no 6 (2024), p. 2667-2721.
\url{https://doi.org/10.5802/aif.3644}

\bibitem{RY:LMS1} R. Richard, A. Yafaev, {\it A Common Generalisation of the André-Oort and André-Pink-Zannier Conjectures.} Preprint, arXiv:2401.03528 (2024). \url{https://doi.org/10.48550/arXiv.2401.03528}


\bibitem{RY:IHES} R. Richard, A. Yafaev, {\it Generalised André-Pink-Zannier conjecture for Shimura varieties of Abelian type.} Publ. math. IHES 141, 249–331 (2025). \url{https://doi.org/10.1007/s10240-025-00154-4}

\bibitem{RY:LMS2} R. Richard, A. Yafaev, {\it Heights on 'Hybrid orbits' in Shimura varieties.} Preprint, arXiv:2406.20013 (2024). \url{https://doi.org/10.48550/arXiv.2406.20013}

\bibitem{RY:SHecke} R. Richard, A. Yafaev, {\it Inner Galois Equidistribution in S-Hecke orbits}, preprint (2017). \url{https://doi.org/10.48550/arXiv.1711.03009}

\bibitem{LNM5}   J.-P. Serre, {\it  Cohomologie galoisienne.} Fifth edition
Lecture Notes in Math., 5. Springer-Verlag, Berlin, 1994.

\bibitem{Serre:MWT} J.-P. Serre, {\it Lectures on the Mordell-Weil Theorem.} 3rd edn. Aspects of Mathematics, Vieweg+Teubner Verlag Wiesbaden (1997). \url{https://doi.org/10.1007/978-3-663-10632-6}

\bibitem{Serre:O4} J.-P. Serre, {\it Oeuvres -- Collected Papers IV: 1985 -- 1998.} Springer Collected Works in Mathematics, Springer Berlin, Heidelberg (2000).


\bibitem{Tsimerman} J. Tsimerman, {\it The André-Oort conjecture for $\mathcal{A}_g$}, Annals of Mathematics, 187 (2018), 379–390. \url{https://doi.org/10.4007/annals.2018.187.2.2}

\bibitem{U} E. Ullmo, {\it Manin-Mumford, Andr´e-Oort, the
equidistribution point of view.} In: Granville, A., Rudnick, Z. (eds) Equidistribution in Number Theory, An Introduction. NATO Science Series, vol 237. Springer, Dordrecht. \url{https://doi.org/10.1007/978-1-4020-5404-4_7} (2007)


\bibitem{UY} E. Ullmo, A. Yafaev, {\it Galois orbits and equidistribution
of special subvarieties: towards
the André-Oort conjecture},
Annals of Mathematics 180 (2014), 823–865
\url{http://dx.doi.org/10.4007/annals.2014.180.3.1}


\bibitem{VdD} L. Van den Dries,  {\it Tame Topology and O-Minimal Structures.} Cambridge University Press; 1998.

\bibitem{YZ} X. Yuan, S.-W. Zhang, {\it On the averaged Colmez conjecture}, Annals of Mathematics, 187 (2018), 533–638. \url{https://doi.org/10.4007/annals.2018.187.2.4}


%\bibitem{OLD} OLD
%
%
%\bibitem{AD} V. Aslanyan, C. Daw {\it A note on unlikely intersections in Shimura varieties}, Journal of Number Theory, Volume 260, July 2024, Pages 212-222.\url{https://doi.org/10.1016/j.jnt.2024.02.015}.
%
%\bibitem{BD} F. Barroero, G. Dill, {\it Hecke orbits and the Mordell-Lang conjecture in distinguished categories} \url{https://arxiv.org/pdf/2401.11193}
%
%%\bibitem{A} Y. Andr\'e, {$G$-functions and geometry.} Publication of the Max-Planck-Institut für Mathematik, Bonn, 1989.
%
%%\bibitem{AM} A. Asan, H. Menken, {On some properties of Liouville Numbers
%%in the non-Archimedean case.} European Journal of Pure and Applied Mathematics,
%%Vol. 6, n. 2, 2013, 239-246.
%%
%%\bibitem{AD} V. Aslanyan, C. Daw {\it A note on unlikely intersections in Shimura varieties}
%%\url{https://arxiv.org/abs/2209.07967}
%%
%%\bibitem{BailyBorel}  W. L. Baily Jr., A.~Borel, {\it Compactification of arithmetic quotients of bounded symmetric domains}, Ann. Math. 84 (1966), pp.~442--528.
%%
%%\bibitem{BKU} G. Baldi, B. Klingler  and R. Ullmo, {\it On the distribution of the Hodge locus.} Invent. math. (2023). \url{https://doi.org/10.1007/s00222-023-01226-0}
%
%\bibitem{BiSYa} G. Binyamini, H. Schmidt, A. Yafaev, A. {\it Lower bounds for Galois orbits of special points on Shimura varieties: a point-counting approach}. Math. Ann. 385, 961–973 (2023). \url{https://doi.org/10.1007/s00208-021-02309-0}
%
%%\bibitem{Bogo} F.  Bogomolov, {\it Sur l'alg\'ebricit\'e des repr\'esentations $l$-adiques}, C.R.A.S. 290, %1980, 701-703.
%
%\bibitem{BorelLAG} A.~Borel, {\it Linear algebraic Groups}. Second Edition, GTM 126, Springer-Verlag, 1991.
%
%%\bibitem{BorelExtension} A.~Borel, {\it Some metric properties of arithmetic quotients of symmetric spaces and an extension theorem}, J. Differential Geometry 6 (1972), pp.~543--560.
%
%\bibitem{BorelJi} A.~Borel, L.~Ji {\it Compactifications of Symmetric and Locally Symmetric Spaces}, Mathematics: Theory and Applications, Birkhauser (2006).
%
%%\bibitem{BorelTits} A.~Borel, J.~Tits, {\it \'{E}l\'{e}ments unipotents et sous-groupes paraboliques de groupes r\'{e}ductifs. I},
%%Invent. Math. 12 (1971), pp. 95--104.
%
%
%%\bibitem{Bost} J.-B. Bost, {\it Algebraic leaves of algebraic foliations over number fields.},
%%Publications Mathématiques de l'IHÉS, Volume 93 (2001) , pp.~161-221.
%%
%%\bibitem{BBK-I-VII} N.~Bourbaki {\it Int\'{e}gration, Chapitre VI: Haar measure.}
%%
%%\bibitem{BBKINT} N.~Bourbaki {\it Int\'{e}gration, Chapitre IX: Mesures sur les espaces topologiques s\'{e}par\'{e}s.}
%%
%%\bibitem{Lie1} N.~Bourbaki {\it Groupes et Alg\`{e}bres de Lie, Chapitre 1: Alg\`{e}bres de Lie.}
%
%\bibitem{CO} C.-L. Chai, F. Oort, {\it
%Abelian varieties isogenous to a Jacobian.}  Ann. of Math. (2)176(2012), no.1, 589–635.
%
%%\bibitem{Ullmospeciale1} L. Clozel and E. Ullmo {\it Équidistribution de sous-variétés spéciales}, Ann. of
%%Math. (2) 161 (2005), no. 3, p. 1571–1588.
%
%\bibitem{CU} L. Clozel and E. Ullmo {\it  Equidistribution adélique des tores et équidistribution des points CM}.
%Doc. Math (2006), volume en l'honneur de J. Coates, p. 233-260.
%
%
%
%%\bibitem{0-Co} C. Cornut, {\it Mazur's conjecture on higher Heegner points.} Invent. Math.148(2002), no.3, 495–523.
%%
%%\bibitem{Daw} C. Daw {\it  Around the André-Oort conjecture.} PhD thesis. University College London, 2014.
%%
%%\bibitem{DO2023}
%%C. Daw and M. Orr, {\it
%%Lattices with skew-Hermitian forms over division algebras and unlikely intersections.} J. Éc. polytech. Math.10(2023), 1097–1156.
%%
%%\bibitem{DO2022}
%%C. Daw and M. Orr, {\it
%%Quantitative reduction theory and unlikely intersections.} Int. Math. Res. Not. IMRN(2022), no.20, 16138–16195.
%%
%%\bibitem{DO2021}
%%C. Daw and M. Orr, {\it Unlikely intersections with E×CM curves in A2.} Ann. Sc. Norm. Super. Pisa Cl. Sci. (5)22(2021), no.4, 1705–1745.
%
%
%
%\bibitem{DR}
%C. Daw, Christopher and J. Ren, {\it
%Applications of the hyperbolic Ax-Schanuel conjecture}. Compos. Math.154(2018), no.9, 1843–1888.
%
%\bibitem{Deligne} P.~Deligne, {\it Preuve des conjectures de Tate et de Shafarevitch}, in Séminaire Bourbaki : volume 1983/84, exposés 615-632, Astérisque, no. 121-122 (1985), Talk no. 616, 17 p. \url{https://www.numdam.org/item/SB_1983-1984__26__25_0/}
%
%%\bibitem{Deligne} P.~Deligne, {\it Vari\'{e}t\'{e}s de Shimura: interpr\'{e}tation modulaire, et techniques de construction de mod\`{e}les canoniques.} in {\it Automorphic forms, representations and L-functions,} Proc. Sympos. Pure Math., XXXIII, Part 2, pp. 247--289. 
%%
%%\bibitem{EMS} A.~Eskin, S.~Mozes, N.~Shah, {\it Non Divergence of Translates of Certain Algebraic Measures }, GAFA, Geom. func., anal. 7 (1997),  48--80.
%
%\bibitem{EdYa} B. Edixhoven and A. Yafaev, {\it
%Subvarieties of Shimura varieties.} Ann. of Math. (2)157(2003), no.2, 621–645.
%
%\bibitem{FaFa2} G. Faltings, {\it Diophantine approximation on abelian varieties.} Ann. of Math. (2)133(1991), no.3, 549–576.
%
%\bibitem{FaFa} G. Faltings, {\it Finiteness theorems for Abelian Varieties}, (translated from the original German paper),
%in Arithmetic Geometry, edited by G. Cornell and J.H. Silverman, Springer-Verlag, 1983. 
%
%%\bibitem{FW} G. Faltings, G. Wustholz (Editors), {\it Rational points.} Seminar Bonn/Wuppertal 1983/83, Publications of the MPIM, Bonn.
%%
%%\bibitem{GW} R.~Godement, {\it Domaines fondamentaux des groupes arithm\'etique},
%%S\'eminaire {B}ourbaki, 1962/63.
%%
%%\bibitem{HRS17} P. Habegger, G. Rémond, T. Scanlon, E. Ullmo and A. Yafaev {\it Around the
%%Zilber-Pink conjecture/Autour de la conjecture de Zilber-Pink}, Panoramas et Synthèses, vol. 52, Société Mathématique de France, Paris,
%%2017, Papers corresponding to courses from the conference "États de la Recherche"
%%held at CIRM, Marseille, May 2011.
%%
%%\bibitem{Helgason} S.~Helgason,  {\it Differential geometry, Lie groups, and symmetric spaces}, Pure and Applied Mathematics, 80. 
%%
%%\bibitem{Klingler} B.~Klingler, {\it Hodge loci and atypical intersections: conjectures.} In Motives and Complex Multiplication (2017, to appear) \url{https://arxiv.org/abs/1711.09387}
%%
%%\bibitem{KUY} B. Klingler, E. Ullmo, A. Yafaev, {\it The hyperbolic Ax-Lindemann-Weierstrass conjecture}
%%Pub. Maths. IHES 123 (2016), 333--360.
%%
%%\bibitem{KY} B. Klingler, A. Yafaev, {\it The Andr\'e-Oort conjecture.}, Annals of Mathematics, 180 (2014), 1-59 
%%
%%\bibitem{GIT} F. Kirwan, J. Fogarty, D. Mumford, {\it Geometric invariant theory.} Ergebnisse der Mathematik und ihrer Grenzgebiete, 1994,
%%
%%\bibitem{MarOpp} G. A. Margulis, {\it Oppenheim conjecture.} Fields Medallists' lectures, 272--327.
%%World Sci. Ser. 20th Century Math., 5.
%%
%%\bibitem{Matsushima} Yoz\^{o} Matsushima, {\it Espaces homogènes de Stein des groupes de Lie complexes}, Nagoya Math. J., 16 (1960) pp. 205–218 
%
%\bibitem{McQuillan} M. McQuillan {\it Diophantine approximations and foliations.} Inst. Hautes Études Sci. Publ. Math.(1998), no.87, 121–174.
%
%%\bibitem{Milne} J. Milne {\it Lectures on Shimura varieties.} Lecture notes are available on the author's web page.
%%
%%\bibitem{Moonen} B.~Moonen, {\it Linearity properties of Shimura varieties. I.},
%%J. Algebraic Geom.~7 (1998), no. 3, pp.~539--567.
%%
%%\bibitem{Orr} M. Orr, {\it Families of abelian varieties with many isogenous fibres.}
%%Journal für die reine und angewandte Mathematik 705, 2015, pp. 211-231
%%
%%\bibitem{OrrThesis} M. Orr, {\it La conjecture de Andr\'e-Pink: orbites de Hecke et sous-vari\'et\'es faiblement sp\'eciales.} 
%%PhD thesis, University Paris-Sud (Orsay), 2013.
%
%\bibitem{AO-PST}
%J. Pila, A. N. Shankar, J. Tsimerman 
%with an appendix by H. Esnault and M. Groechenig,
%{\it Canonical Heights on Shimura Varieties and the André-Oort conjecture}. \url{https://arxiv.org/abs/2109.08788v3}
%
%\bibitem{P} J. Pila, {\it Point-counting and the Zilber-Pink conjecture.}
%Cambridge Tracts in Math., 228, Cambridge University Press, Cambridge, 2022.
%%
%%\bibitem{Pinky} R. Pink, {\it A Combination of the Conjectures of Mordell-Lang and Andr\'e-Oort}
%%In: Geometric Methods in Algebra and Number Theory, (Bogomolov, F., Tschinkel, Y., Eds.), Progress in Mathematics 253, Birkhäuser (2005), 251-282.
%%
%%\bibitem{PR} V.~Platonov, A.~Rapinchuk, {\it Algebraic groups and number theory.} (Translated from the 1991 Russian original by Rachel Rowen) Pure and Applied Mathematics, 139. Academic Press, Inc., Boston, MA, (1994)
%%
%%\bibitem{0-Ra} M. Ratner, {\it Interactions between ergodic theory, Lie groups, and number theory.} Proceedings of the International Congress of Mathematicians, Vol. 1, 2 (Zürich, 1994), 157–182. Birkhäuser Verlag, Basel, 1995
%
%\bibitem{Poonen} B. Poonen, {\it Mordell-Lang plus Bogomolov} Invent. math. 137, 413–425 (1999). \url{https://doi.org/10.1007/s002220050331}
%
%
%\bibitem{Raynaud} M. Raynaud, {\it Sous-variétés d'une variété abélienne et points de torsion.} Arithmetic and geometry, Vol. I, 327–352.
%Progr. Math., 35 Birkhäuser Boston, Inc., Boston, MA, 1983
%
%
%%\bibitem{RU} R. Richard {\it R\'esultat g\'eom\'etrique sur les rep\'esentations de groupes r\'eductifs sur un corps ultram\'etrque.} Preprint 2017, available on \url{https://arxiv.org/abs/1706.00301}.
%%
%%\bibitem{R} R. Richard {\it Sur quelques questions d'\'equidistribution en g\'eom\'etrie arithm\'etique.}
%%Ph.D. thesis, 2009. Available on the Internet at https://tel.archives-ouvertes.fr/tel-00438515/.
%%
%%
%%\bibitem{RS} R. Richard, N. Shah {\it Geometric results on linear representations of reductive groups for applications
%%to homogeneous dynamics}, Erg. Th. Dyn. Sys. (2016)
%%
%%\bibitem{RUll} R. Richard, E. Ullmo {\it Équidistribution de sous-variétés spéciales et o-minimalité: André-Oort géométrique}
%%(with an appendix with Jiaming Chen), to appear in 
%%Annales de l’Insitut Fourier \url{https: // arxiv. org/ abs/ 2104. 04439}
%
%\bibitem{APZ2} R. Richard, A. Yafaev,  
%{\it Generalised André-Pink-Zannier conjecture for Shimura varieties of Abelian Type.}  \url{https://arxiv.org/abs/2111.11216}
%
%\bibitem{RYLMS}R. Richard, A. Yafaev,  
%{\it A Common Generalisation of the André-Oort and André-Pink-Zannier Conjectures.} \url{https://arxiv.org/abs/2401.03528}
%
%
%%\bibitem{APZ2-CRAS} R. Richard, A. Yafaev,  
%%{\it Generalised André-Pink-Zannier conjecture for Shimura varieties of Abelian Type.}  Submitted to Comptes Rendus Math. Acad. Sci. Paris.
%
%\bibitem{APZ1} R. Richard, A. Yafaev,  
%{\it Height functions on Hecke orbits and the generalised André-Pink-Zannier conjecture.}  \url{https://arxiv.org/abs/2109.13718}
%
%\bibitem{RY:Hybrid} R. Richard, A. Yafaev,  
%{\it Heights on `hybrid orbits' in Shimura varieties}  \url{https://arxiv.org/abs/2406.20013}
%
%%\bibitem{APZ1-CRAS} R. Richard, A. Yafaev,  
%%{\it On the Generalised André-Pink-Zannier conjecture.} To appear in Comptes Rendus Math. Acad. Sci. Paris.
%%
%%
%%\bibitem{RY-CRAS} R. Richard, A. Yafaev,  
%%{\it Topological and equidistributional refinement of the André-Pink-Zannier conjecture at finitely many places.} C. R. Math. Acad. Sci. Paris357(2019), no.3, 231–235. \url{https://doi.org/10.1016/j.crma.2019.01.013}.
%%
%%
%%\bibitem{RZ} R. Richard, T. Zamojski,  
%%{\it Limit distributions of translated pieces of possibly irrational leaves 
%%in $S$-arithmetic homogeneous spaces.}  Preprint 2016, \url{https://arxiv.org/abs/1604.08494} .
%%
%%
%%\bibitem{RZ-CRAS} R. Richard, T. Zamojski,  
%%{\it Stabilité analytique et convergence locale de translatées en dynamique homogène S-arithmétique.} C. R. Math. Acad. Sci. Paris357(2019), no.3, 241–246. \url{https://doi.org/10.1016/j.crma.2019.02.005}.
%
%\bibitem{Richardson} Richardson, {\it Conjugacy classes of~$n$-tuples in Lie algebras and Lie groups}, Duke Math. Journal, 57 (1988).
%
%%\bibitem{RichardsonMatsushima} Richardson, {\it Affine coset spaces of reductive algebraic groups} Bull. London Math. Soc. , 9 (1977) pp. 38–41.
%%
%%\bibitem{Satake} I. Satake, {\it Algebraic structures of symmetric domains.}
%% Kano Memorial Lectures, 4. Iwanami Shoten, Tokyo; Princeton University Press, Princeton, N.J., 1980. xvi+321 pp.
%

%
%\bibitem{SerreCR} J.-P. Serre, {\it Compl\`{e}te r\'{e}ductibilit\'{e}}, S\'{e}m. Bourb. \no 934 (2004).
%
%\bibitem{SerreLMW} J.-P. Serre, {\it 
%Lectures on the Mordell-Weil theorem. }
%Translated from the French and edited by Martin Brown from notes by Michel Waldschmidt. With a foreword by Brown and Serre. Third edition
%Aspects Math.
%
%%\bibitem{Serre} J.-P. Serre, {\it Lie groups and Lie algebras.}
%%LNM, Vol 1500, 1992. 
%
%\bibitem{SerreOEuvres4} J.-P. Serre, {\it {\OE}uvres, Collected Papers, Volume IV, 1985-1998.}
%
%%\bibitem{Springer} T. A. Springer {\it Linear Algebraic Groups.} Modern Birkhauser Classics, 1998.
%
%\bibitem{TsiAG} J. Tsimerman, {\it
%The André-Oort conjecture for $\mathcal{A}_g$.}
%Ann. of Math. (2)187(2018), no.2, 379–390.
%
%
%\bibitem{Tsi} J. Tsimerman, {\it
%Brauer-Siegel for arithmetic tori and lower bounds for Galois orbits of special points.} J. Amer. Math. Soc.25(2012), no.4, 1091–1117.
%
%%\bibitem{UApp} E. Ullmo, {\it 
%%Applications du théorème d'Ax-Lindemann hyperbolique.} Compos. Math.150(2014), no.2, 175–190.
%
%\bibitem{U} E. Ullmo, {\it Equidistribution des sous-vari\'et\'es sp\'eciales II. }, 
%J. Reine Angew. Math. 606 (2007), p. 193-216.
%
%%\bibitem{U1} E. Ullmo, {\it Quelques applications du th\'eor\`eme d'Ax-Lindemann hyp\'erbolique.}, Compositio Mathematica. 150 (2014), no. 2, 175–190.
%%
%%\bibitem{UY} E. Ullmo, A. Yafaev, {\it Hyperbolic Ax-Lindemann theorem in the co-compact case.}
%%Duke Mathematical Journal, 163 (2014), no. 2, 433-463.
%%
%%\bibitem{UY1} E. Ullmo, A. Yafaev {\it Generalised Tate, Mumford-Tate and Shafarevich conjecture.}
%%Annales scientifiques du Qu\'ebec. 37 (2013) p. 255-284
%%
%%
%%\bibitem{UY2} E. Ullmo, A. Yafaev {\it Algebraic flows on Shimura varieties.} 
%%Preprint, submitted, arXiv:1610.01487
%%
%%\bibitem{UY3} E. Ullmo, A. Yafaev, {\it A characterisation of weakly special subvarieties.}
%%Mathematika, 57, (2011), no. 2, 263 --273
%
%\bibitem{UY4} E. Ullmo, A. Yafaev
%{\it Nombre de classes des tores de multiplication complexe et bornes inférieures pour les orbites galoisiennes de points spéciaux.}
%Bull. Soc. Math. France 143 (2015), no. 1, 197–228.
%
%%\bibitem{UY5} E. Ullmo, A. Yafaev
%%{\it Galois orbits and equidistribution of special subvarieties: towards the André-Oort conjecture.} Ann. of Math. (2)180(2014), no.3, 823–865.
%%
%%\bibitem{0-Va}
%%V. Vatsal {\it Uniform distribution of Heegner points}, Invent. Math. 148 (2002),
%%no. 1, p. 1--46.


\end{thebibliography}



\end{document}



\end{document}
